\providecommand{\zerJ}{\mathcal{J}}
\providecommand{\zerK}{\mathcal{K}}
\providecommand{\zersinc}{\operatorname{sinc}}
\providecommand{\zergE}{\gamma_{\mathrm{E}}}
\providecommand{\zerLL}{\mathcal{L}}
\providecommand{\zerconv}{\operatorname{conv}}
\providecommand{\zerdiag}{\operatorname{diag}}
\providecommand{\zerSym}{\mathfrak{D}}
\providecommand{\zerLog}{\operatorname{Log}}
\providecommand{\anapos}[1]{\bigl(#1\bigr)_{+}}
\providecommand{\ananeg}[1]{\bigl(#1\bigr)_{-}}
\providecommand{\anaone}{\mathbf{1}}
\providecommand{\anaE}{\mathcal{E}}
\providecommand{\anaEx}{\mathbb{E}}
\providecommand{\anaPr}{\mathbb{P}}
\providecommand{\anaM}{\mathcal{M}}
\providecommand{\anaW}{\mathcal{W}}
\providecommand{\anaZ}{\mathcal{Z}}
\providecommand{\anaC}{\mathcal{C}}
\section{Positive filters that cancel zeros}\label{sec:fil}\label{rec:sec:dilation}
\providecommand{\filStir}[2]{\genfrac\{\}{0pt}{}{#1}{#2}}

A probability measure $\mu$ on the logarithmic scale acts on a power $x^s$ by the multiplier
$F_\mu(s)=\int e^{su}\dd\mu(u)$, and a positive average preserves inequalities between the averaged
prime sums. Averaging a prime statistic at the scales $xe^u$ against $\mu$ multiplies the contribution
of a zero $\rho$ by $F_\mu(\rho)$; on the critical line it is convenient to write
$\widehat\mu(\omega)=\int e^{i\omega u}\dd\mu(u)$, so that $F_\mu(\tfrac12+i\gamma)$ is the Fourier
transform of the tilted measure $e^{u/2}\dd\mu/\int e^{v/2}\dd\mu$ at $\gamma$. This section asks how
short such a positive average can be if it must annihilate prescribed zero contributions. It treats,
in turn: exact optima for finite sets of complex exponents, the central multiplicity that makes a
wide reflected band optimal, and the analytic dependence of the optima on the exponents, with their exact second variation at every strictly
positive mirror configuration; the
mean-frequency lower bound for arbitrary spectra and the exact quadratic cost of splitting a
multiple frequency asymmetrically; the logarithmic support, sharp up to a bounded error, of a filter
that cancels every zero up to a height $B$; one filter cancelling every noncentral zero, with the
exact supremal coefficient $\pi d/Q_*$ of its double-exponential decay; the resulting averages of
prime sums, among them an explicit compact kernel for all residue classes; and finite positive grids
of prime sums that read off the central multiplicity (\cref{fil:sec:grids}). All extremal statements
for prescribed frequencies are unconditional; \cref{src:thm:dgrh} enters only to put the zeros of
Dirichlet $L$-functions on one vertical line, where they become real Fourier frequencies of normalized
prime errors. Two frequencies in arbitrary ratio, progressions of frequencies and zero-free intervals
of positive trigonometric sums are in \cref{sec:zfg}; balancing and oscillation of prime errors at a
single centre are in \cref{sec:flu}.

\subsection{Exact optima for reflection-paired spectra}\label{fil:sec:reflect}
\sources{quadratic_spectra; sharp_positive_zero_filters; optimal_zero_filters; optimal_dirichlet_zero_filters; grh_cayley_and_filters; erdelyi_ellipse_and_proveit_extensions; ProveIt_research_2026-11-27; proveit_three_extensions_2026-11-27_1}

Fix distinct $s_1,\dots,s_r$ in the open upper half-plane and positive integers $m_j$, put
$n=\sum_jm_j$, and require of a compactly supported probability measure $\mu$
\begin{equation}\label{rec:eq:fjets}
 F_\mu^{(k)}(s_j)=\int u^ke^{s_ju}\dd\mu(u)=0\qquad(1\le j\le r,\ 0\le k<m_j).
\end{equation}
Since $\mu$ is real, the conjugate conditions hold automatically. Let $D(\boldsymbol s,\boldsymbol m)$
be the infimum of $\operatorname{diam}\supp\mu$ over such $\mu$. Translation multiplies $F_\mu$ by a
nowhere vanishing exponential, so it preserves \eqref{rec:eq:fjets} and the diameter. The feasible
set is nonempty: at $s=\beta+i\omega$ the masses at $0$ and $\pi/\omega$ in ratio $e^{\beta\pi/\omega}:1$
give a simple zero, and convolutions of such two-point measures give any prescribed orders. The
infimum is attained: translate a minimizing sequence into a fixed compact interval and take a weak
limit, using continuity of the test functions $u^ke^{s_ju}$ and of the support bound. A
\emph{certificate} on $[0,L]$ is a real function $g$ in the real span of the functions
$\Re(u^ke^{s_ju})$, $\Im(u^ke^{s_ju})$ ($k<m_j$) with $g\ge0$ on $[0,L]$; every admissible measure
supported on $[0,L]$ integrates $g$ to zero and therefore lives on the zero set of $g$.

\begin{theorem}[Reflection-paired spectra; \stU]\label{fil:thm:reflect}
Fix $h>0$. Suppose $0<\Im s_j<2\pi/h$ for all $j$ and that the multiset of targets, with
multiplicities $m_j$, is invariant under $s\mapsto\bar s+2\pi i/h$. Put
\[
 P(z)=\prod_j\bigl(z-e^{hs_j}\bigr)^{m_j}=\sum_{k=0}^np_kz^k .
\]
Then $P$ has real coefficients, $P(0)>0$, $P(1)>0$, and $D(\boldsymbol s,\boldsymbol m)\ge nh$.
Equality holds if and only if every $p_k\ge0$; in that case the unique minimizer up to translation is
\[
 \mu_*=\frac1{P(1)}\sum_{k=0}^np_k\,\delta_{kh}.
\]
If some $p_k<0$, then $D(\boldsymbol s,\boldsymbol m)>nh$. If $h\Im s_j\in(\pi/2,3\pi/2)$ for every
$j$, all coefficients are strictly positive and $\mu_*$ has $n+1$ atoms. No commensurability and no
common real part is assumed.
\end{theorem}

\begin{proof}
Put $q_j=-e^{hs_j}$ and $Q(z)=(-1)^nP(-z)=\prod_j(z-q_j)^{m_j}$. Because $0<\Im(hs_j)<2\pi$, the
$q_j$ lie in $\C\setminus(-\infty,0]$, and the reflection invariance makes their multiset closed under
conjugation. Real members are positive and the others form conjugate pairs; hence $P$ is real,
$Q(-t)=(-1)^n\abs{Q(-t)}$ for $t\ge0$, $P(0)=\prod_jq_j^{m_j}>0$ and $P(1)=\prod_j(1+q_j)^{m_j}>0$.
Distinct targets give distinct roots, since the imaginary parts of the $hs_j$ lie in an open interval of
length $2\pi$.

With the principal logarithm let
\[
 E(x)=\sum_j\Res_{z=q_j}\frac{z^{x-1}}{Q(z)},\qquad G(u)=(-1)^{n+1}\sin(\pi u/h)\,E(u/h).
\]
The residue formula, at repeated poles too, makes $E$ a combination of $x^ke^{x\log q_j}$ with
$k<m_j$, real for real $x$. Since $\log q_j=hs_j-i\pi$, multiplication by $\sin(\pi u/h)$ produces the
exponents $s_j$ and $s_j-2\pi i/h=\bar s_{j'}$, $j'$ the reflected partner, so $G$ is a real
combination of $\Re(u^ke^{s_ju})$ and $\Im(u^ke^{s_ju})$, $k<m_j$. For $0<x<n$ the residue theorem
on a keyhole about the negative axis (the circles contribute $O(R^{x-n})$ and $O(\epsilon^x)$, the jump
across the cut is $2i\sin(\pi(x-1))t^{x-1}/Q(-t)$) gives
\[
 E(x)=\frac{(-1)^{n+1}\sin(\pi x)}\pi\int_0^\infty\frac{t^{x-1}}{\abs{Q(-t)}}\dd t,\qquad
 G(hx)=\frac{\sin^2(\pi x)}\pi\int_0^\infty\frac{t^{x-1}}{\abs{Q(-t)}}\dd t .
\]
The integral is positive and finite, equal to $1/(\abs{Q(0)}x)+O(1)$ as $x\downarrow0$ and to
$1/(n-x)+O(1)$ as $x\uparrow n$. Hence $G\ge0$ on $[0,nh]$, with double zeros at the interior nodes
$h,\dots,(n-1)h$, simple zeros at the ends, no other zeros, and
\begin{equation}\label{fil:eq:slopes}
 G'(0)=\frac\pi{h\abs{Q(0)}},\qquad G'(nh)=-\frac\pi h .
\end{equation}
A feasible measure of diameter at most $nh$, translated into $[0,nh]$, integrates $G$ to zero and so
lives on $\set{0,h,\dots,nh}$. Its transform is $A(e^{hs})$ for a real polynomial $A$ of degree at
most $n$ with nonnegative coefficients and $A(1)=1$; since $s\mapsto e^{hs}$ has nonzero derivative,
$A$ vanishes to order $m_j$ at $e^{hs_j}$, so $A=P/P(1)$, including its positive endpoint
coefficients. Thus a feasible measure of diameter at most $nh$ exists precisely when every
$p_k\ge0$, it is unique, and its diameter is exactly $nh$. Attainment of the infimum gives the strict
inequality otherwise. If $h\Im s_j\in(\pi/2,3\pi/2)$ then $\Re e^{hs_j}<0$: real roots give factors
$z-r_j$ with positive coefficients and pairs give $z^2-2\Re(r_j)z+\abs{r_j}^2$, so every coefficient of
the product is strictly positive.
\end{proof}

The nonnegativity criterion can hold beyond the open band; at its boundary some coefficients
vanish and the theorem still applies, but the perturbation theory below needs strictly positive
masses. The following special cases were proved earlier by the same residue certificate. In each
of them one checks the reflection invariance and $p_k\ge0$ and applies \cref{fil:thm:reflect}; for a
common real part, \cref{fil:cor:pair,fil:cor:triple,fil:cor:cluster} are also equality cases of
\cref{fil:thm:meanbound} below.

\begin{example}[One frequency, several real parts; \stU]\label{rec:thm:fwidth}
Fix $\tau>0$, distinct reals $\beta_1,\dots,\beta_r$, positive integers $m_j$, $s_j=\beta_j+i\tau$,
$m=\sum m_j$ and $h=\pi/\tau$. Each target is fixed by $s\mapsto\bar s+2\pi i/h$, and
$e^{hs_j}=-e^{h\beta_j}$. So every solution of \eqref{rec:eq:fjets} has support diameter at least
$m\pi/\tau$, and up to translation exactly one has that diameter:
$\mu_*=\sum_{\ell=0}^mw_\ell\delta_{\ell h}$ with
$\sum_\ell w_\ell z^\ell=\prod_j\bigl((z+e^{h\beta_j})/(1+e^{h\beta_j})\bigr)^{m_j}$. The width bound alone
also follows from Obrechkoff's sector inequality \cite{EremenkoFryntov2016Remarks}.
\end{example}

\begin{example}[A block of consecutive harmonics; \stU]\label{rec:thm:hblock}
Fix integers $p,H\ge1$, $\tau>0$, distinct $\beta_j$ and $m_j\ge1$, and put $B=2p+H-1$,
$h=2\pi/(B\tau)$, $n=H\sum m_j$, $a_j=e^{h\beta_j}$, $\omega=e^{2\pi i/B}$. If a compactly supported
probability measure satisfies $F_\mu^{(k)}(\beta_j+i\ell\tau)=0$ for $p\le\ell\le p+H-1$, $k<m_j$ and all
$j$, its support has diameter at least $nh$, and the unique minimizer up to translation has masses
$w_t$ at $th$, $0\le t\le n$, with
\[
 \sum_tw_tz^t=\frac{P(z)}{P(1)},\qquad P(z)=\prod_j\Bigl[\prod_{\ell=p}^{p+H-1}(z-a_j\omega^\ell)\Bigr]^{m_j}.
\]
Indeed, the reflection $s\mapsto\bar s+2\pi i/h$ maps $\beta_j+i\ell\tau$ to $\beta_j+i(B-\ell)\tau$ and the
block $p\le\ell\le p+H-1$ to itself, and $0<\ell\tau<B\tau=2\pi/h$. The coefficients are positive:
with $\theta_s=(2s-H+1)\pi/B$, each inner factor is $a_j^HR(z/a_j)$ for
$R(y)=\prod_s(1+ye^{i\theta_s})=\sum c_ty^t$, and $(1+\alpha y)R(\omega y)=(1+\alpha\omega^Hy)R(y)$ gives
$c_t/c_{t-1}=\sin((H-t+1)\pi/B)/\sin(t\pi/B)>0$.

For $p=1$, $H=2$, $h=1$, $e^{\beta_1}=1/2$, $e^{\beta_2}=2$ with multiplicities $1,2$ the optimal width
is $6$, and the masses are
\[
 (w_0,\dots,w_6)=\tfrac1{343}(16,48,108,92,57,18,4),
\]
the coefficients of $(4z^2+2z+1)(z^2+2z+4)^2/343$; the exact check is recorded in
\begin{center}\artifact{quadratic_spectra/filters_example.txt}.\end{center}
\end{example}

\begin{example}[Two frequencies with equal vanishing orders; \stU]\label{fil:cor:pair}
Let $0<\alpha<\gamma\le3\alpha$, $\beta\in\R$ and $m\ge1$. The least diameter of a probability measure
with $F_\mu^{(k)}(\beta+i\alpha)=F_\mu^{(k)}(\beta+i\gamma)=0$ for $k<m$ is
\[
 D_m=\frac{4\pi m}{\alpha+\gamma}.
\]
With $h=2\pi/(\alpha+\gamma)$ and $a=e^{h\beta}$, the unique minimizer up to translation has masses at
$0,h,\dots,2mh$ given by the coefficients of
\[
 \frac{(z^2-2a\cos(h\alpha)z+a^2)^m}{(1-2a\cos(h\alpha)+a^2)^m}.
\]
Indeed, the two targets are reflected partners for this $h$, and $h\alpha+h\gamma=2\pi$, so
$P(z)=(z^2-2a\cos(h\alpha)z+a^2)^m$; since $\gamma\le3\alpha$ gives $h\alpha\in[\pi/2,\pi)$, we have
$\cos(h\alpha)\le0$, and every coefficient is nonnegative. The width is the fraction
$4\alpha\gamma/(\alpha+\gamma)^2$ of the width $m\pi(1/\alpha+1/\gamma)$ of two separately designed filters.
For $\gamma<3\alpha$ all $2m+1$ masses are positive; at $\gamma=3\alpha$ the odd coefficients vanish.
\end{example}

The formula is the $m$-fold convolution of the simple optimum, but \cref{fil:thm:reflect} shows that no
other positive measure, continuous or not, is shorter. For $m=2$ and $c=\cos(h\alpha)$, $\beta=0$, the weights
on $0,h,\dots,4h$ are $(1,-4c,2+4c^2,-4c,1)/(4(1-c)^2)$; for $\alpha=1$, $\gamma=2$ the width is $8\pi/3$
and the weights are $(1,2,3,2,1)/9$. Two-mode positive input shaping contains the low-order cases
\cite{CrainSinghoseSeering1996Derivation,SinghoseCrainSeering1997Convolved,LauPao2003Input,PaoSinghose1998Robust,Singhose1997Command}.
For the first two zeta ordinates the ratio is below $3$, so with $\beta=\tfrac12$ the least
logarithmic widths for order $m=1,2,3$ at both are $0.357438197$, $0.714876394$, $1.072314591$
(scale ratios $1.42966$, $2.04393$, $2.92214$), against $0.7434068$ for two separate order-two
filters; see
\begin{center}
\artifact{optimal_zero_filters/results.json}\\
\artifact{optimal_dirichlet_zero_filters/checks.json}
\end{center}

\begin{example}[A three-frequency band; \stU]\label{fil:cor:triple}
Let $0<\alpha<\gamma\le5\alpha$, $b=(\alpha+\gamma)/2$ and $\beta\in\R$. For simple vanishing at
$\beta+i\alpha$, $\beta+ib$, $\beta+i\gamma$ the least diameter is $6\pi/(\alpha+\gamma)$. With $h$, $a$ as in
\cref{fil:cor:pair} and $\theta=\pi-h\alpha$, the minimizer has the coefficients of
\[
 P(z)=(z+a)(z^2+2a\cos\theta\,z+a^2)=z^3+a(1+2\cos\theta)z^2+a^2(1+2\cos\theta)z+a^3 .
\]
Here the middle target is fixed by the reflection and the outer two are partners, and
$1+2\cos\theta\ge0$ is equivalent to $\gamma\le5\alpha$. For $\gamma<5\alpha$ every coefficient is
positive, at $\gamma=5\alpha$ only the endpoints remain, and for $\gamma>5\alpha$ this width is not
attained. The $m$th power of $P$ gives the least diameter $6\pi m/(\alpha+\gamma)$ for equal order $m$ at
all three, in the same band; for $\gamma>5\alpha$ the coefficient of degree $3m-1$ of $P^m$ is negative.
For $\cos\theta=-1/3$, $\beta=0$, the weights are $(3,1,1,3)/8$; here $\gamma/\alpha>3$, so
\cref{fil:cor:pair} alone does not give the result, and the angles need not be commensurate.
\end{example}

\begin{example}[Symmetric clusters; \stU]\label{fil:cor:cluster}
Let $\kappa>0$, and let the frequencies be $\kappa$ (order $m_0\ge0$) and $\kappa\pm\delta_j$ (order
$m_j\ge1$ each), $0<\delta_1<\dots<\delta_s\le\kappa/2$, with $N=m_0+2\sum_jm_j\ge1$ and $T=\pi/\kappa$.
With $h=T$ the reflection is $\lambda\mapsto2\kappa-\lambda$ on the frequencies, all of which lie in
$(0,2\kappa)$, and $\cos(\delta_jT)\ge0$, so \cref{fil:thm:reflect} applies to $s=i\lambda$: the least
diameter of a probability measure whose Fourier transform vanishes at each frequency to the given
order is $NT=N\pi/\kappa$, and the unique minimizer with left endpoint $0$ is
$\sum_{j=0}^Nw_j\delta_{jT}$ with
\[
 \sum_jw_jz^j=\frac{(1+z)^{m_0}}{2^{m_0}}\prod_{j=1}^s\Bigl(\frac{z^2+2\cos(\delta_jT)z+1}{2+2\cos(\delta_jT)}\Bigr)^{m_j},
\]
whose centred transform is
\[
 \Phi(\omega)=\cos(T\omega/2)^{m_0}\prod_{j=1}^s\Bigl(\frac{\cos(T\omega)+\cos(\delta_jT)}{1+\cos(\delta_jT)}\Bigr)^{m_j}.
\]
All weights are positive if $\delta_s<\kappa/2$, in particular if $N\delta_s<\kappa$. Cancelling the $N$
frequencies separately by two-point filters needs diameter $\pi\sum_\lambda1/\lambda$, strictly more as
soon as a noncentral pair occurs, by strict convexity of $1/x$. For the frequencies $10,12,14$ with
orders $1,2,1$ the width is $\pi/3$ and the weights are
$(\tfrac12-\tfrac{\sqrt3}4,\tfrac14,-\tfrac12+\tfrac{\sqrt3}2,\tfrac14,\tfrac12-\tfrac{\sqrt3}4)$.
\end{example}

\paragraph{Central multiplicity for wide bands.}
\Cref{fil:cor:pair,fil:cor:triple} show that a pair of frequencies spanning a ratio above $3$ loses the
reflection optimum, and that one central root restores it up to ratio $5$. In general, fix $\kappa>0$,
$h=\pi/\kappa$, frequencies $0<\alpha_j<\kappa$ ($1\le j\le m$, repetitions allowed) with their reflected
partners $2\kappa-\alpha_j$, and the central frequency $\kappa$ with multiplicity $r\ge0$; put
$\theta_j=h\alpha_j\in(0,\pi)$ and $N=r+2m$. By \cref{fil:thm:reflect} applied to $s=i\lambda$, the least
diameter of a probability measure whose Fourier transform vanishes at these frequencies to the prescribed
orders is at least $Nh$, with equality if and only if
\[
 P(z)=(1+z)^r\prod_{j=1}^m\bigl(1-2\cos\theta_j\,z+z^2\bigr)
\]
has nonnegative coefficients; the minimizer then has the masses $[z^\ell]P/P(1)$ at $\ell h$. A common real
part of the targets tilts these masses and changes nothing else. Factors with $\cos\theta_j\le0$ have
nonnegative coefficients; the question is how many central roots compensate the others.

\begin{theorem}[Central multiplicity criterion; \stU]\label{fil:thm:central}
If $a_1,\dots,a_m\in[0,1]$ and $r\ge0$ is an integer, then $(1+z)^r\prod_j(1-a_jz+z^2)$ has
nonnegative coefficients if and only if $r\ge\sum_ja_j$. Consequently, if $\pi/3\le\theta_j\le\pi/2$ for
all $j$, that is $\kappa/3\le\alpha_j\le\kappa/2$, the width $Nh$ is attained if and only if
$r\ge2\sum_j\cos\theta_j$, and the least central multiplicity with this property is
$\lceil2\sum_j\cos\theta_j\rceil$.
\end{theorem}

\begin{proof}
The coefficient of $z$ is $r-\sum_ja_j$, which gives necessity. For sufficiency induct on $m$; the cases
$m\le1$ are covered by the last step below. Fix all parameters but two fractional ones $u,v$ and vary these
with $u+v$ fixed. The product of their two factors is $(1+z^2)^2-(u+v)z(1+z^2)+uvz^2$, and the remaining
factors times $(1+z)^r$ have nonnegative coefficients by induction, their parameters summing to at most $r$.
Moving $(u,v)$ to the endpoint of its admissible segment that minimizes $uv$ makes one of them $0$ or $1$
and decreases every coefficient of the whole product weakly. Repeating, one reaches $b$ parameters equal to
$1$, some equal to $0$, and at most one fractional $a$, with a product dominated coefficientwise by the
original one. Pair $b$ binomials with the unit parameters by $(1+z)(1-z+z^2)=1+z^3$; if $0<a<1$, integrality
gives $r\ge b+1$ and $(1+z)(1-az+z^2)=1+(1-a)z+(1-a)z^2+z^3$. All remaining factors have nonnegative
coefficients. For the consequence take $a_j=2\cos\theta_j\in[0,1]$.
\end{proof}

For one pair and $r=1$ this is \cref{fil:cor:triple}. Below $\theta=\pi/3$ the linear coefficient no longer
decides. For a repeated pair let $Q_c(z)=1-2cz+z^2$ and let $R_m(c)$ be the least $r\ge0$ for which
$(1+z)^rQ_c^m$ has nonnegative coefficients; $R_m(c)=0$ for $c\le0$. For $0<c<1$ put
\[
 G_c(a)=\frac{1+c}{(a+1)(1-c)}\Bigl(\frac{a(1+c)}{2(a+1)}\Bigr)^a\qquad(a>0),
\]
and let $W$ be the inverse of $w\mapsto we^w$ on $(0,\infty)$.

\begin{theorem}[Phase transition for a repeated pair; \stU]\label{fil:thm:centralphase}
For $0<c\le\tfrac12$, $R_m(c)=\lceil2cm\rceil$ for every $m\ge1$. For $\tfrac12<c<1$ there is a unique
$a_c>c$ with $G_c(a_c)=1$, and
\[
 R_m(c)>2a_cm\quad(m\ge1),\qquad \tau(c)=\lim_{m\to\infty}\frac{R_m(c)}m=2a_c,\qquad
 \lim_{c\uparrow1}(1-c)\tau(c)=4W(e^{-1}).
\]
Also $R_1(c)=2\lceil(1-c)^{-1}\rceil-3$ for $0<c<1$.
\end{theorem}

The transition at $c=\tfrac12$ is at $\alpha=\kappa/3$, where the upper frequency is five times the lower; up
to that ratio the linear coefficient detects every failure, and beyond it a modulus test at a nonreal point
of the unit circle forces more central roots.

\begin{proof}
The first assertion is \cref{fil:thm:central} with $a_j=2c$. For $m=1$ and $D=r+2$ the $j$th coefficient divided
by $\binom Dj$ is $1-2(1+c)j(D-j)/(D(D-1))$, minimal at $j=\lfloor D/2\rfloor$; the two parities of $D$ give the
formula for $R_1$. So $R_m\le mR_1$ is finite, and $R_m$ is subadditive in $m$, whence
$\lim R_m/m=\inf R_m/m$.

A polynomial with nonnegative coefficients is bounded in modulus by its value at $\abs z$, so any admissible
$(r,m)$ satisfies, with $\lambda=r/m$,
\begin{equation}\label{fil:eq:modtest}
 \Bigl(\frac{\abs{1+z}}{1+\abs z}\Bigr)^\lambda\frac{\abs{Q_c(z)}}{Q_c(\abs z)}\le1\qquad(z\in\C).
\end{equation}
Conversely, if \eqref{fil:eq:modtest} holds at some $\lambda_0\ge2c$ and $A/B>\lambda_0$ is rational, then
$P=(1+z)^AQ_c^B$ satisfies $\abs{P(z)}<P(\abs z)$ off $[0,\infty)$, has endpoint coefficients $1$ and linear and
penultimate coefficients $A-2cB>0$; by De Angelis's theorem on positive powers \cite{DeAngelis2003Asymptotic}
all coefficients of $P^k$ are positive for large $k$, so $R_{Bk}\le Ak$, and subadditivity gives
$\lim R_m/m\le\lambda_0$.

Let $\tfrac12<c<1$, $d=1-c$, $b=(1+c)/2$, $z=xe^{it}$ with $x>0$, and $v=4x/(1+x)^2$, $s=v\sin^2(t/2)$,
$u=1-bv$; it suffices to take $0<s<1$. The squared binomial ratio in \eqref{fil:eq:modtest} is $1-s$, and
\[
 \frac{\abs{Q_c(z)}^2}{Q_c(x)^2}=1+\frac{2s}u+\frac{s(s-2d)}{u^2},\qquad \frac d2\le u\le1-bs,
\]
since $Q_c(x)/(1+x)^2=1-bv$ and $(\abs{Q_c(z)}^2-Q_c(x)^2)/(1+x)^4=s(s+2c-2bv)$; every pair in this range occurs.
One needs the supremum of the logarithm of this ratio divided by $-\log(1-s)$. At fixed $s$ it is a quadratic
in $1/u$; an interior maximum has $u=2d-s$ and value $1/(1-s/(2d))$, admissible for $0<s\le3d/2$, and
$\log(1-s/(2d))/\log(1-s)$ increases there (it is a weighted average of the increasing function
$(1-t)/(2d-t)$), so interior maxima are dominated by the endpoint $u=d/2$. At $u=1-bs$ (that is $z=-x$) the
unsquared ratio is $(1-(1-b)s)/(1-bs)$, whose logarithm is $\sum_j(b^j-(1-b)^j)s^j/j\le c\sum_js^j/j$; this
endpoint costs at most $\lambda=2c$. At $u=d/2$ (that is $\abs z=1$) the unsquared ratio is $\abs{1-2s/d}$, and
with $\lambda=2a$ the condition is $(1-s)^a\abs{1-2s/d}\le1$ on $[0,1]$. It holds for $s\le d/2$; beyond, the
left side has its maximum $G_c(a)$ at $s=(2+ad)/(2(a+1))$. $G_c$ decreases strictly from $(1+c)/(1-c)>1$ to
$0$, and $G_c(c)=(c/2)^c/(1-c)>1$ for $c>\tfrac12$, since $c\log(c/2)-\log(1-c)$ is strictly convex, vanishes
at $\tfrac12$ and has positive derivative there. This gives the unique root $a_c>c$, \eqref{fil:eq:modtest} at
$\lambda=2a_c$ and its failure below. At $\lambda=2a_c$ equality occurs at a nonreal point of the unit circle,
where a polynomial with positive constant and linear coefficients has strict triangle inequality; as
$2a_cm-2cm>0$, this excludes $r=2a_cm$, so $R_m(c)>2a_cm$, and the limit follows from the previous paragraph.
Finally, with $a=2w/d$, $G_{1-d}(2w/d)\to e^{-1-w}/w$ as $d\downarrow0$, a strictly decreasing function equal to
$1$ only at $w=W(e^{-1})$; monotonicity gives $da_c/2\to W(e^{-1})$.
\end{proof}

\begin{center}
\small
\begin{tabular}{@{}lrrrrrr@{}}
\toprule
$c$ & $R_1$ & $R_2$ & $R_5$ & $R_{10}$ & $R_{20}$ & $\tau(c)$\\
\midrule
$1/4$ & 1 & 1 & 3 & 5 & 10 & 0.5000\\
$1/2$ & 1 & 2 & 5 & 10 & 20 & 1.0000\\
$3/5$ & 3 & 5 & 10 & 17 & 33 & 1.5449\\
$3/4$ & 5 & 9 & 19 & 35 & 67 & 3.1998\\
$9/10$ & 17 & 27 & 57 & 106 & 206 & 9.8688\\
\bottomrule
\end{tabular}
\end{center}
The entries are exact integer computations, reproduced by
\begin{center}\artifact{erdelyi_ellipse_and_proveit_extensions/02_positive_filters/check_filters.py};\end{center}
the last column is $2a_c$ rounded. At $c=\tfrac34$, $m=1$ the linear coefficient would allow $r=2$, while the minimum is $5$.

\begin{corollary}[Lorentz degree of a centred quadratic power; \stU]\label{fil:cor:lorentz}
Let $d_L(p)$ be the least $D$ for which $p(X)$ is a nonnegative combination of $(1-X)^j(1+X)^{D-j}$,
$0\le j\le D$. For $0<\eps<1$ and $m\ge1$,
\[
 d_L\bigl((X^2+\eps^2)^m\bigr)=2m+R_m\Bigl(\frac{1-\eps^2}{1+\eps^2}\Bigr),
\]
which is $2m+\lceil2m(1-\eps^2)/(1+\eps^2)\rceil$ for $1/\sqrt3\le\eps<1$, and
\[
 \lim_{\eps\downarrow0}\eps^2\lim_{m\to\infty}\frac{d_L((X^2+\eps^2)^m)}m=2W(e^{-1})=0.556929\ldots .
\]
\end{corollary}

\begin{proof}
With $X=(1-z)/(1+z)$, $d_L(p)\le D$ for $D\ge\deg p$ exactly when $(1+z)^Dp((1-z)/(1+z))$ has nonnegative
coefficients (multiplying by $1+z$ raises the degree), and $(1+z)^2((1-z)^2/(1+z)^2+\eps^2)=(1+\eps^2)Q_c(z)$ with
$c=(1-\eps^2)/(1+\eps^2)$, which is at most $\tfrac12$ exactly when $\eps\ge1/\sqrt3$. Apply
\cref{fil:thm:centralphase}, with $1-c\sim2\eps^2$.
\end{proof}

Erd\'elyi bounded this degree above and below by multiples of $m/\eps^2$ \cite{Erdelyi1991Estimates}; the
corollary gives the exact first-order constant in the centred case. For $c>\tfrac12$ the size of
$R_m(c)-m\tau(c)$ is not known, and for distinct angles below $\pi/3$ different evaluation points constrain
different quadratic factors.

\paragraph{The mean-frequency bound without symmetry.}
For real frequencies the lower bound of \cref{fil:cor:cluster} needs no reflection symmetry: the mean
frequency alone bounds the width, and the equality cases are exactly the strictly or weakly positive
mirror configurations of \cref{fil:thm:reflect}. The proof rests on an identity for exponential sums.

\begin{lemma}[Strip identity; \stU]\label{fil:lem:strip}
Let $0=t_0<t_1<\dots<t_M=L$ and let $w_0,\dots,w_M\ge0$ be weights of total mass one with $w_0,w_M>0$.
Put $G(s)=\sum_jw_je^{t_js}$. For $T>0$,
\[
 J(T):=\int_\R\log\frac{G(x)}{\abs{G(x+iT)}}\dd x=\frac{LT^2}2-2\pi\sum_{0<\Im\rho<T}(T-\Im\rho),
\]
the sum running over the zeros $\rho$ of $G$ with multiplicity. The integrand is nonnegative, so
$J(T)\ge0$.
\end{lemma}

\begin{proof}
Each end of the exponential sum has a single dominant term uniformly in $\Im s$, so the integrand decays
exponentially in $\abs x$, the zeros of $G$ have bounded real parts, and the logarithmic singularities are
integrable; one may first truncate the $x$-integral outside the zero-containing strip and then pass to the
limit. In distributions $\Delta\log\abs G=2\pi\sum_\rho\delta_\rho$, by local factorization and
$\Delta\log\abs s=2\pi\delta_0$. Integrating the $xx$-derivative in $x$ gives the difference $L-0$ of the
endpoint slopes of $\log\abs G$, so $J''(T)=L-2\pi\sum_\rho\delta_{\Im\rho}(T)$ for $T>0$. Integration on
$[0,T]$ gives the identity, since $J(0)=0$, and $J'(0)=0$ by conjugation ($J$ is even) and since $G>0$ on
the real axis. Nonnegativity is the triangle inequality $\abs{G(x+iT)}\le G(x)$.
\end{proof}

This is the exponential-sum analogue of the integrated angular inequality of Eremenko and Fryntov
\cite[Thm.~1]{EremenkoFryntov2016Remarks} for polynomials with nonnegative coefficients.

\begin{theorem}[Mean-frequency bound and its equality cases; \stU]\label{fil:thm:meanbound}
Let $\lambda_1,\dots,\lambda_N>0$ be a multiset of frequencies, $\kappa=N^{-1}\sum_j\lambda_j$, and let
$B(\boldsymbol\lambda)$ be the least diameter of the support of a probability measure whose Fourier
transform vanishes at each $\lambda_j$ to the order of its repetition. Then
\[
 \kappa B(\boldsymbol\lambda)\ge N\pi .
\]
Put $h=\pi/\kappa$. Equality holds if and only if every $\lambda_j<2\kappa$, the multiset is invariant under
$\lambda\mapsto2\kappa-\lambda$, and the polynomial $A(z)=\prod_j(z-e^{ih\lambda_j})$ has nonnegative
coefficients; then the coefficients of $A(z)/A(1)$ are the masses of the optimizer on $0,h,\dots,Nh$, unique
up to translation.

More generally, let $s_1,\dots,s_N$ be targets in the open upper half-plane, listed with multiplicity and
with arbitrary real parts, put $\kappa=N^{-1}\sum_j\Im s_j$ and $h=\pi/\kappa$, and let $D$ be the least
width in \eqref{rec:eq:fjets}. Then $D\ge N\pi/\kappa$, with equality if and only if $\Im s_j<2\kappa$ for
every $j$, the multiset is invariant under $s\mapsto\bar s+2i\kappa$, and $P(z)=\prod_j(z-e^{hs_j})$ has
nonnegative coefficients; the optimizer is then the measure of \cref{fil:thm:reflect}, unique up to
translation.
\end{theorem}

\begin{proof}
An optimum exists, as noted before \cref{fil:thm:reflect}. For a feasible measure the real and imaginary jet
vector of the constraints is continuous on its support, and its mean lies in the convex hull of its values
there, so Carath\'eodory's theorem replaces the measure by a finite one with the same constraints and no larger
diameter. Translate such an optimal finite measure to start at $0$; it has atoms at both ends of $[0,L]$,
$L$ the optimal width, and $G(s)=\int e^{st}\dd\mu(t)$ satisfies the hypotheses of \cref{fil:lem:strip}. Its zeros
include the targets $s_j$ with their orders; the first statement is the case $s_j=i\lambda_j$. At $T=2\kappa$ the
zero sum is at least $\sum_j(T-\Im s_j)_+\ge NT-\sum_j\Im s_j=N\kappa$, so $J(T)\ge0$ gives $2\kappa^2L\ge2\pi N\kappa$,
which is the bound. (It is also the case $m=n$ of \cref{zfg:prop:prefix}.)

At equality, $J(T)=0$ and $\Im s_j\le T$ for every $j$. Then the triangle inequality is an equality at every
real $x$, by continuity (a zero of $G$ on the line $\Im s=T$ would make the integral positive). At $x=0$ all
phases $e^{it_jT}$ equal $e^{it_0T}=1$, so every node lies in $h\Z$ and $G(x+iT)=G(x)>0$; this excludes
$\Im s_j=T$. Now $G(s)=A_0(e^{hs})$ for a real polynomial $A_0$ of degree exactly $N$ with nonnegative
coefficients and $A_0(1)=1$. The map $s\mapsto e^{hs}$ is injective on the strip $0<\Im s<T$, which contains the
targets, and locally invertible, so $A_0$ vanishes at the $e^{hs_j}$ to the prescribed orders, and $A_0=P/P(1)$.
Real coefficients make the root multiset closed under conjugation; since $\overline{e^{hs}}=e^{h(\bar s+2i\kappa)}$,
this is the invariance under $s\mapsto\bar s+2i\kappa$, and the leading coefficient and the masses being
nonnegative give nonnegative coefficients of $P$. Conversely, under these three conditions the targets satisfy
the hypotheses of \cref{fil:thm:reflect} with $n=N$ and this $h$, which gives a minimizer of width $Nh$, unique
up to translation. For $s_j=i\lambda_j$ the three conditions are those of the first statement, with $P=A$.
\end{proof}

Thus a cluster that is not mirror symmetric about its mean always costs more than $N\pi/\kappa$, and within the
mirror locus the boundary of the equality region is the coefficient nonnegativity of the explicit polynomial
$A$. The bound does not depend on the real parts of the targets; attaining it does. It contains the lower
bound of \cref{fil:thm:reflect}, since reflection invariance forces $\kappa=\pi/h$. For $N=2$ the criterion is
$\lambda_2\le3\lambda_1$, as in \cref{fil:cor:pair}.

For prime sums, let $W$ be smooth with compact support in the interval $(1,2)$, and put
$S_W(x)=\sum_n\Lambda(n)W(n/x)$. A zeta zero contributes $-x^\rho\widetilde W(\rho)$ to $S_W$, with
$\widetilde W$ the Mellin transform, and a positive average of $S_W(xe^u)$ multiplies this by
$F_\mu(\rho)$. For $\rho=\tfrac12+i\gamma$, $\gamma>0$, and $h=\pi/\gamma$, the average
\[
 \frac{e^{h/2}S_W(x)+S_W(e^hx)}{1+e^{h/2}}
\]
cancels the conjugate pair; its $m$-fold iterate has a zero of order $m$ at $\rho$, logarithmic span
$mh$, minimal by \cref{rec:thm:fwidth}, and multiplier $O(\abs\eta^m)$ at a zero displaced by $\eta$.
A zero of order $m$ also annihilates $x^\rho(\log x)^k$ for $k<m$, since
$\int(\log x+u)^ke^{\rho u}\dd\mu(u)=0$, which matters for differentiated explicit formulas and for
uncertain zero locations; mere cancellation of a multiple zero needs only a simple zero of $F_\mu$.

\subsection{Independent perturbations, splitting and stability}\label{fil:sec:perturb}
\sources{sharp_positive_zero_filters; optimal_zero_filters; optimal_dirichlet_zero_filters; quadratic_spectra; gram-euler-and-zero-filters; grh_cayley_and_filters; grh_hidden_forms_and_filters_2026-11-27; ProveIt_research_2026-11-27; deep_padic_roots_2026-11-27; proveit_three_extensions_2026-11-27_1}

Exact reflection pairing is not natural for actual zero ordinates. It can be removed locally. Write
$\sigma=(s_1,\dots,s_r)$ with fixed multiplicities and $Y_\sigma(u)\in\R^{2n}$ for the real and
imaginary parts of $u^ke^{s_ju}$, $k<m_j$.

\begin{theorem}[Analytic continuation of the optimum; \stU]\label{fil:thm:perturb}
Let the base targets satisfy \cref{fil:thm:reflect} with all coefficients of $P$ strictly positive.
There is an open neighbourhood $U$ of the base, in the $2r$ real coordinates of the targets, in which
the minimum-width problem \eqref{rec:eq:fjets} has a unique minimizer up to translation; with left
endpoint $0$ it is
\[
 \mu_\sigma=\sum_{i=0}^nw_i(\sigma)\delta_{u_i(\sigma)},\qquad0=u_0<u_1<\dots<u_n,\quad w_i>0,
\]
and the weights, the nodes and the width $L(\sigma)=u_n(\sigma)$ are real analytic on $U$. A real
analytic family of certificates is positive off the nodes, with double interior and simple endpoint
zeros. The targets may move independently in both directions; no reflection symmetry is needed at
the perturbed point.

The same holds in coordinates in which targets may collide or split. Let $\kappa=\pi/h$, translate the
real parts of the base to mean $0$ (exponential tilting changes no width), and describe a multiset of $n$
targets with mean $i\kappa$ by the coefficients $c=(c_2,\dots,c_n)\in\C^{n-1}$ of its centred polynomial
\[
 P_c(y)=\prod_{s}\bigl(y-(s-i\kappa)\bigr)=y^n+\sum_{r=2}^nc_ry^{n-r},
\]
the product running over the multiset. On a neighbourhood of the base coefficients $c^0$ in the real
coordinates $(\Re c,\Im c)$ the minimizer is again unique up to translation, has $n+1$ positive atoms,
and has nodes, weights and width real analytic in $c$, including where $P_c$ has repeated roots. Tilting
and positive rescaling of the targets remove the restriction on the mean.
\end{theorem}

\begin{proof}
The unknowns $(w_0,\dots,w_n,u_1,\dots,u_n)$ satisfy the $2n+1$ real equations
\begin{equation}\label{fil:eq:system}
 \sum_iw_i=1,\qquad\sum_iw_iY_\sigma(u_i)=0 .
\end{equation}
At the base ($u_i=ih$, $w_i=p_i/P(1)$) let a variation lie in the kernel of the Jacobian. The variation
of the transform is $\delta F(s)=A(e^{hs})+sB(e^{hs})$ with real polynomials
$A(z)=\sum_i\delta w_iz^i$, $B(z)=\sum_iw_i\delta u_iz^i$, $A(1)=0$, $B(0)=0$. All target jets of
$\delta F$ vanish; conjugating the equations of the reflected partner, so do those of
$\delta F(s-2\pi i/h)$. The difference is $(2\pi i/h)B(e^{hs})$, so $P\mid B$; as $\deg B\le n$ and
$B(0)=0\ne P(0)$, $B=0$. Then $P\mid A$ and $A(1)=0\ne P(1)$ give $A=0$, and positivity of the $w_i$
gives $\delta u_i=0$. The real analytic implicit function theorem yields an analytic solution with
positive weights and ordered nodes.

For global optimality, the columns $Y(u_i)$ ($0\le i\le n$) and $Y'(u_i)$ ($1\le i<n$) span a hyperplane
of $\R^{2n}$: the weight relation is a dependence, and adding $Y'(u_n)$ gives rank $2n$ by the full
Jacobian. Let $v_\sigma$ be its normal with $v_\sigma\cdot Y'(u_n)=-1$, analytic in $\sigma$, and
$g_\sigma=v_\sigma\cdot Y_\sigma$. It vanishes at the nodes, with vanishing derivative at interior
nodes; at the base it is a positive multiple of the residue certificate $G$, since the normal space is
one-dimensional and $G'(nh)<0$. The strict signs of $g'$ at the ends and of $g''$ at interior nodes,
and positivity off neighbourhoods of the nodes, persist under perturbation, so after shrinking $U$,
$g_\sigma\ge0$ on $[0,L(\sigma)]$ with exactly the nodes as zeros. A feasible measure of width at most
$L(\sigma)$ is forced onto the nodes, where \eqref{fil:eq:system} fixes the weights (the weight columns
with the mass row are independent), and both endpoint weights are positive.

In coefficient coordinates fix a contour $C$ about the roots of $P_{c^0}$ and put
\[
 (\mathcal R_cf)(y)=\frac1{2\pi i}\oint_C\frac{f(z)}{P_c(z)}\,\frac{P_c(z)-P_c(y)}{z-y}\dd z
\]
for entire $f$. By Cauchy's formula this is $f(y)$ minus $P_c(y)$ times a holomorphic function, so it is the
polynomial of degree below $n$ that agrees with $f$ at the roots of $P_c$ to their multiplicities; it is
holomorphic in $c$, and \eqref{rec:eq:fjets} is equivalent to the vanishing of the $n$ complex coefficients of
$\mathcal R_c[F_\mu(i\kappa+\cdot)]$, also at collisions. These replace the jet equations in
\eqref{fil:eq:system}. A variation in the kernel of the Jacobian at the base now makes
$A(-e^{hy})+(y+i\kappa)B(-e^{hy})$ divisible by $P_{c^0}(y)$, with real $A$, $B$ as before. As $P_{c^0}$ is real,
conjugating coefficients gives the same with $-i\kappa$, and subtraction gives $P_{c^0}(y)\mid B(-e^{hy})$; since
$y\mapsto-e^{hy}$ is injective with nonzero derivative on $\abs{\Im y}<\kappa$, this means $P\mid B$, and the
argument concludes as before without using distinct roots. The certificate is built in the same way from
the real and imaginary parts $Y_c(u)$ of the coefficients of $\mathcal R_c[e^{(i\kappa+\cdot)u}]$.
\end{proof}

\begin{proposition}[Sensitivity of the minimum; \stU]\label{fil:prop:sense}
On $U$, for an infinitesimal change $\dot\sigma$ of the targets,
\[
 \dot L=-\frac{\sum_{i=0}^nw_i\,v_\sigma\cdot(\partial_\sigma Y_\sigma)(u_i)[\dot\sigma]}{w_n\,g_\sigma'(L)} .
\]
\end{proposition}

\begin{proof}
Differentiate the moment equation of \eqref{fil:eq:system} and take the scalar product with $v_\sigma$:
the weight terms and the interior node terms vanish, and $u_0$ is fixed.
\end{proof}

At a mirror configuration the first-order change is universal.

\begin{theorem}[Differential at a mirror configuration; \stU]\label{fil:thm:differential}
In \cref{fil:thm:perturb} write $s_j=s_j^0+\Delta\beta_j+i\Delta\gamma_j$. Then
\[
 D(\boldsymbol s)=nh-\frac{h^2}\pi\sum_jm_j\Delta\gamma_j+O_{\boldsymbol s^0,\boldsymbol m}\bigl(\norm{\Delta\boldsymbol s}^2\bigr).
\]
All derivatives in the real parts vanish, even when they vary independently.
\end{theorem}

\begin{proof}
Put $\kappa_0=\pi/h$ and $a=n^{-1}\sum_jm_j\Delta\gamma_j$. At the base $\sum_jm_j\Im s_j^0=n\kappa_0$, so the targets
divided by $1+a/\kappa_0$ have weighted mean imaginary part $\kappa_0$. On that hyperplane $D\ge n\pi/\kappa_0=nh$ by
\cref{fil:thm:meanbound}, with equality at the base, and $D$ is real analytic there by \cref{fil:thm:perturb}; so its
derivative along the hyperplane vanishes at the base. Positive scaling of the targets divides the width by the
scale, whence $D(\boldsymbol s)=nh(1+a/\kappa_0)^{-1}+O(\norm{\Delta\boldsymbol s}^2)=nh-(h^2/\pi)\,na+O(\norm{\Delta\boldsymbol s}^2)$.
Changes of the real parts do not move the mean imaginary part.
\end{proof}

Perturbations keeping the weighted mean imaginary part fixed thus have no first-order cost; their
second-order cost is determined below (\cref{fil:thm:mirrorhessian,fil:cor:mirrorlabel}).

\begin{example}[Detuned harmonics and slices off the mirror locus]\label{fil:ex:detune}
For frequencies $(1+\epsilon_1,2+\epsilon_2,3+\epsilon_3)$ (so $h=\pi/2$, $P=1+z+z^2+z^3$),
$L=\tfrac{3\pi}2-\tfrac\pi4(\epsilon_1+\epsilon_2+\epsilon_3)+O(\norm\epsilon^2)$; nearby rationally
independent triples thus have a unique analytic global optimizer. For $b$ close to $2$ the frequencies
$(1,b,3)$ have exact optimal width $\pi+\pi/b$, also for irrational $b$: the convolution of the centred
equal-mass filters at $1$ and $b$ annihilates all three (the filter at $1$ also kills $3$), has nodes
$\pm(\pi/2\pm\pi/(2b))$ of mass $\tfrac14$, equals the base optimizer at $b=2$ and solves
\eqref{fil:eq:system}, so it is the optimizer of \cref{fil:thm:perturb}. In units $h=1$ the base exponents
$i(2\pi/3,\pi,4\pi/3)$ give $P=(z+1)(z^2+z+1)$ and masses $(1,2,2,1)/6$; for
$a_1=2\pi i/3+\epsilon/3$, $a_2=i(\pi+\epsilon)$, $a_3=4\pi i/3$ and small real $\epsilon$ the unique
shortest filter has width $3-\epsilon/\pi+O(\epsilon^2)$ (numerically $(D-3+\epsilon/\pi)/\epsilon^2\approx0.0429$,
\artifact{optimal_zero_filters/results.json}).
\end{example}

\paragraph{Second order at a mirror configuration.}
By \cref{fil:thm:differential,fil:thm:meanbound} the width is stationary, and minimal, along perturbations
that keep the weighted mean imaginary part fixed. Its second variation there is an explicit positive
quadratic form whose kernel is exactly the tangent space of the reflection locus. Use the coefficient
coordinates of \cref{fil:thm:perturb}: $\kappa=\pi/h$, a base $P_{c^0}$ with real coefficients whose roots $y$,
listed with multiplicity, lie in $\abs{\Im y}<\kappa$, and
\[
 P(z)=\prod_y\bigl(z+e^{hy}\bigr)=\sum_{k=0}^np_kz^k,\qquad p_k>0\quad(0\le k\le n),
\]
the polynomial of \cref{fil:thm:reflect} for the targets $i\kappa+y$. Let $D(c)$ be the least width at the
roots of $P_c$ shifted by $i\kappa$. This multiset is reflection-paired exactly when $\Im c=0$, and for real $c$
near $c^0$ the coefficients of its polynomial stay positive, so that $D(c)=nh$ there. Let
$G_c(s)=\sum_kw_k(c)e^{t_k(c)s}$, $t_0=0$, be the transform of the minimizer, and $w_k=p_k/P(1)$ its base weights.

\begin{lemma}[Exact deficit; \stU]\label{fil:lem:mirrordefect}
After shrinking the neighbourhood, $G_c$ has exactly the $n$ prescribed zeros in $0<\Im s<2\kappa$, and with
$\varpi_k(x,c)=w_k(c)e^{xt_k(c)}/G_c(x)$
\begin{equation}\label{fil:eq:mirrordefect}
 D(c)-nh=\frac1{2\kappa^2}\int_\R-\log\Bigl\lvert\sum_{k=0}^n\varpi_k(x,c)e^{2i\kappa t_k(c)}\Bigr\rvert\dd x .
\end{equation}
\end{lemma}

\begin{proof}
On a small closed neighbourhood the endpoint weights and the gaps between adjacent nodes are bounded below.
For $\Re s>R$ division by the last term bounds the sum of the absolute values of the others by $Ce^{-gR}$, and
the first term does the same for $\Re s<-R$; so for large $R$ no zero with $0\le\Im s\le2\kappa$ lies outside
$\abs{\Re s}\le R$, uniformly in $c$. At the base $G_{c^0}(s)=P(e^{hs})/P(1)$ has exactly $n$ zeros in the strip,
one logarithm of each root of $P$, and none on its edges, where $G_{c^0}(x)=G_{c^0}(x+2i\kappa)>0$. Rouch\'e's
theorem on the rectangle keeps the count, and the $n$ prescribed zeros account for all of them; their imaginary
parts sum to $n\kappa$. \Cref{fil:lem:strip} at $T=2\kappa$ gives $J(2\kappa)=2\kappa^2(D(c)-nh)$, and dividing the
integrand by $G_c(x)$ gives \eqref{fil:eq:mirrordefect}.
\end{proof}

\begin{theorem}[Width Hessian at a mirror configuration; \stU]\label{fil:thm:mirrorhessian}
In this setting the following hold on a neighbourhood of $c^0$.
\begin{enumerate}[label=\textup{(\roman*)}]
\item Let $V$ be a real polynomial of degree at most $n-2$, let $P_{c(\eps)}=P_{c^0}+i\eps V$, and write the nodes
as $t_k=kh+\eps\sigma_k+O(\eps^2)$. Then $\sigma_0=\sigma_n=0$, and $B_V(z)=\sum_kw_k\sigma_kz^k$ is the unique real
polynomial of degree below $n$ with
\begin{equation}\label{fil:eq:hermite}
 \kappa B_V(-e^{hy})\equiv V(y)H(y)\pmod{P_{c^0}(y)},\qquad H(y)=\frac{P(-e^{hy})}{P(1)P_{c^0}(y)},
\end{equation}
where $H$ is holomorphic and nonzero at the roots of $P_{c^0}$; also $B_V(0)=0$. The map
$V\mapsto(\sigma_1,\dots,\sigma_{n-1})$ is a linear isomorphism. If the roots $y_j$ of $P_{c^0}$ are simple, then
with $z_j=-e^{hy_j}$
\[
 B_V(z)=\frac h{\kappa P(1)}\sum_{j=1}^n\frac{z_jV(y_j)}{P_{c^0}'(y_j)}\,\frac{P(z)}{z-z_j}.
\]
\item Along every real analytic path of multisets with mean imaginary part $\kappa$ through the base, with node
velocities $\sigma_k$, the width is $D=nh+\eps^2Q+O(\eps^3)$, where, with $\varpi_x(k)=p_ke^{khx}/P(e^{hx})$,
\begin{equation}\label{fil:eq:mirrorvar}
 Q=\int_\R\operatorname{Var}_{\varpi_x}(\sigma)\dd x
 =\frac1h\sum_{0\le j<k\le n}p_jp_k(\sigma_j-\sigma_k)^2\int_0^\infty\frac{z^{j+k-1}}{P(z)^2}\dd z .
\end{equation}
For the path of \textup{(i)} write $Q=Q(V)$; then $Q(V)>0$ for $V\ne0$.
\item In the coordinates $(\Re c,\Im c)$ the Hessian of $D$ at $c^0$ vanishes on the real and the mixed
blocks and is $2Q$ on the imaginary block, so its kernel is exactly the tangent space $\Im c=0$ of the
reflection locus. There are $0<C_1<C_2$ with
\[
 C_1\norm{\Im c}^2\le D(c)-nh\le C_2\norm{\Im c}^2 ,
\]
so that $D(c)=nh$ exactly when $\Im c=0$, and a real analytic change of coordinates
$(\Re c,\Im c)\mapsto(\Re c,\eta)$, with $\eta=0$ exactly on $\Im c=0$, gives
\[
 D(c)=nh+\norm\eta^2 .
\]
\end{enumerate}
\end{theorem}

\begin{proof}
(i) Along the path the mean stays $i\kappa$, so $D\ge nh=D(c^0)$ by \cref{fil:thm:meanbound}, and analyticity gives
$\dot D=0$, that is $\sigma_n=0$; anchoring gives $\sigma_0=0$. The transform of the minimizer at $i\kappa+y$ is
divisible by $P_{c(\eps)}(y)$, and at the base it is $P(-e^{hy})/P(1)=P_{c^0}(y)H(y)$. Differentiating at $\eps=0$,
with $W(z)=\sum_k\dot w_kz^k$,
\[
 W(-e^{hy})+(y+i\kappa)B_V(-e^{hy})\equiv iV(y)H(y)\pmod{P_{c^0}(y)} .
\]
$H$ is holomorphic and nonzero at the roots of $P_{c^0}$, since $y\mapsto-e^{hy}$ has nonzero derivative. All of
$P_{c^0}$, $H$, $W$, $B_V$, $V$ are real, so conjugating coefficients replaces $i\kappa$ and $i$ by $-i\kappa$ and
$-i$, and subtraction gives \eqref{fil:eq:hermite}. If a real polynomial $B$ of degree below $n$ has
$P_{c^0}(y)\mid B(-e^{hy})$, then $P\mid B$ and $B=0$; this gives uniqueness. If $B_V=0$, then $P_{c^0}\mid VH$,
hence $P_{c^0}\mid V$ and $V=0$; so the map is injective between spaces of dimension $n-1$. For simple roots,
$H(y_j)=hz_jP'(z_j)/(P(1)P_{c^0}'(y_j))$, and Lagrange interpolation at the $z_j$ gives the displayed formula;
its value at $0$ is $-hP(0)(\kappa P(1))^{-1}\sum_jV(y_j)/P_{c^0}'(y_j)=0$, as $\deg V\le n-2$.

(ii) For a probability vector $\varpi$ and real $\sigma$,
$-\log\abs{\sum_k\varpi_ke^{2i\kappa\eps\sigma_k}}=2\kappa^2\eps^2\operatorname{Var}_\varpi(\sigma)+O(\eps^3)$. At the base all
phases $e^{2i\kappa kh}$ equal $1$, so the variation of the weights and of $\varpi$, and the second-order terms of
the nodes, do not enter the coefficient of $\eps^2$, the total mass being one. On each tail one endpoint carries
all but $O(e^{-g\abs x})$ of the mass, uniformly near the base; after factoring out its phase the integrand of
\eqref{fil:eq:mirrordefect} has three parameter derivatives bounded by $C(1+\abs x)^3e^{-g\abs x}$, and on the central
compact interval the sum inside the modulus stays away from $0$. Differentiating under the integral in
\cref{fil:lem:mirrordefect} (which uses only the mean imaginary part) gives the first form of $Q$, and
$\operatorname{Var}_\varpi(\sigma)=\sum_{j<k}\varpi_j\varpi_k(\sigma_j-\sigma_k)^2$ with $z=e^{hx}$ gives the second; the integrals
converge since $1\le j+k\le2n-1$. All edge weights are positive, so $Q(V)=0$ forces every $\sigma_k$ to equal
$\sigma_0=0$, whence $B_V=0$ and $V=0$.

(iii) Write $c=E+iO$ with $E$, $O$ real. Since $D(E,0)=nh$ and $D\ge nh$, $\partial_OD(E,0)=0$, and Taylor's formula
in $O$ gives $D(E,O)-nh=O^{\mathsf T}M(E,O)O$ with $M(E,O)=\int_0^1(1-t)\,\partial^2_{OO}D(E,tO)\dd t$ real analytic and
symmetric. So the real and mixed blocks of the Hessian vanish at $c^0$, and by \textup{(ii)} $M(c^0)$ is the matrix of
$Q$, which is positive definite. Its eigenvalues stay between two positive constants nearby, which gives the
two-sided bound. The derivative $Z\mapsto XZ+ZX$ of $X\mapsto X^2$ is invertible at a positive definite $X$ (in an
eigenbasis it multiplies the entries by $\lambda_j+\lambda_k>0$), so $M$ has a real analytic positive definite
square root, and $\eta=M(E,O)^{1/2}O$ is the required coordinate.
\end{proof}

The coefficients, not individual roots, are the regular coordinates. At a collision a root displacement
changes the coefficients only at high order: along a root path on which $\Im c$ first becomes nonzero at
order $r$ the width excess has exact order $2r$, as in \cref{fil:cor:ray}. At the coalesced base
$P_{c^0}=y^n$, $P=(1+z)^n$, real frequencies near $\kappa$ have $c_r\in i^r\R$, so their odd coefficients are the
imaginary ones; \cref{fil:thm:hessian} below evaluates $Q$ there in closed form. For labelled targets the
velocity map and the full Hessian are explicit.

\begin{corollary}[Labelled targets; \stU]\label{fil:cor:mirrorlabel}
Let the base satisfy \cref{fil:thm:reflect} with all coefficients of $P$ positive, with distinct targets $s_j^0$ of
multiplicities $m_j$, reflected partners $j^*$, $\kappa_0=\pi/h$ and $z_j=e^{hs_j^0}$, and move the targets as
$s_j^0+\eps v_j$ ($v\in\C^r$) with fixed multiplicities. Put $d_j=(v_j-\overline{v_{j^*}})/i$.
\begin{enumerate}[label=\textup{(\roman*)}]
\item If $\sum_jm_j\Im v_j=0$, the node velocities $u_k$ satisfy $u_0=u_n=0$ and
\[
 \sum_{k=0}^nw_ku_kz^k=-\frac h{2\kappa_0P(1)}\sum_{j=1}^rm_jd_jz_j\frac{P(z)}{z-z_j},
\]
and $D=nh+\eps^2Q(v)+O(\eps^3)$, with $Q(v)$ given by \eqref{fil:eq:mirrorvar} for $\sigma=u$. Moreover $Q(v)=0$ if and
only if $v_j=\overline{v_{j^*}}$ for every $j$.
\item For arbitrary $v$ put $a=n^{-1}\sum_jm_j\Im v_j$ and $\hat v=v-(a/\kappa_0)\boldsymbol s^0$. Then
\[
 D(\boldsymbol s^0+\eps v)=nh-\frac{nha}{\kappa_0}\eps+\Bigl(\frac{nha^2}{\kappa_0^2}+Q(\hat v)\Bigr)\eps^2+O(\eps^3),
\]
so the second derivative along every line is nonnegative and vanishes exactly when $v_j=\overline{v_{j^*}}$ for
every $j$.
\item With $\kappa=n^{-1}\sum_jm_j\Im s_j$ and
$\mathcal E(\boldsymbol s)^2=\sum_jm_j\abs{s_j-\overline{s_{j^*}}-2i\kappa}^2$ there are $0<c<C$ with
$c\,\mathcal E(\boldsymbol s)^2\le D(\boldsymbol s)-n\pi/\kappa\le C\,\mathcal E(\boldsymbol s)^2$ near the base.
\end{enumerate}
\end{corollary}

\begin{proof}
(i) After tilting away the real mean the path has mean imaginary part $\kappa_0$, so
\cref{fil:thm:mirrorhessian}\textup{(ii)} gives the expansion; $u_0=0$ by anchoring and $u_n=0$ by
\cref{fil:thm:differential}. Let $U(z)=\sum_kw_ku_kz^k$; the first variation of the transform is
$\delta F(s)=A_1(e^{hs})+sU(e^{hs})$ with $A_1$ real, and $F_0(s)=P(e^{hs})/P(1)$. Differentiating the zero of order
$m_j$ at the moving target, $\delta F+v_jF_0'$ vanishes to order $m_j$ at $s_j^0$. Conjugating the same statement
for $j^*$ and shifting by $2i\kappa_0$ (note $e^{2i\kappa_0h}=1$), subtraction shows that
$2\kappa_0U(z)+(hd_j/P(1))zP'(z)$ is divisible by $(z-z_j)^{m_j}$ for every $j$. The displayed right side satisfies
these congruences: modulo $(z-z_j)^{m_j}$ only its $j$th term survives, and $zP'(z)\equiv m_jz_jP(z)/(z-z_j)$. It is
real by the pairing ($d_{j^*}=\overline{d_j}$, $z_{j^*}=\overline{z_j}$), has degree below $n$, and its constant term is a
multiple of $\sum_jm_jd_j=2\sum_jm_j\Im v_j=0$. As $U$ also has degree below $n$, the difference is divisible by $P$
and vanishes. Finally $Q(v)=0$ exactly when $U=0$, that is $\sum_jm_jd_jz_j/(z-z_j)=0$, which for distinct nonzero
poles means $d_j=0$ for all $j$.

(ii) The rescaled path $\boldsymbol s^0+\eps v$ divided by $1+\eps a/\kappa_0$ has mean imaginary part $\kappa_0$ and equals
$\boldsymbol s^0+\eps\hat v+O(\eps^2)$. Positive scaling divides the width by the scale, and the term of order $\eps^2$ in
the rescaled path does not affect the expansion, because the gradient of $D$ on the fixed-mean hyperplane vanishes
at the base. Expanding gives the formula, and a velocity with $v_j=\overline{v_{j^*}}$ has $a=0$.

(iii) The projection $s_j^{\mathrm{mir}}=(s_j+\overline{s_{j^*}}+2i\kappa)/2$ preserves the weighted mean and lands on the
reflection locus, inside the positive region near the base; $\delta=\boldsymbol s-\boldsymbol s^{\mathrm{mir}}$ has weighted
norm $\mathcal E/2$ and satisfies $\delta_{j^*}=-\overline{\delta_j}$, a direction on which the quadratic form of \textup{(i)} at
$\boldsymbol s^{\mathrm{mir}}$ is positive definite, uniformly on a compact neighbourhood. Taylor's formula about
$\boldsymbol s^{\mathrm{mir}}$, with a uniform cubic remainder, gives both bounds.
\end{proof}

The quadratic form has elementary explicit bounds. With $w_{\min}=\min_kw_k$ and $n\ge2$,
\[
 \frac{2e^{-2n}w_0w_{\min}}h\sum_{k=1}^{n-1}u_k^2\le Q
 \le\frac1h\sum_{k=1}^{n-1}w_k\Bigl(\frac1{w_0k}+\frac1{w_n(n-k)}\Bigr)u_k^2 :
\]
for the lower bound, $\varpi_x(j)\ge e^{-n}w_j$ on $\abs{hx}\le1$, and one keeps the pairs $(0,k)$ in the variance; for
the upper bound, $\operatorname{Var}(u)\le\sum_k\varpi_x(k)u_k^2$, with $\varpi_x(k)\le(w_k/w_n)e^{-(n-k)hx}$ for $x\ge0$ and
$\varpi_x(k)\le(w_k/w_0)e^{khx}$ for $x\le0$. Together with the explicit linear maps of
\cref{fil:thm:mirrorhessian,fil:cor:mirrorlabel} they give computable coercivity constants.

\begin{example}[Explicit second variations]\label{fil:ex:mirror}
(i) For $\boldsymbol s^0=(i,2i,3i)$, all simple, $\kappa_0=2$, $h=\pi/2$, $P=1+z+z^2+z^3$ and all $w_k=\tfrac14$. Write
$v_j=\alpha_j+i\beta_j$ with $\beta_1+\beta_2+\beta_3=0$, and $X=\beta_1+\beta_3=-\beta_2$, $Y=\alpha_1-\alpha_3$. Then
$u_1=\pi(X-Y)/4$, $u_2=-\pi(X+Y)/4$, and with $\int_0^\infty P^{-2}\dd y=\pi/8$,
$\int_0^\infty yP^{-2}\dd y=\pi/8-\tfrac14$, $\int_0^\infty y^2P^{-2}\dd y=\tfrac12-\pi/8$ (the others follow by $y\mapsto1/y$),
\[
 D(\boldsymbol s^0+\eps v)=\frac{3\pi}2+\frac\pi{16}\bigl(3X^2+(\pi-1)Y^2\bigr)\eps^2+O(\eps^3).
\]
Asymmetry of the frequencies and unequal outer real parts cost separately; common changes of the real parts,
changes of the middle real part and opposite shifts of the outer frequencies stay on the mirror locus.

(ii) For real frequencies $(1-d,1,1+d)$, $0<d<\tfrac23$, $\kappa=1$, and the direction $(1,-2,1)$, put $c=\cos\pi d>-\tfrac12$.
Then $P(z)=(z+1)(z^2+2cz+1)$, $u_1=-u_2=2\pi(1-c)/(1+2c)$, and
\[
 D(1-d+\eps,1-2\eps,1+d+\eps)=3\pi+\eps^2Q_d+O(\eps^3),\qquad
 Q_d=6\pi\Bigl(1-\frac{3\pi d\cos\pi d}{(1+2\cos\pi d)\sin\pi d}\Bigr).
\]
Partial fractions reduce the variance integral to $\int_0^\infty\dd z/(z^2+2cz+1)=\arccos c/\sqrt{1-c^2}$ and
$\int_0^\infty z\,\dd z/(z^2+2cz+1)^2=(1-c\arccos c/\sqrt{1-c^2})/(2(1-c^2))$. Thus $Q_{1/2}=6\pi$,
$Q_d=\tfrac{3\pi^5}{10}d^4+\tfrac{3\pi^7}{28}d^6+O(d^8)$ as $d\downarrow0$, in agreement with \cref{fil:cor:sixth} (here
$e_3=2d^2\eps+O(\eps^3)$), and $Q_d\sim4\pi/(\tfrac23-d)$ as $d\uparrow\tfrac23$, where the two inner masses tend to $0$. The
neighbourhoods of \cref{fil:thm:perturb} therefore cannot be chosen with control uniform up to a vanishing weight.

(iii) Two targets $s_\pm=\pm\beta+i(\kappa\pm\delta)$ collide at $\beta=\delta=0$. Their centred polynomial is
$y^2-(\beta+i\delta)^2$, so $\Im c_2=-2\beta\delta$; at the base $P=(1+z)^2$, for $V=1$ one finds $B_V(z)=-h^2z/(4\kappa)$,
$\sigma_1=-h^2/(2\kappa)$ and $Q(1)=4\sigma_1^2/(3h)=\pi^3/(3\kappa^5)$. Hence, for $\abs\beta+\abs\delta$ small relative to $\kappa$,
\[
 D=\frac{2\pi}\kappa+\frac{4\pi^3}{3\kappa^5}\beta^2\delta^2\Bigl[1+O\Bigl(\frac{\beta^2+\delta^2}{\kappa^2}\Bigr)\Bigr],
\]
and $D=2\pi/\kappa$ exactly when $\beta\delta=0$: pure damping and pure splitting cost nothing, their interaction a
positive quartic.

(iv) Mixed entries have no fixed sign. Let $\kappa=1$, $h=\pi$, $a=\log2/\pi$ and base roots $y=-2a,-a,3a$ (frequency
one, different real parts), so that $P_{c^0}=y^3-7a^2y-6a^3$, $P(z)=(z+\tfrac14)(z+\tfrac12)(z+8)$, and the base weights are
$(8,49,70,8)/135$. For $V=v_2y+v_3$ write $Q(V)=(v_2,v_3)\mathsf Q(v_2,v_3)^{\mathsf T}$. The interpolation formula of
\cref{fil:thm:mirrorhessian} gives $(\sigma_1,\sigma_2)=(-\tfrac{173\pi}{980a},-\tfrac{7\pi}{50a})$ for $V=y$ and
$(\tfrac\pi{20a^2},-\tfrac{13\pi}{350a^2})$ for $V=1$, and exact rational integration gives
\begin{gather*}
 \mathsf Q_{22}=\frac\pi{a^2}\Bigl(\frac{49}{200}-\frac{6646079\log2}{22785000}\Bigr),\\
 \mathsf Q_{33}=\frac\pi{a^4}\Bigl(\frac{21}{200}-\frac{21933\log2}{155000}\Bigr),\\
 \mathsf Q_{23}=\frac\pi{a^3}\Bigl(-\frac{27}{200}+\frac{29441\log2}{155000}\Bigr),
\end{gather*}
approximately $2.7633$, $9.1705$ and $-0.9776$; $\mathsf Q_{23}<0$ since $\log2<\tfrac7{10}$. Time reversal of a positive
measure sends targets $s$ to $-\bar s$ and preserves the width; it replaces $P_c(y)$ by $(-1)^n\overline{P_c(-\bar y)}$
and $(v_2,v_3)$ by $(-v_2,v_3)$. So the base with roots $2a,a,-3a$ has the opposite mixed entry, and along the bases
with roots $ty$, $-1\le t\le1$, which are all strictly positive mirror configurations (the roots of $P$ are
negative), the mixed entry passes through zero. The strict signs of all minors in \cref{fil:thm:hessian} concern
real frequencies at a coalesced base.
\end{example}

The formulas of (i), (ii) and (iv) are checked in exact arithmetic and against high-precision solutions of
the moment equations in
\begin{center}\small
\artifact{ProveIt_research_2026-11-27/04_mirror_hessian/mirror_hessian_checks.json}\\
\artifact{deep_padic_roots_2026-11-27/data/filter_hessian_results.json}\\
\artifact{proveit_three_extensions_2026-11-27_1/papers/positive_filters/mirror_checks.json}
\end{center}
for example, in (iv) the ratios of the computed excess to $\eps^2Q$ in the directions $(1,0)$, $(0,1)$, $(1,1)$ are
$1-O(\eps)$ at $\eps=10^{-3},10^{-4},10^{-5}$. The proofs do not depend on these computations.

\Cref{fil:thm:perturb,fil:thm:mirrorhessian} apply in particular near a single frequency of high order, where
the prescribed frequencies split.

\begin{theorem}[Splitting a multiple frequency; \stU]\label{fil:thm:split}
For every $N\ge1$ there is $\eps_N>0$ such that, if $\lambda_1,\dots,\lambda_N>0$ satisfy
$\max_j\abs{\lambda_j/\kappa-1}<\eps_N$ for some $\kappa>0$, the least-diameter probability measure whose
Fourier transform vanishes at the multiset $\set{\lambda_j}$ (coincident entries giving the combined
order) is unique up to translation. With left endpoint $0$ it has exactly $N+1$ positive atoms
$0=t_0<\dots<t_N$; nodes and weights are real analytic in the entries, including at coincidences,
symmetric in them, equal to $t_j=j\pi/\kappa$, $w_j=2^{-N}\binom Nj$ on the diagonal, and the minimizer is
symmetric about $t_N/2$. Its width $B(\set{\lambda_j})=t_N$ is at least $N\pi/\bar\lambda$,
$\bar\lambda=N^{-1}\sum_j\lambda_j$, by \cref{fil:thm:meanbound}, and for $N\ge3$ it is determined through
sixth order by \cref{fil:cor:sixth}.
\end{theorem}

\begin{proof}
Scale to $\kappa=1$. For small $\eps_N$ the targets $i\lambda_j$ have mean $i\bar\lambda$ close to $i$, and after division by
$\bar\lambda$ their centred polynomial $\prod_j(y-i(\lambda_j/\bar\lambda-1))$ is close to $y^N$. This is the base of
\cref{fil:thm:reflect} with $h=\pi$, all $N$ targets at $i$ and $P(z)=(1+z)^N$, whose coefficients are all positive. The
coefficient form of \cref{fil:thm:perturb} gives uniqueness, the $N+1$ positive atoms, and nodes and weights real
analytic in the coefficients; these are symmetric polynomials in the entries, so nodes and weights are symmetric real
analytic functions of the entries, also at coincidences. On the diagonal the minimizer is the binomial filter of
\cref{rec:thm:fwidth}. For real frequencies the time reversal $t\mapsto-t$ preserves the vanishing conditions and the
diameter, so uniqueness makes the minimizer symmetric about $t_N/2$.
\end{proof}

\paragraph{The cost of asymmetric splitting.}
Fix $N\ge3$. By scaling, the normalized width $b=\kappa B$ depends only on the relative deviations
$u_j=\lambda_j/\kappa-1$ from the mean $\kappa=N^{-1}\sum_j\lambda_j$; put $d=\max_j\abs{u_j}$ and
\[
 P(x)=\prod_j(x-u_j)=x^N+\sum_{r=2}^Nc_rx^{N-r},\qquad c_r=(-1)^re_r(u),
\]
with $e_r$ the elementary symmetric functions; the coefficient $c_1=-e_1$ vanishes. Write $E$ for the vector of
even-indexed and $O$ for the vector of odd-indexed coefficients, and $\mathcal I_N=\set{3,5,\dots}\cap[3,N]$ for the
odd indices. Working with coefficients avoids differentiating individual roots at a multiple root.

\begin{lemma}[Coefficient continuation; \stU]\label{fil:lem:coeff}
Fix a small circle $C$ about $0$. For real $c=(c_2,\dots,c_N)$ near $0$ and entire $f$ put
\[
 R_cf(x)=\frac1{2\pi i}\oint_C\frac{f(1+z)}{P_c(z)}\,\frac{P_c(z)-P_c(x)}{z-x}\dd z .
\]
It is a polynomial of degree below $N$, analytic in $c$, it agrees with $f(1+\cdot)$ at the roots of $P_c$ to
their multiplicities, and $R_0f$ is the Taylor polynomial of $f$ at $1$ of degree $N-1$. There are a
neighbourhood $U$ of $0$ and real analytic $w_j(c)>0$, $0=t_0(c)<\dots<t_N(c)$ on $U$ with $\sum_jw_j=1$ and
$R_c\widehat{\mu_c}=0$ for $\mu_c=\sum_jw_j\delta_{t_j}$, and $t_j(0)=j\pi$, $w_j(0)=2^{-N}\binom Nj$. When the roots of
$P_c$ are real, $\mu_c$ is the minimizer of \cref{fil:thm:split} for the frequencies $1+u_j$. Thus $b(c)=t_N(c)$
continues the normalized width real analytically to all of $U$, also where the roots of $P_c$ are not all real.
\end{lemma}

\begin{proof}
The interpolation property is that of $\mathcal R_c$ in \cref{fil:thm:perturb}, in the variable $1+x$. With $\kappa=1$ the
targets $i(1+u_j)$ have centred polynomial $\prod_j(y-iu_j)=y^N+\sum_ri^rc_ry^{N-r}$, so real $c$ near $0$ form a slice
of the coefficient coordinates of \cref{fil:thm:perturb} at the base $y^N$ of \cref{fil:thm:split}, and the vanishing
of $R_c\widehat\mu$ is the vanishing of $\mathcal R$ there. The coefficient form of \cref{fil:thm:perturb} restricted to
this slice gives the branch, with the binomial filter at $c=0$; $\mu_c$ is the minimizer for every $c\in U$, in
particular that of \cref{fil:thm:split} at real roots.
\end{proof}

\begin{lemma}[Defect integral; \stU]\label{fil:lem:defect}
After shrinking $U$, for $c\in U$ put $G_c(s)=\sum_jw_j(c)e^{t_j(c)s}$ and $p_j(x,c)=w_j(c)e^{xt_j(c)}/G_c(x)$. Then
\begin{equation}\label{fil:eq:defect}
 b(c)-N\pi=\frac12\int_\R-\log\Bigl\lvert\sum_jp_j(x,c)e^{2it_j(c)}\Bigr\rvert\dd x\ \ge0 .
\end{equation}
Moreover $b(E,0)=N\pi$, $\partial_Ob(E,0)=0$, and $b-N\pi=\sum_{a,a'\in\mathcal I_N}c_ac_{a'}G_{aa'}(c)$ with $G_{aa'}$ analytic
near $0$: every monomial of the Taylor series of $b-N\pi$ contains at least two odd coefficients.
\end{lemma}

\begin{proof}
This is the coalesced case of \cref{fil:lem:mirrordefect,fil:thm:mirrorhessian}, with $\kappa=1$, $n=N$ and base $y^N$, in
the slice of real $c$ of \cref{fil:lem:coeff}, where the coefficients of the centred polynomial are $i^rc_r$: the even
$c_r$ give real and the odd ones imaginary coefficients. Thus \eqref{fil:eq:defect} is \eqref{fil:eq:mirrordefect}; the
integrand is nonnegative since $\sum_jp_j=1$. The locus $O=0$ is the reflection locus, where $b=N\pi$, and the
factorization $b-N\pi=O^{\mathsf T}M(E,O)O$ in the proof of \cref{fil:thm:mirrorhessian}\textup{(iii)}, restricted to the
slice, gives $\partial_Ob(E,0)=0$ and the last assertion.
\end{proof}

\begin{theorem}[Width Hessian at a coalesced frequency; \stU]\label{fil:thm:hessian}
There is a symmetric matrix $H_N$ indexed by $\mathcal I_N$ such that near $c=0$
\[
 b(E,O)=N\pi+O^{\mathsf T}H_NO+O_N\bigl((\norm E+\norm O)\norm O^2\bigr);
\]
thus $2H_N$ is the Hessian of $b$ in the odd coordinates at $0$. For all sets $I,J\subseteq\mathcal I_N$ of equal positive
size, listed increasingly,
\[
 \det(H_N)_{I,J}>0 .
\]
Explicitly, define the Bernoulli cumulants $K_r(p)$ by $\log(1-p+pe^t)=\sum_{r\ge1}K_r(p)t^r/r!$, let $\mathcal P_j$ be the
Legendre polynomial with $\mathcal P_j(1)=1$, and for $r\in\mathcal I_N$ put
\[
 M_r(p)=\frac{(-1)^{(r-1)/2}}{(r-1)!}K_r(p),\qquad M_r'(p)=\sum_{n\in\mathcal I_N}a_{nr}\mathcal P_{n-1}(2p-1).
\]
Then $a_{nr}=0$ for $n>r$, and for $n\le r$, with $\filStir rk$ the Stirling numbers of the second kind,
\[
 a_{nr}=\frac{(2n-1)(-1)^{(r-1)/2}}{(r-1)!}\sum_{k=n}^r(-1)^{k-1}k!\filStir rk\frac{((k-1)!)^2}{(k-n)!\,(k+n-1)!}.
\]
Finally
\begin{equation}\label{fil:eq:hessfactor}
 (H_N)_{rt}=\pi^{r+t-1}\sum_{n\in\mathcal I_N}\beta_{N,n}a_{nr}a_{nt},\qquad
 \beta_{N,n}=\frac{(N+n-1)!\,(N-n)!/((N-1)!\,N!)-1}{(2n-1)n(n-1)}>0 .
\end{equation}
\end{theorem}

\begin{proof}
\emph{The quadratic coefficient.} Let $c=\eps(0,v)$ with $v$ supported on the odd indices and
$s_j=\partial_\eps t_j\vert_{\eps=0}$. By \cref{fil:thm:mirrorhessian}\textup{(ii)} the coefficient of $\eps^2$ in $b$ is
$\int_\R\operatorname{Var}_{p(x,0)}(s)\dd x$. At the base the weights $p_j(x,0)$ are binomial with
$p=e^{\pi x}/(1+e^{\pi x})$, so this coefficient is
\begin{equation}\label{fil:eq:variance}
 Q(v)=\frac1\pi\int_0^1\frac{\operatorname{Var}(s_J)}{p(1-p)}\dd p,\qquad J\sim\mathrm{Bin}(N,p).
\end{equation}

\emph{The velocities are Bernoulli cumulants.} Take a unit variation of $c_r$, $r$ odd, put
$C_0=(-1)^N2^{-N}$, $z=e^{i\pi x}$, $\ell=\log z$, and $T_r(z)=\sum_j(-1)^jw_j(0)s_jz^j$. At the base
$\widehat{\mu_0}(1+x)=C_0(z-1)^N$, and a variation of the transform at $1+x$ has the form $A(z)+(i+\ell/\pi)T(z)$
with real polynomials $A$, $T$ of degree at most $N$, $T(0)=0$ from anchoring. Differentiating the
divisibility of $\widehat{\mu_c}(1+x)$ by $x^N+c_rx^{N-r}$, the first variation is $\equiv C_0(z-1)^N/x^r$ modulo $x^N$, and
$x=\ell/(i\pi)$. Taking imaginary parts as formal series in real $z-1$,
\[
 T_r(z)\equiv(-1)^{(r-1)/2}\pi^rC_0\frac{(z-1)^N}{(\log z)^r}\pmod{(z-1)^N},\qquad T_r(0)=0,
\]
which determines a polynomial of degree at most $N$ uniquely. Define the Eulerian numerator $D_{r-1}$ by
$\sum_{m\ge1}m^{r-1}z^m=D_{r-1}(z)/(1-z)^r$ for $\abs z<1$; it vanishes at $0$ and has degree at most $r-1$. Applying
$(\dd/\dd y)^{r-1}$ to $e^y/(1-e^y)=-1/y+(\text{holomorphic})$ gives, for odd $r$,
$D_{r-1}(z)/(r-1)!=(z-1)^r/(\log z)^r+O((z-1)^r)$, so
$T_r(z)=(-1)^{(r-1)/2}\pi^rC_0(z-1)^{N-r}D_{r-1}(z)/(r-1)!$. The binomial mean of the velocities is
$2^N(1-p)^NT_r(-p/(1-p))$, and with $z=-p/(1-p)$ the generating function of the cumulants gives
$K_r(p)=-D_{r-1}(z)/(1-z)^r$ (first for $p<\tfrac12$, expanding $\log(1-p+pe^t)$ in powers of $ze^t$, then by polynomial
identity). Hence
\begin{equation}\label{fil:eq:means}
 \mathbb E[s_J]=\pi^rM_r(p).
\end{equation}
One family of polynomials thus serves every cluster size $N\ge r$. Since $M_r$ vanishes at $0,1$ and
$M_r(1-p)=-M_r(p)$, the derivative $M_r'$ is symmetric about $\tfrac12$ with zero integral, so its Legendre
indices are among the even degrees $2,4,\dots,r-1$, i.e.\ $n\in\mathcal I_N$, $n\le r$.

\emph{Orthogonal decomposition of the variance.} If $f$ is real on $\set{0,\dots,N}$ and $m(p)=\mathbb Ef(J)$, then,
with $(N)_s=N!/(N-s)!$,
\begin{equation}\label{fil:eq:hoeffding}
 \operatorname{Var}f(J)=\sum_{s=1}^N\frac{[p(1-p)]^s}{s!\,(N)_s}\bigl(m^{(s)}(p)\bigr)^2 .
\end{equation}
Indeed write $J=\sum_iX_i$ with independent Bernoulli variables of mean $p$; the multilinear extension of
$f(\sum_iX_i)$, expanded at $(p,\dots,p)$, has the common coefficient $m^{(s)}(p)/(N)_s$ at each product
$\prod_{i\in S}(X_i-p)$ with $\abs S=s$, and these products are orthogonal. If $m'(p)=\sum_nh_n\mathcal P_{n-1}(2p-1)$, the
$s$th term of \eqref{fil:eq:hoeffding}, after integration in \eqref{fil:eq:variance}, is diagonal in the $h_n$: for
$1\le s\le n$,
\[
 \int_0^1[p(1-p)]^{s-1}\Bigl(\frac{\dd^{s-1}}{\dd p^{s-1}}\mathcal P_{n-1}(2p-1)\Bigr)^2\dd p=\frac{(n+s-2)!}{(2n-1)(n-s)!},
\]
and mixed degrees are orthogonal, by $s-1$ integrations by parts in the Legendre equation (equivalently the
associated Legendre norms after $x=2p-1$). Hence the coefficient of $h_n^2/\pi$ is
$\frac1{2n-1}\sum_{s=1}^n\frac{(n+s-2)!}{s!\,(n-s)!\,(N)_s}$. Multiplied by $n(n-1)$ the summands become
$\binom ns\binom{n+s-2}s/\binom Ns$; including $s=0$, the terminating Chu--Vandermonde identity evaluates the sum as
$(-N-n+1)^{\overline n}/(-N)^{\overline n}$ (rising factorials; $n\le N$), the factorial ratio in \eqref{fil:eq:hessfactor},
and subtracting the term $s=0$ gives $\beta_{N,n}$, which is positive since all terms are. With \eqref{fil:eq:means},
$m'=\sum_rv_r\pi^rM_r'$, and polarization gives \eqref{fil:eq:hessfactor}. For $a_{nr}$ expand
$K_r(p)=\sum_{k=1}^r(-1)^{k-1}(k-1)!\filStir rkp^k$ and use Rodrigues' formula,
$\int_0^1p^{k-1}\mathcal P_{n-1}(2p-1)\dd p=((k-1)!)^2/((k-n)!(k+n-1)!)$ for $k\ge n$ and $0$ for $k<n$. Jacobi
diagonalizations of related Bernstein--Durrmeyer operators are classical \cite{Waldron2003Generalised}; here the
cumulant identity \eqref{fil:eq:means} decides which perturbations of the extremal atoms enter which modes.

\emph{Strict positivity of every minor.} The cumulants satisfy $K_1(p)=p$ and $K_{r+1}=p(1-p)K_r'$ (differentiate
$\log(1+e^\theta)$ in $\theta=\log(p/(1-p))$), so $K_r'$ arises from $\mathcal P_0$ by $r-1$ applications of
$f\mapsto(p(1-p)f)'$. With $x=2p-1$ this operator sends $\mathcal P_l$ to
\begin{equation}\label{fil:eq:walk}
 \frac{l(l-1)}{2(2l+1)}\mathcal P_{l-1}-\frac{(l+1)(l+2)}{2(2l+1)}\mathcal P_{l+1}
\end{equation}
(omit the first term for $l=0$), by the Legendre derivative identity and the three-term recurrence. On the
lattice with time increasing by one and height $l\ge0$, give the up and down steps the absolute values of the
coefficients in \eqref{fil:eq:walk} and delete steps of weight zero; positive down steps start at height at least
two. The coefficient of $\mathcal P_{n-1}$ in $K_r'$ is the signed sum of the paths from $(1-r,0)$ to $(0,n-1)$, each
with exactly $(r+n-2)/2$ up steps, so the sign of $a_{nr}$ is $(-1)^{(n-1)/2}$ for $n\le r$ and $(r-1)!\abs{a_{nr}}$ is the
unsigned path weight. For increasing lists $r_1<\dots<r_k$ and $n_1<\dots<n_k$ in $\mathcal I_N$, the determinant of
the unsigned path matrix is the sum of the weights of vertex-disjoint families joining corresponding indices,
by the tail-swapping involution at a first intersection \cite{GesselViennot1985Binomial}. Only the order-preserving
matching contributes: a path starting earlier from height $0$ has positive height when a later one starts, since
no step down to height $0$ has positive weight, so if it ended below the later path the two would meet; heights
at common times have equal parity, so paths cannot cross in the middle of an edge. If $n_i\le r_i$ for all $i$, a
disjoint family exists: let path $i$ take all its up steps first, so that its height at time $\tau\in[1-r_i,0]$
is $\min(\tau+r_i-1,n_i-1-\tau)$; at common times both arguments for the higher index exceed those for the lower
one by at least two. If some $n_i>r_i$ the order-preserving connection is impossible by length. Hence
$\det(\abs{a_{n_ir_j}})\ge0$, with strict inequality if and only if $n_i\le r_i$ for all $i$. The signs of the
$a_{nr}$ depend only on $n$, so after absorbing the powers of $\pi$ in positive diagonal scalings the matrix
\eqref{fil:eq:hessfactor} is $\abs a^{\mathsf T}\operatorname{diag}(\beta_{N,n})\abs a$. The Cauchy--Binet formula expands
each minor in nonnegative terms, and the term with $n_i=2i+1$ ($1\le i\le k$) is strictly positive, since the
$i$th entry of an increasing list in $\mathcal I_N$ is at least $2i+1$.

\emph{Remainder.} By \cref{fil:lem:defect}, $b(E,O)-N\pi=O^{\mathsf T}H(E)O+O(\norm O^3)$ with $2H(E)=\partial^2_{OO}b(E,0)$
and $H(E)=H_N+O(\norm E)$, where $O^{\mathsf T}H_NO=Q(O)$ by the first step.
\end{proof}

Since there are finitely many minors and $b$ is real analytic, the same strict signs hold for the Hessian
$\partial^2_{OO}b(E,O)$ on a neighbourhood of $0$. They concern the interaction of perturbations and go beyond
positive definiteness. Every entry of $H_N^{-1}$ has the strict sign $(-1)^{i+j}$, $i,j$ the positions in the
increasing list $\mathcal I_N$, by the cofactor formula. For distinct odd coordinate directions $\mathbf e_i,\mathbf e_j$
and small $h,k>0$, integrating the positive mixed derivative over the rectangle gives
\[
 b(E,O+h\mathbf e_i+k\mathbf e_j)-b(E,O+h\mathbf e_i)-b(E,O+k\mathbf e_j)+b(E,O)>0;
\]
at corners whose polynomials are real-rooted each term is a genuine optimal width.

\begin{corollary}[Two-sided bound and the sixth-order law; \stU]\label{fil:cor:sixth}
Let $N\ge3$ and $S=\sum_{r\in\mathcal I_N}e_r^2$. There are $0<c_N<C_N$ such that for all sufficiently small $d$
\[
 c_NS\le\kappa B(\boldsymbol\lambda)-N\pi\le C_NS,
\]
and, uniformly,
\[
 \kappa B(\boldsymbol\lambda)=N\pi+\frac{N\pi^5}{20(N-1)(N-2)}\,e_3^2+O_N(d^8).
\]
Equivalently, with $H=\max_j\abs{\lambda_j-\kappa}$,
\[
 B(\boldsymbol\lambda)=\frac{N\pi}\kappa+\frac{N\pi^5}{180(N-1)(N-2)}\,\frac{\bigl(\sum_j(\lambda_j-\kappa)^3\bigr)^2}{\kappa^7}
 +O_N\Bigl(\frac{H^8}{\kappa^9}\Bigr).
\]
The first five orders of detuning beyond the mean leave the optimum unchanged; there is no term of degree
seven, and the sixth-order term depends only on the squared third central moment.
\end{corollary}

\begin{proof}
All principal minors of $H_N$ are positive, so $H_N$ is positive definite, and the remainder in
\cref{fil:thm:hessian} is $o(\norm O^2)$; since $c_r^2=e_r^2$, this gives the two-sided bound. With $c_r=O_N(d^r)$ the
remainder is $O_N(d^2d^6)$, and every term of $O^{\mathsf T}H_NO$ except $(H_N)_{33}c_3^2$ is $O_N(d^8)$. From
\eqref{fil:eq:hessfactor}, $\beta_{N,3}=N/(5(N-1)(N-2))$ and $a_{33}=-\tfrac12$, so $(H_N)_{33}=N\pi^5/(20(N-1)(N-2))$. Newton's
identity with $e_1=0$ gives $\sum_ju_j^3=3e_3$; rescale by $\kappa$.
\end{proof}

For $r=3$ the velocities are explicit: $s_j=\pi^3j(N-j)(2j-N)/(2N(N-1)(N-2))$ for a unit change of $c_3$, with binomial
mean $\pi^3p(1-p)(2p-1)/2$, and \eqref{fil:eq:hoeffding} has three nonzero terms, whose integrals give
$\frac1{20N}+\frac3{20N(N-1)}+\frac1{5N(N-1)(N-2)}=\frac N{20(N-1)(N-2)}$. A direct second-variation computation along the
coefficient path $y^N-ty^{N-3}$ gives the same constant. For $N=3,4$ the only odd coordinate is $e_3$, so near a
coalesced frequency the excess is comparable to $e_3^2$, and it vanishes exactly for clusters that are mirror
symmetric about their mean (one frequency at the mean and two equidistant from it, or two reflected pairs), as
\cref{fil:thm:meanbound} requires. Cancelling the frequencies separately by two-point filters costs
$\pi\sum_j1/\lambda_j=N\pi/\kappa+(\pi/\kappa)\sum_ju_j^2+O_N(d^3/\kappa)$: the optimum has no quadratic penalty, and its
first possible penalty, of order six, depends only on the asymmetric part of the cluster.

\begin{corollary}[Every asymmetric ray; \stU]\label{fil:cor:ray}
Let $v$ be a fixed centred real vector whose multiset is not invariant under $v\mapsto-v$, and let $r$ be the least
odd index with $e_r(v)\ne0$. For $u_j=hv_j$ and $h\to0$,
\[
 \kappa B-N\pi=(H_N)_{rr}\,e_r(v)^2h^{2r}+O_{N,v}(h^{2r+2}).
\]
\end{corollary}

\begin{proof}
A centred real multiset has all odd coefficients zero precisely when it is invariant under reflection, so $r$
exists. Here $E=O_{N,v}(h^2)$, the first nonzero odd coefficient is $-e_r(v)h^r$, and the next possible one has degree
$r+2$; apply \cref{fil:thm:hessian}.
\end{proof}

\begin{example}[Unequal orders and exceptional rays]\label{fil:ex:exceptional}
(i) For small real $\eps\ne0$, order two at $1-\eps$ and order one at $1+2\eps$ give $\sum_ju_j^3=6\eps^3$ and
\[
 B(1-\eps,1-\eps,1+2\eps)=3\pi+\tfrac{3\pi^5}{10}\eps^6+O(\eps^8)>3\pi .
\]
For $\eps=0.01$ a high-precision solution of the three moment equations gives the excess
$9.1902306864\cdot10^{-11}$ against the leading term $9.1805905436\cdot10^{-11}$; for $(1-2t,1-t,1+3t)$, where $e_3=6t^3$, the
quotient $(B-3\pi)/(36t^6)$ divided by its limit $3\pi^5/40$ is $1.0612$ at $t=1/20$ and $1.00024$ at $t=1/320$. These are
floating-point illustrations; radii of global optimality are not certified.

(ii) For $N=5$ and $v=(-9,-5,-1,7,8)$ the entries and their cubes sum to zero, so $e_3(v)=0$, and $e_5(v)=-2520$. For a
unit change of $c_5$ one finds $T/(\pi^5C_0)=z(z+1)(z^2+10z+1)/24$, velocities $\pi^5(0,\tfrac1{120},-\tfrac{11}{240},\tfrac{11}{240},-\tfrac1{120},0)$
with binomial mean $\pi^5p(p-1)(2p-1)(12p^2-12p+1)/24$, and $(H_5)_{55}=25\pi^9/6912$, so the first correction has degree ten:
$\kappa B-5\pi\sim\frac{25\pi^9}{6912}\,2520^2\,t^{10}$ along $u=tv$.

(iii) For $N=7$ and $D=\operatorname{diag}(\pi^3,\pi^5,\pi^7)$ in the order $3,5,7$,
\[
 \pi D^{-1}H_7D^{-1}=
 \begin{pmatrix}
 7/600&1/1440&7/216000\\
 1/1440&11/17280&371/5702400\\
 7/216000&371/5702400&189091/1010880000
 \end{pmatrix}.
\]
The roots $v_1,\dots,v_7$ of $x(x^2-1)(x^2-4)(x^2-9)+1$ are real and simple: at either end of each interval
$[j-\tfrac14,j+\tfrac14]$, $j=-3,\dots,3$, the unshifted product has absolute value at least
$\tfrac14(\tfrac34)^2(\tfrac74)^2(\tfrac{11}4)^2>1$ (each positive integer distance from $j$ occurs at most twice among the other
roots) and changes sign across the interval, and adding one keeps the seven sign changes. The centred polynomial
has $c_3=c_5=0$, $c_7=1$, so the frequencies $\kappa(1+hv_j)$ satisfy
\[
 \kappa B-7\pi=\frac{189091\,\pi^{13}}{1010880000}\,h^{14}+O(h^{16}),
\]
a family whose leading width response occurs at order fourteen. The minor of $H_7$ on rows $3,5$ and columns $5,7$
compares two mixed responses with a third, which is invisible if one coefficient varies at a time.
\end{example}

The entries for $3\le N\le9$, all their minors, the coefficient of \cref{fil:cor:sixth} for symbolic $N$, the
tenth-order example and the numerical illustrations are checked in exact or high-precision arithmetic in
\begin{center}\small
\artifact{gram-euler-and-zero-filters/filters/verification.txt}\\
\artifact{grh_cayley_and_filters/filters/filter_checks.txt}\\
\artifact{grh_hidden_forms_and_filters_2026-11-27/data/spectral_skewness_verification.json}
\end{center}
The proofs above cover every $N$.

\begin{remark}[Actual clusters]\label{fil:rem:clusters}
For each fixed $N$ and relative radius there are arbitrarily high clusters of $N$ distinct zeta
ordinates within that radius: otherwise every multiplicative interval of that radius would contain
fewer than $N$ distinct ordinates, there would be $O_N(\log T)$ distinct ordinates below $T$, and the unit
interval bound $O(\log T)$ on multiplicities would give $O_N(\log^2T)$ zeros, against the
Riemann--von Mangoldt formula. Under \cref{src:thm:dgrh} \cref{fil:thm:split} therefore applies to genuine
zero clusters, without any bound on gaps. Numerical solutions for the nonsymmetric clusters
$(0.98,1,1.05)$ and $(0.98,0.98,1.01,1.06)$ have widths agreeing with $N\pi/\bar\lambda$ to $10^{-8}$
(\artifact{optimal_dirichlet_zero_filters/checks.json}); the neighbourhood $\eps_N$ is not certified.
In absolute units, by the Riemann--von Mangoldt formula $N_\zeta(t)=\frac t{2\pi}\log\frac t{2\pi e}+O(\log t)$ there is
a constant $C_N$ such that every window $[T,T+C_N]$ with $T$ large contains at least $N$ ordinates counted with
multiplicity, and by \cref{fil:cor:sixth} any $N$ of them (repeated ordinates prescribing repeated vanishing)
have a filter of optimal width $N\pi/\kappa+O_N(C_N^6/\kappa^7)$ as $T\to\infty$. For a zero
$\rho=\beta+i\gamma$ in an explicit formula the tilt by $e^{\beta t}$ turns the multiplier $\int e^{\rho t}\dd\mu$ into the
Fourier condition at $\gamma$ with the same support, so on the critical line, under \cref{src:thm:dgrh} for Dirichlet
$L$-functions, these widths are the exact least costs of cancelling tight clusters of zero terms by positive
averages.
\end{remark}

The optimizers are stable, and their intervals cannot be shortened cheaply. Use any norm on the moment
space, its dual $\norm\cdot_*$, and $\eps(\nu)=\norm{\int Y\dd\nu}$.

\begin{theorem}[Sharp stability; \stU]\label{fil:thm:stable}\label{rec:prop:fstable}
Let $\mu_*$ be the minimizer, with left endpoint $0$ and width $L$, at a configuration of
\cref{fil:thm:perturb} (including a strictly positive base of \cref{fil:thm:reflect}) or of
\cref{fil:thm:split}. For probability measures $\nu$ on $[0,L]$ with $\eps(\nu)\le1$,
$W_1(\nu,\mu_*)\le C\eps(\nu)^\alpha$ with $\alpha=1$ for $n=1$ and $\alpha=\tfrac12$ for $n\ge2$; these
exponents are best possible, and after shrinking $U$ the constant is uniform on $U$. For the minimizers
of \cref{fil:cor:cluster} with $N\delta_s<\kappa$, of \cref{fil:cor:pair} with $m=2$, $\beta=0$,
$\gamma<3\alpha$, and of \cref{fil:thm:split}, the quantity
$e(d)=\inf\set{\max_{\lambda,k}\abs{\int t^ke^{i\lambda t}\dd\nu}:\supp\nu\subset[d/2,L-d/2]}$ satisfies
$cd\le e(d)\le C'd$ for small $d>0$.
\end{theorem}

\begin{proof}
Let $g$ be the certificate and $G$ its zero set. For $n\ge2$, $g(u)\ge c\operatorname{dist}(u,G)^2$ on
$[0,L]$ (linear in the distance for $n=1$), and $\int g\dd\nu\le\norm v_*\eps$. Moving each point of $\nu$
to a nearest node costs $O(\sqrt\eps)$ in $W_1$ by Cauchy--Schwarz (or $O(\eps)$), and moves the moments
by as much; the mass row and the moment columns at the nodes are injective, so the rounded weights are
within $O(\sqrt\eps)$ of the optimal ones, and moving mass between nodes costs at most $L$ times that.
Uniformity on $U$ follows because the certificates vary continuously and the derivative bounds can be
chosen uniform. For sharpness split an interior atom into two half-atoms at $\pm a$: the moments change by
$O(a^2)$ while the transport cost is at least the mass times $a$ (test with the $1$-Lipschitz
$\operatorname{dist}(\cdot,G)$); for $n=1$ move an endpoint atom inward by $a$. For $e(d)$, round as
before: when $e$ is small the rounded endpoint weights stay above half of $w_0,w_N$, and on the cells
rounding to the endpoints $g(t)\ge c_3\min(t,L-t)$, so these cells contribute at least $c_4d$ to
$\int g\dd\nu$. The upper bound follows by pushing $\mu_*$ forward under $t\mapsto d/2+(1-d/L)t$.
\end{proof}

\begin{theorem}[Anchored shortening law; \stU]\label{fil:thm:short}
At a configuration of \cref{fil:thm:perturb}, let $M(d)$ be the minimum of $\eps(\nu)$ over probability
measures on $[0,L-d]$. Then
\[
 M(d)=\frac{w_n\,[-g'(L)]}{\norm v_*}\,d+O(d^2).
\]
At a reflection-paired base, \eqref{fil:eq:slopes} gives $w_n[-G'(L)]=\pi/(hP(1))$.
\end{theorem}

\begin{proof}
Put $a=-g'(L)>0$ and $K=w_na$. The values $Y(u_i)$ have rank $n$ with unique relation proportional to
$(w_i)$, so some $h_0=z\cdot Y$ takes the value $K$ at every node but the last and $K-a$ at the last.
Near interior nodes the quadratic lower bound gives $g+dh_0\ge dK-O(d^2)$, near $0$ the positive
derivative does the same, near $L$ the minimum on $[0,L-d]$ is $dK+O(d^2)$, and elsewhere $g$ is bounded
below; integrating, $(\norm v_*+d\norm z_*)\eps(\nu)\ge dK-O(d^2)$. Conversely choose $r$ with $v\cdot r=K$
and $\norm r=K/\norm v_*$; moving the last node by $-d$ changes the moments by $-dw_nY'(L)$, whose
product with $v$ is $dK$, and the first-order variations of the weights and interior nodes span $v^\perp$,
so the total moment can be made $dr+O(d^2)$.
\end{proof}

For the single-frequency problem the symmetric shortening also has an exact rate. With $b=mh$ in
\cref{rec:thm:fwidth}, the real coordinates $Y(u)=(\Re,\Im)(u^ke^{s_ju})_{j,k}$ and
$\eps(\mu)=\norm{\int Y\dd\mu}$, let $g=v\cdot Y$ be the certificate.

\begin{theorem}[First-order cost of shortening at one frequency; \stU]\label{rec:thm:fleak}
For $0<d<b$ let $D(d)=\min\set{\eps(\mu):\supp\mu\subset[d/2,b-d/2]}$. Then
\[
 \lim_{d\downarrow0}\frac{D(d)}d=\frac{w_0g'(0)-w_mg'(b)}{2\norm v_*}>0 .
\]
\end{theorem}

\begin{proof}
By convex duality, $D(d)=\max_{\norm z_*\le1}\min_{u\in I_d}z\cdot Y(u)$. A direction with $z\cdot Y\ge0$
on $I$ integrates to zero against $\mu_*$, so it vanishes at the nodes with zero derivative inside;
writing $z\cdot Y=p\cos\tau u+r\sin\tau u$ with $p,r$ in the real span $V$ of $u^ke^{\beta_ju}$ forces
$p=0$ and $r\in\R_{\ge0}q$, where $q\in V$ vanishes at $h,\dots,(m-1)h$, so the supporting directions form
the ray of $v$. With $a_0=g'(0)/2$, $a_m=-g'(b)/2$, $a_\ell=0$ otherwise and $K=\sum w_\ell a_\ell>0$,
evaluating near-maximizers on the grid with the two endpoints moved inward gives $\limsup D(d)/d\le K$.
Conversely, grid evaluation $V\to\R^{m+1}$, $p\mapsto((-1)^\ell p(\ell h))$, has image the hyperplane
orthogonal to $(w_\ell)$, so one can prescribe $(-1)^\ell p(\ell h)=K-a_\ell$; then $g+dp\cos\tau u$ is at
least $dK-O(d^2)$ on $I_d$ and has dual norm $1+O(d)$.
\end{proof}

For $\tau=\pi$, $(\beta_1,\beta_2)=(0,\log2)$ and simple conditions the weights are $(1/3,1/2,1/6)$, the
certificate is $(2-2^u)\sin\pi u$, and for the maximum of the two complex moduli $D(d)=(\pi/9)d+o(d)$.
For the consecutive-harmonic certificate of \cref{rec:thm:hblock}, $g'(0)=\pi/(h\abs{Q(0)})$,
$g'(b)=-\pi/h$ and $w_0/w_n=\abs{Q(0)}$, so the symmetric slope is $\pi w_n/(h\norm v_*)$, the same as the
anchored slope of \cref{fil:thm:short}.

The perturbation theorem gives no explicit radius. For a concrete triple a separate finite proof does.

\begin{theorem}[The first three zeta ordinates; \stU]\label{fil:thm:three}
Let
\[
 \lambda_1\in[14.1347251417,\,14.1347251418],\quad\lambda_2\in[21.0220396387,\,21.0220396388],\]
\[\lambda_3\in[25.0108575801,\,25.0108575802].
\]
For every such triple the probability measure of least
diameter with $\widehat\mu(\lambda_j)=0$ ($j=1,2,3$) is unique up to translation; centred, it is even,
with nodes $\pm t,\pm T$, inner total mass $w$ and equal masses within each pair, where $(t,T,w)$ is the
unique solution of
\[
 w\cos(\lambda_jt)+(1-w)\cos(\lambda_jT)=0\qquad(j=1,2,3)
\]
in the max-norm box of radius $10^{-8}$ about
\[
 (0.0794453460700645,\ 0.2349726063584985,\ 0.6943858328459642).
\]
Its diameter $D=2T$ satisfies $0.46994519<D<0.46994524$.
\end{theorem}

\begin{proof}
Let $f(t,T,w)$ be the three left sides and $A$ the rational inverse of a midpoint matrix for $Df(x_0)$.
Exact interval arithmetic, uniform over the frequency box, gives $\sup_X\norm{I-ADf}_\infty<4\cdot10^{-6}$
and $\norm{Af(x_0)}_\infty<5.8\cdot10^{-12}$, so $x\mapsto x-Af(x)$ contracts the box into itself and its
unique fixed point solves $f=0$, with $0<t<T$, $0<w<1$. For this solution put
$a=(\cos(\lambda_jt))_j$, $b=(-\lambda_j\sin(\lambda_jt))_j$, $c=\pm(a\times b)$ with $\sum_jc_j>0$, and
$g(u)=\sum_jc_j\cos(\lambda_ju)$. Then $g(t)=g'(t)=0$ and $g(T)=0$. Interval bounds give $g''>0$ on
$[0.075,0.084]$, $g'<0$ on $[0.230,T_{\max}]$, and $g>0$ on $[0,0.075]$ and $[0.084,0.230]$ (308 cells).
So $g\ge0$ on $[-T,T]$ with zeros exactly $\pm t,\pm T$, and a feasible measure of diameter at most $2T$
lives there. Total mass and the first cosine equation fix the pair totals, since
$\cos(\lambda_1t)>0>\cos(\lambda_1T)$, and the certified negative determinant
$\sin(\lambda_1t)\sin(\lambda_2T)-\sin(\lambda_1T)\sin(\lambda_2t)$ makes the masses within each pair equal.
Both endpoint masses are positive.
\end{proof}

The certificate uses integers and exact rationals only, with outward rounding on a grid of step
$10^{-30}$ and Taylor bounds for sine and cosine:
\begin{center}\artifact{sharp_positive_zero_filters/certify_three.py}\\
\artifact{sharp_positive_zero_filters/certificate.txt}\end{center}
The box contains the first three zeta ordinates;
it does not itself certify them. Tilted at $\beta=\tfrac12$, the least logarithmic window is
$\approx0.469945213$, scale ratio $\approx1.5999065$, against $\pi\sum_{i\le3}1/\gamma_i\approx0.497$ for three
separate two-atom filters; the centred masses are $\approx0.15280708$ at each endpoint and
$\approx0.34719292$ at each inner node.

\subsection{Robustness and deletion of prescribed modes}\label{fil:sec:robust}
\sources{sharp_positive_zero_filters; optimal_zero_filters; optimal_dirichlet_zero_filters; grh_extensions_2026-11-27_2}

A zero of order $m$ of $F_\mu$ also controls nearby exponents. For a probability measure on $[0,L]$
with $F_\mu$ vanishing to order $m$ at $\rho$, expanding $e^{\eta u}$ through degree $m-1$ gives
\begin{equation}\label{fil:eq:taylor}
 \abs{F_\mu(\rho+\eta)}\le e^{(\Re\rho+\abs\eta)_+L}\,\frac{(\abs\eta L)^m}{m!},
\end{equation}
and for a measure on $[-H,H]$ with $\widehat\mu$ vanishing to order $m$ at real $\lambda$,
$\abs{\widehat\mu(\lambda+\eta)}\le\min(1,(H\abs\eta)^m/m!)$. Two cancelled frequencies also suppress the gap
between them.

\begin{proposition}[Suppression between cancelled frequencies; \stU]\label{fil:prop:detune}
If an even probability measure on $[-L,L]$ has $\widehat\mu(\tau)=\widehat\mu(\sigma)=0$ for
$0<\tau<\sigma$, then for real $\omega$
\[
 \abs{\widehat\mu(\omega)}\le\min\Bigl\{1,\frac{L^2}2\abs{\omega-\tau}\abs{\omega-\sigma}\Bigr\},\qquad
 \abs{\widehat\mu(\omega)}\le\min\Bigl\{1,\frac{L^4}{24}\abs{(\omega^2-\tau^2)(\omega^2-\sigma^2)}\Bigr\};
\]
in particular $\abs{\widehat\mu}\le L^2(\sigma-\tau)^2/8$ on $[\tau,\sigma]$. For the minimizer of
\cref{fil:cor:pair} with $\beta=0$, centred, with $\kappa=(\alpha+\gamma)/2$, $\delta=(\gamma-\alpha)/2$ and
$T=\pi/\kappa$,
\[
 \abs{\widehat\mu(\omega)}\le\min\Bigl\{1,\Bigl[\frac{T^2\abs{\omega-\alpha}\abs{\omega-\gamma}}{2(1+\cos\delta T)}\Bigr]^m\Bigr\}.
\]
\end{proposition}

\begin{proof}
The transform is real on $\R$ and its $j$th derivative is at most $L^j$; apply the real interpolation
remainder to the linear interpolant at $\tau,\sigma$, or to the cubic interpolant at $\pm\tau,\pm\sigma$.
For the pair filter the centred transform is $\bigl((\cos T\omega+\cos\delta T)/(1+\cos\delta T)\bigr)^m$ (as
in \cref{fil:cor:cluster}), and
$\cos T\omega+\cos\delta T=2\sin(T(\omega-\alpha)/2)\sin(T(\omega-\gamma)/2)$.
\end{proof}

For $q\ge1$ let $Z_a$, $v_a$, $\Gamma_q$ be as in \eqref{zer:eq:va}, so that \eqref{zer:eq:Zexp} holds for
$t\ge\log2$ under \cref{src:thm:dgrh}.

\begin{corollary}[Deleting prescribed modes of ramp errors; \stG]\label{fil:cor:rampdelete}
Let $\mu_c$ be one of the following filters, centred, of half-width $H$: the minimizer of
\cref{rec:thm:twofreq} for two ordinates $0<\tau<\sigma$ in $\Gamma_q$ (rescaled), of \cref{fil:cor:pair} or
\cref{fil:cor:cluster} with real part $0$, or of \cref{fil:thm:split}, for positive ordinates or approximations
to them. For $t\ge H+\log2$, simultaneously for all reduced classes $a$,
\[
 \int Z_a(t+u)\dd\mu_c(u)=\sum_{\substack{\gamma\in\Gamma_q\\\gamma\ne0}}\frac{v_a(\gamma)\,\widehat\mu_c(\gamma)}{1/4+\gamma^2}e^{i\gamma t}
 +O_q\bigl((1+t+H)e^{-(t-H)/2}\bigr).
\]
Exact targets and their negatives contribute zero, whatever their multiplicities or coincidences across
characters; an ordinate within $\eta$ of a target of order $m$ has multiplier at most
$\min(1,(H\eta)^m/m!)$. The rule uses the finitely many scales $e^{t+u}$, $u\in\supp\mu_c\subset[-H,H]$, and no
translation-independent positive rule with the same cancellation and derivative requirements has a smaller
logarithmic window.
\end{corollary}

\begin{proof}
Absolute and uniform convergence in \eqref{zer:eq:Zexp} permits termwise integration, and the multiplier at
$\gamma$ is $\widehat\mu_c(\gamma)$; the remainder is bounded uniformly for $\abs u\le H$. The Taylor bound
uses $\abs{\widehat\mu_c^{(m)}}\le H^m$. Optimality is the extremal theorem in each case.
\end{proof}

No simplicity, separation from other ordinates or linear independence of zeros is used, and the
untargeted frequencies remain. For raw zero terms $c_\rho x^\rho$, as in the smoothed sums
$S_{\chi,W}(X)=\sum_n\Lambda(n)\chi(n)W(n/X)$, the tilted filter $e^{-\beta u}\dd\nu/\int e^{-\beta v}\dd\nu$
with $\beta=\tfrac12$ has the same support and orders, so the least multiplicative range
$X_{\max}/X_{\min}$ is $e^D$; for two ordinates $\alpha<\gamma\le3\alpha$ of one or two Dirichlet
$L$-functions and order $m$ at both, it is $\exp(4\pi m/(\alpha+\gamma))$, for every smooth weight at once.
A pole term is multiplied by $F_\mu(1)>0$, and a central real zero is not cancelled by a positive average.

\subsection{Optimal logarithmic support for all low zeros}\label{fil:sec:logsupport}
\sources{logarithmic_dirichlet_zero_filters_2026-11-26; grh_optimal_positive_filters; sharp_zero_filters; sharp_zeta_filters; grh_zero_filter_tradeoff; grh_bounded_support_2026-11-27}

The theorems above prescribe a fixed finite set of frequencies. The complementary problem prescribes all
zero ordinates up to a height $B$. Here $\widehat\mu(z)=\int e^{izt}\dd\mu(t)$, and a zero of order at least
$r$ at $\lambda$ means $\widehat\mu^{(j)}(\lambda)=0$ for $j<r$. For a finite symmetric multiset
$\Lambda\subset\R\setminus\set0$ write $m_\Lambda$ for its multiplicity function. A first construction,
with explicit constants for arbitrary finite multisets, colours the frequencies into separated classes.

\begin{lemma}[A positive interpolating density; \stU]\label{rec:lem:sepfilter}
Let $\Gamma\subset\R\setminus\set0$ be finite and symmetric with the distinct points of $\Gamma\cup\set0$ at least $\delta>0$
apart, $A=4\pi/\delta$, $g_A(t)=A^{-1}(1-\abs t/A)_+$ and $f_A(z)=(\sin(Az/2)/(Az/2))^2$. There is an even
probability density $w_\Gamma$ on $[-A,A]$ with $\tfrac56g_A\le w_\Gamma\le\tfrac65g_A$ and $\widehat w_\Gamma=0$ on $\Gamma$.
\end{lemma}

\begin{proof}
With $M_{\lambda\kappa}=f_A(\lambda-\kappa)$, $b_\lambda=f_A(\lambda)$ and $0\le f_A(x)\le4/(A^2x^2)$, separation gives
$\norm{M-I}_1\le\frac8{A^2\delta^2}\sum_j j^{-2}=\tfrac1{12}$ and $\norm b_1\le\tfrac1{12}$, so $a=-M^{-1}b$ has $\norm a_1\le\tfrac1{11}$, is real,
and satisfies $a_{-\lambda}=a_\lambda$ by uniqueness. Then $Q(t)=1+\sum_\lambda a_\lambda e^{-i\lambda t}\in[\tfrac{10}{11},\tfrac{12}{11}]$, its
$g_A$-weighted integral $Z$ lies in the same range, and $w_\Gamma=g_AQ/Z$ has
$\widehat w_\Gamma(\kappa)=Z^{-1}(f_A(\kappa)+\sum_\lambda a_\lambda f_A(\kappa-\lambda))=0$.
\end{proof}

Giving up the two-sided comparison with $g_A$, a cosine taper interpolates at about a third of this
half-width.

\begin{lemma}[Cosine-window interpolation; \stU]\label{fil:lem:cosine}
Let $y_*\in(0,\pi/2)$ be the solution of $y_*\cot y_*=\tfrac12$, so $y_*=1.16556\ldots$. Let
$\Gamma\subset\R\setminus\set0$ be finite and symmetric with the distinct points of $\Gamma\cup\set0$ at least $\delta>0$
apart. There is an even probability density on $[-A,A]$, $A=\pi^2/(2y_*\delta)=4.2338\ldots/\delta$, whose
Fourier transform vanishes on $\Gamma$.
\end{lemma}

\begin{proof}
The function $y\cot y$ decreases strictly from $1$ to $0$ on $(0,\pi/2)$, its derivative being
$(\sin y\cos y-y)/\sin^2y<0$, so $y_*$ is unique. Let $\alpha=\pi/(2A)$ and $u=\alpha/\delta=y_*/\pi<\tfrac12$. The even
density $g(t)=\tfrac\alpha2\cos(\alpha t)\mathbf1_{[-A,A]}(t)$ has transform $f(x)=\alpha^2\cos(Ax)/(\alpha^2-x^2)$ (removable at
$x=\pm\alpha$), so $f(0)=1$ and $\abs{f(x)}\le\alpha^2/(x^2-\alpha^2)$ for $\abs x\ge\delta$. Put $M_{\lambda\eta}=f(\lambda-\eta)$ and $b_\lambda=f(\lambda)$
for $\lambda,\eta\in\Gamma$. On either side of a point of $\Gamma\cup\set0$ the $j$-th closest point is at distance at least $j\delta$, so
the column-sum norm and the $\ell^1$ norm satisfy
\[
 \norm{M-I}_1,\ \norm b_1\le2\sum_{j\ge1}\frac{u^2}{j^2-u^2}=1-\pi u\cot(\pi u)=\tfrac12,
\]
the series identity being the logarithmic derivative of $\sin(\pi u)/(\pi u)=\prod_{j\ge1}(1-u^2/j^2)$. Hence $M$ is
invertible and $c=-M^{-1}b$ has $\norm c_1\le1$; it is real with $c_{-\lambda}=c_\lambda$ by reflection symmetry and uniqueness.
Then $Q(t)=1+\sum_{\lambda\in\Gamma}c_\lambda e^{-i\lambda t}$ is real, even and nonnegative, and it does not vanish identically on an
interval, distinct exponentials being linearly independent; so $Z=\int gQ>0$, and $gQ/Z$ is an even probability
density whose transform at $\eta\in\Gamma$ is $(b_\eta+(Mc)_\eta)/Z=0$.
\end{proof}

\begin{theorem}[A positive filter with explicit support; \stU]\label{rec:thm:finitefilter}
Let $\Lambda\subset[-B,B]\setminus\set0$ be a finite symmetric multiset whose positive elements are at least $\delta>0$,
with local occupancy $K=\sup_{x\ge0}\sum_{x<\lambda\le x+\delta,\lambda>0}m_\Lambda(\lambda)$. For every integer $N\ge0$ there is an even
probability measure $\mu$ on $[-H,H]$, $H=\pi^2K/(2y_*\delta)+N/B$, with $\ord_\lambda\widehat\mu\ge m_\Lambda(\lambda)$ on $\Lambda$ and
$\abs{\widehat\mu(\omega)}\le(\sin1)^N$ for real $\abs\omega\ge B$.
\end{theorem}

\begin{proof}
Colour the positive occurrences in nondecreasing order greedily with $K$ colours, reusing a colour only
if its last point is at distance at least $\delta$; equal occurrences get different colours, and at most $K-1$
earlier occurrences lie within distance $\delta$ of the current one. Reflect each colour class to a symmetric
set, which with $0$ is $\delta$-separated, apply \cref{fil:lem:cosine}, and convolve, which adds the vanishing
orders. Convolving further with $N$ copies of the uniform measure on $[-1/B,1/B]$ multiplies the transform by
$(\sin(\omega/B)/(\omega/B))^N$, and $\sup_{\abs x\ge1}\abs{\sin x/x}=\sin1$.
\end{proof}

With \cref{rec:lem:sepfilter} in place of \cref{fil:lem:cosine} the half-width is $4\pi K/\delta+N/B$, but each
colour factor is then within the factor $\tfrac65$ of a triangle, which the residual certificate after
\cref{rec:thm:uniformavg} uses. Replacing the $N$ uniform factors by the measure $\sigma_{B,s}$ of
\cref{fil:prop:cheb} gives half-width $\pi^2K/(2y_*\delta)+s$ and the tail bound $\operatorname{sech}(sB)$.

With $c_0=-\log\sin1$ and $N=\ceil{c_0^{-1}\log(1/\eps)}$ the tail is at most $\eps$. If $\widehat\mu$ has a zero of order
$r$ at real $\lambda$ and $\mu$ is a probability measure on $[-H,H]$, then $\int P(T+t)e^{i\lambda(T+t)}\dd\mu(t)=0$ for every
polynomial $P$ of degree below $r$ and real $T$.

\paragraph{The extremal function.}
Let $\Lambda$ be a locally finite symmetric multiset of nonzero reals whose positive counting function
(with multiplicity) satisfies, for fixed $a>0$ and $b$,
\begin{equation}\label{fil:eq:count}
 N(T)=aT\log T+bT+O(\log(T+2))\qquad(T\ge2).
\end{equation}
For $B\ge2$ and $L\ge0$ let $H_\Lambda(B,L)$ be the least $H\ge0$ for which some probability measure $\mu$ on
$[-H,H]$ has $\ord_\lambda\widehat\mu\ge m_\Lambda(\lambda)$ for $0<\abs\lambda\le B$ and
$\sup_{\abs x\ge B}\abs{\widehat\mu(x)}\le e^{-L}$; by weak compactness the infimum is attained, and
reflecting and averaging allows $\mu$ even. For a family $\mathcal F$ write $H_{\mathcal F}(B,L)$ when $\Lambda$
is its aggregate multiset of nonzero zero ordinates.

\begin{theorem}[Slit-plane attenuation inequality; \stU]\label{fil:thm:slit}
Let $\mu$ be a probability measure on $[-H,H]$ with transform $F$, and let $\Lambda_0$ be a finite symmetric
multiset in $(-B,B)\setminus\set0$ with $\ord_\lambda F\ge m_{\Lambda_0}(\lambda)$. Then
\[
 HB\ge\sum_{\lambda\in\Lambda_0}\operatorname{arcosh}\frac B{\abs\lambda}
 +\frac{2B}\pi\int_B^\infty\frac{\log(1/\abs{F(x)})}{x\sqrt{x^2-B^2}}\dd x .
\]
In particular, if $\abs{F(x)}\le e^{-L}$ for $\abs x\ge B$, then
$H\ge L/B+(2/B)\sum_{0<\lambda\le B}m_{\Lambda_0}(\lambda)\operatorname{arcosh}(B/\lambda)$.
\end{theorem}

\begin{proof}
On $D=\C\setminus((-\infty,-B]\cup[B,\infty))$ let $Q(z)=\sqrt{B^2-z^2}$, positive at $0$; the map
$w=z/(B+Q(z))$, with inverse $z=2Bw/(1+w^2)$ and $Q=B(1-w^2)/(1+w^2)$, is a conformal bijection onto the unit
disc. For $z=x+iy$, $(\Re Q)^2=\tfrac12(\abs{B^2-z^2}+B^2-x^2+y^2)\ge y^2$, so $G(w)=F(z(w))e^{-HQ(z(w))}$ is
bounded by one, $G(0)=e^{-HB}$, and each prescribed zero $\lambda$ gives a zero of the same order at
$w(\lambda)$, with $\log(1/\abs{w(\lambda)})=\operatorname{arcosh}(B/\abs\lambda)$. Jensen's formula at radii
$r<1$ and reverse Fatou for $\log\abs G\le0$ give
$-HB+\sum\log(1/\abs{w(\lambda)})\le\frac1{2\pi}\int_0^{2\pi}\log\abs{F(B/\cos\theta)}\dd\theta$, since
$\Re Q=0$ on the boundary. Changing variables on the four quadrants, with
$\int_B^\infty\dd x/(x\sqrt{x^2-B^2})=\pi/(2B)$, gives the inequality.
\end{proof}

With $x=B\sec\theta$ the attenuation term is $\frac2\pi\int_0^{\pi/2}\log(1/\abs{F(B\sec\theta)})\dd\theta$: the weight is harmonic
measure on the two rays, not a uniform average over frequencies, and the bound adds the costs of the zeros and
of the attenuation, whereas the separate bounds $HB\ge L$ and \cref{rec:lem:jensensupport} do not imply their
sum. A prescribed zero at $\pm B$ may be included; it contributes $0$. Since
$\int_0^{\pi/2}\log\sec\theta\dd\theta=\frac\pi2\log2$, a power tail $\abs{F(x)}\le e^{-L}(B/\abs x)^s$ on $\abs x\ge B$, with $L,s\ge0$,
forces $H\ge(L+s\log2)/B+(2/B)\sum_{0<\lambda\le B}m_{\Lambda_0}(\lambda)\operatorname{arcosh}(B/\lambda)$.

Under \eqref{fil:eq:count}, partial summation and $t=B\sin\theta$ give
$\frac2B\sum_{0<\lambda\le B}m_\Lambda(\lambda)\operatorname{arcosh}(B/\lambda)=2\int_0^B\frac{N(t)\dd t}{t\sqrt{B^2-t^2}}
=\pi a\log B+\pi(b-a\log2)+O_\Lambda(\log^2B/B)$, using $\int_0^{\pi/2}\log\sin\theta\dd\theta=-\frac\pi2\log2$.
The circular Jensen bound is weaker by $\pi a(1-\log2)$ in this asymptotic, but its weights are different
term by term, and it is kept:

\begin{lemma}[Jensen support bound; \stU]\label{rec:lem:jensensupport}
If a probability measure on $[-H,H]$ has Fourier transform vanishing on a finite symmetric multiset
$\Lambda\subset[-B,B]\setminus\set0$ to the given multiplicities, then
$H\ge\frac\pi{2B}\sum_{\lambda\in\Lambda}\log\frac B{\abs\lambda}=\frac\pi B\int_0^B\frac{N_\Lambda(t)}t\dd t$, with $N_\Lambda$ the positive counting
function.
\end{lemma}

\begin{proof}
$F=\widehat\mu$ has $F(0)=1$ and $\abs{F(Re^{i\theta})}\le e^{HR\abs{\sin\theta}}$, so Jensen's formula gives
$\sum_{\abs z<R}\log(R/\abs z)\le2HR/\pi$; keep the prescribed zeros and let $R\downarrow B$. Translation handles an
uncentred interval.
\end{proof}

Combined with the two-constants theorem in each half-plane, where $[-B,B]$ has harmonic measure $\tfrac12$
seen from $Be^{i\theta}$, the same argument gives $\sum_{\abs z<B}\ord_zF\log(B/\abs z)\le2HB/\pi-D/2$ if
$\abs F\le e^{-D}$ on $\abs x\ge B$; \cref{fil:thm:slit} charges the attenuation more, $D/B$ instead of
$\pi D/(4B)$. Under RH (\stG), Montgomery's proved pair-sum formula gives
$\sum_{0<\gamma\le T}m_\gamma^2\le(\tfrac43+o(1))N(T)$ over distinct ordinates
\cite{Montgomery1973Pair,ConreyGhoshGonek1998Simple}, so at least $(\tfrac56-o(1))N(T)$ ordinates are distinct, by
$m-1\le(m^2-m)/2$, and even simple cancellation at the distinct zeta ordinates needs
$H\ge(\tfrac5{12}-o(1))\log B$.

\paragraph{Near-lattice interpolation.}
The construction that matches the slit bound sorts the ordinates by their phase within long cells, so that
each class is nearly a translate of a lattice; the interpolating kernel then has half-width close to
$\pi/\delta$ instead of $4\pi/\delta$.

\begin{lemma}[Near-lattice interpolation; \stU]\label{fil:lem:phase}
Let $0<\eta\le1$, $M=\ceil{100/\eta}$, $C=M+1$, $\delta>0$ and $0\le j<M$. Let $\Gamma_+$ be a finite set of
positive frequencies with at most one point in each bin $[\delta(k+j/M),\delta(k+(j+1)/M))$, $k\ge C$, and no
other points, and $\Gamma=\Gamma_+\cup(-\Gamma_+)$. With $g$ the convolution of the uniform probability densities
on $[-\pi/\delta,\pi/\delta]$ and $[-\eta/\delta,\eta/\delta]$, there is an even probability density
$w=Z^{-1}g\,(1+\sum_{\lambda\in\Gamma}a_\lambda e^{-i\lambda t})$ on $[-A,A]$, $A=(\pi+\eta)/\delta$, with
$\sum\abs{a_\lambda}\le\tfrac13$ and $\widehat w=0$ on $\Gamma$.
\end{lemma}

\begin{proof}
$\widehat g(x)=f(x)=\operatorname{sinc}(\pi x/\delta)\operatorname{sinc}(\eta x/\delta)$. For two points of the same sign
in cells $k,\ell$ the normalized difference $y=(\lambda-\kappa)/\delta$ satisfies
$\abs{y-(k-\ell)}\le1/M$, so for $m=\abs{k-\ell}\ge1$, $\abs y\ge m/2$, $\abs{\sin\pi y}\le\pi/M$ and
$\abs{f(\lambda-\kappa)}\le1/(M\eta y^2)\le4/(M\eta m^2)$; with at most two cells at each distance, the same-sign
column sums are at most $16/(M\eta)\le0.16$. Opposite signs have $\abs y\ge C+k$, giving at most
$1/(\pi\eta(2C-1))<0.01$, and $\sum_\lambda\abs{f(\lambda)}\le2/(\pi\eta(C-1))<0.01$. Hence the matrix
$D=(f(\lambda-\kappa))$ is within $\tfrac14$ of the identity in the column-sum norm, $a=-D^{-1}(f(\lambda))_\lambda$
has $\norm a_1\le\tfrac13$, is real and even by uniqueness, and $1+\sum a_\lambda e^{-i\lambda t}\in[\tfrac23,\tfrac43]$.
\end{proof}

\paragraph{Nested-bin allocation.}
A single grid of long cells is cheap for the high frequencies but cannot hold the lowest ones, and paying
for those separately costs the excess $\log\log B$. Instead the occurrences are allocated to nested bins,
level by level, with a fixed capacity per unit length at each level; saturation of the earlier levels
absorbs the low frequencies instead of charging them twice. Assume \eqref{fil:eq:count} and write
$n(I)$ for the number of occurrences, with multiplicity, of positive elements of $\Lambda$ in an interval
$I$. There are constants $c\ge\abs{a+b}$ and $K\ge1$ such that for all $x\ge0$ and $h\ge1$
\begin{equation}\label{fil:eq:bincount}
 a\log(x+2)-c-\frac Kh\log(x+h+2)\le\frac{n([x,x+h))}h\le a\log(x+h+2)+c+\frac Kh\log(x+h+2),
\end{equation}
by integrating the derivative $a\log(x+2)+a+b$ of the model $a(x+2)\log(x+2)+bx$, which differs from
$N$ by $O(\log(x+2))$, and charging the endpoint errors to $K\log(x+h+2)$. Fix an integer
$h_0\ge\max(1,64K/a)$ and, for levels $j\ge1$,
\[
 h_j=h_02^{j-1},\qquad \eta_j=j^{-2},\qquad M_j=100j^2,\qquad \delta_j=M_jh_j,\qquad R_j=(M_j+1)\delta_j ,
\]
an integer $D$ with $D\ge16(c+a\log2+1)$, $D/a\ge\log32$ and $e^{3D/(4a)}\ge4(R_2+h_2+2)$, and
$T_j=\exp((j-\tfrac14)D/a)$. Since $R_{j+1}/R_j\le32$, these choices give
\begin{equation}\label{fil:eq:Tlarge}
 T_j\ge4(R_{j+1}+h_{j+1}+2)\qquad(j\ge1).
\end{equation}

\begin{lemma}[Bin capacities; \stU]\label{fil:lem:capacity}
\textup{(i)} If $k=j=1$, or $j\ge2$ and $k=j-1$, every interval of length $h_k$ with right endpoint at most
$T_j+h_j$ contains at most $Djh_k$ occurrences. \textup{(ii)} Every interval $[x,x+h_i)$ with
$x\ge T_{i+1}/2$ contains at least $Dih_i$ occurrences.
\end{lemma}

\begin{proof}
(i) By \eqref{fil:eq:Tlarge}, $h_j+2\le T_j$, so the right side of \eqref{fil:eq:bincount} is at most
$Dj-D/4+a\log2+c+(K/h_k)((j-\tfrac14)D/a+\log2)$. For the indicated $k$, $j/h_k\le2/h_0$, so the $D$-term of
the last summand is at most $D/32$ and $K\log2/h_k\le a\log2/64$; the choice of $D$ makes the total at most
$Dj$. (ii) Here $h_i\le x+2$, so $\log(x+h_i+2)\le\log(x+2)+\log2$, the coefficient $a-K/h_i$ of
$\log(x+2)$ is positive, and $\log(x+2)\ge(i+\tfrac34)D/a-\log2$; the left side of \eqref{fil:eq:bincount} is at
least $Di+3D/4-a\log2-c-(K/h_i)(i+\tfrac34)D/a\ge Di$, as $(i+1)/h_i\le2/h_0$.
\end{proof}

Call the occurrences below $R_1$ the base set. For $j=1,2,\dots$ partition $[0,\infty)$ into the bins
$[kh_j,(k+1)h_j)$ and, in every bin with left endpoint at least $R_j$, allocate to level $j$ up to $Dh_j$ of
the occurrences not allocated before (all of them if there are fewer); the choice within a bin is arbitrary,
and the occurrences of a repeated frequency are distinct items. The grids are nested and $R_j$ is a multiple
of $h_j$.

\begin{lemma}[Saturation and coverage; \stU]\label{fil:lem:allocate}
For every $j\ge1$: \textup{(a)} each level-$j$ bin with left endpoint at least $T_{j+1}/2$ allocates exactly
$Dh_j$ occurrences; \textup{(b)} every occurrence at height at most $T_j$ lies in the base set or is allocated
at one of the levels $1,\dots,j$.
\end{lemma}

\begin{proof}
(a) A level-$j$ bin $I$ starting at $x\ge T_{j+1}/2$ is eligible, lies beyond the base set and contains at least
$Djh_j$ occurrences by \cref{fil:lem:capacity}(ii). For each $i<j$ it is a union of level-$i$ bins starting
beyond $T_{i+1}/2$, which by induction allocated exactly $Dh_j$ occurrences of $I$ in total at level $i$. So at
least $Dh_j$ remain, and exactly $Dh_j$ are allocated.

(b) For $j=1$, an occurrence $\lambda\le T_1$ outside the base set lies in an eligible $h_1$-bin with right
endpoint at most $T_1+h_1$, which has at most $Dh_1$ occurrences by \cref{fil:lem:capacity}(i). For $j\ge2$ let
$\lambda\le T_j$ be unallocated after level $j-1$; by induction $\lambda>T_{j-1}$, so its level-$j$ bin $I$ starts
beyond $T_{j-1}-h_j\ge T_{j-1}/2$ and is eligible, with all sub-bins of levels $i\le j-2$ starting beyond
$T_{i+1}/2$. Each of the two level-$(j-1)$ halves $I'$ of $I$ has $n(I')\le Djh_{j-1}$ by
\cref{fil:lem:capacity}(i), and the levels $i\le j-2$ allocated exactly $Dh_{j-1}$ in it each, by (a), so after level
$j-1$ at most $\max(0,n(I')-D(j-1)h_{j-1})\le Dh_{j-1}$ remain in $I'$. Hence at most $Dh_j$ remain in $I$, and
level $j$ allocates all of them.
\end{proof}

\begin{theorem}[Multiplicity-preserving positive kernels with bounded excess; \stU]\label{fil:thm:cancel}
Under \eqref{fil:eq:count} there are even probability measures $\nu_B$ with $\ord_\lambda\widehat\nu_B\ge m_\Lambda(\lambda)$
for $0<\abs\lambda\le B$, half-width at most $\pi a\log B+C_\Lambda$, and
$\int e^{ru}\dd\nu_B(u)\le e^{C_\Lambda r^2/2}$ for all real $r$, with $C_\Lambda$ independent of $B$.
\end{theorem}

\begin{proof}
For each base occurrence $\lambda\le B$ take the two-atom measure $\tfrac12(\delta_{-\pi/(2\lambda)}+\delta_{\pi/(2\lambda)})$;
the base set is finite, so these contribute a bounded half-width $H_{\mathrm{base}}$ and a bounded sum
$V_{\mathrm{base}}$ of squared half-widths. Let $J=\lceil a\log B/D+\tfrac14\rceil$, so $B\le T_J$ and, by
\cref{fil:lem:allocate}(b), every occurrence up to $B$ outside the base set is allocated at a level $j\le J$.
Write each eligible level-$j$ bin index as $kM_j+\rho$ with $0\le\rho<M_j$; its left endpoint is
$\delta_j(k+\rho/M_j)$ and eligibility means $k\ge M_j+1$. Give the at most $Dh_j$ occurrences allocated in
a bin distinct colours in $\set{1,\dots,Dh_j}$. For fixed $\rho$ and colour the allocated occurrences up to $B$
have at most one point in each cell of \cref{fil:lem:phase} with $\eta=\eta_j$, $M=M_j=\ceil{100/\eta_j}$,
$\delta=\delta_j$ and $j$ there replaced by $\rho$; apply the lemma to each such class. Level $j$ has at most
$M_jDh_j=D\delta_j$ factors, each of half-width $(\pi+\eta_j)/\delta_j$, so it costs at most $D(\pi+\eta_j)$:
the cost is proportional to the incremental density $D$, not to the full density below the current height.
Let $\nu_B$ be the convolution of all these factors; the transforms multiply, so the vanishing orders add and
no spacing across classes or levels is needed. Its half-width is at most
\[
 H_{\mathrm{base}}+\pi DJ+D\sum_{j\le J}\eta_j\le\pi a\log B+H_{\mathrm{base}}+\frac{5\pi D}4+\frac{\pi^2D}6 .
\]
An even measure on $[-A,A]$ has $\int e^{ru}\le\cosh(rA)\le e^{r^2A^2/2}$, so the variance proxy is at most
$V_{\mathrm{base}}+D(\pi+1)^2\sum_{j\ge1}\delta_j^{-1}<\infty$.
\end{proof}

The bound holds for arbitrarily severe clusters of ordinates, since only the counts
\eqref{fil:eq:bincount} enter. The constants are large and the construction does not compete with direct
numerical filters at small heights.

The phase bins are what bring the constant down to $\pi a$. For the zero ordinates of a family of $d$
Dirichlet $L$-functions as in \cref{fil:cor:dirsupport}, colouring alone, by \cref{rec:thm:finitefilter} with the
separation $\delta=\sqrt{\log B}$ on $(\delta,B]$ and the two-atom factors above for the ordinates up to $\delta$, gives
only the half-width $\kappa_*d\log B+O(\sqrt{\log B}+(\log\log B)^2)$, $\kappa_*=\pi/(4y_*)=0.6738\ldots$, instead of
$\tfrac d2\log B$.

\begin{proposition}[Optimal positive stopband factor; \stU]\label{fil:prop:cheb}
For $B>0$ and $s\ge0$ there is an even probability measure $\sigma_{B,s}$ on $[-s,s]$ with transform
\[
 P_{B,s}(z)=\frac{\cos\bigl(s\sqrt{z^2-B^2}\bigr)}{\cosh(sB)}=\frac{\cosh\bigl(s\sqrt{B^2-z^2}\bigr)}{\cosh(sB)} .
\]
Among probability measures on $[-s,s]$ the least possible value of $\sup_{\abs x\ge B}\abs{F(x)}$ is
$\operatorname{sech}(sB)$, attained by $\sigma_{B,s}$, and $\int e^{ru}\dd\sigma_{B,s}\le\exp(sr^2/(2B))$.
\end{proposition}

\begin{proof}
$C_t(z)=\cos(t\sqrt{z^2-B^2})$ solves $C_t(z)=\cos(tz)+B^2\int_0^t\frac{\sin((t-v)z)}zC_v(z)\dd v$. Starting from
$\alpha_{0,t}=\tfrac12(\delta_{-t}+\delta_t)$ and putting $\alpha_{n+1,t}=B^2\int_0^tU_{t-v}*\alpha_{n,v}\dd v$, with
$U_w=\tfrac12\mathbf1_{[-w,w]}\dd u$ (transform $\sin(wz)/z$), gives even positive measures on $[-t,t]$ of mass
$(Bt)^{2n}/(2n)!$, whose sum has mass $\cosh(Bt)$ and transform $C_t$. Equivalently, $\sigma_{B,s}$ is the law of
the displacement over time $s$ of a particle with velocity $\pm1$ switching at the times of a Poisson process of
rate $B$, conditioned on equal initial and final velocities. On $\abs x\ge B$ the numerator is a cosine of a real
number. The moment bound is $\log\cosh(s\sqrt{B^2+r^2})-\log\cosh(sB)\le s(\sqrt{B^2+r^2}-B)\le sr^2/(2B)$. For
optimality, symmetrize a competitor with bound $\eps$; $G(w)=F(\sqrt{w^2+B^2})$ is even entire of exponential
type at most $s$ (as $\abs{\Im\sqrt{w^2+B^2}}\le\abs{\Im w}$), bounded by $\eps$ on $\R$, so iterated Bernstein
inequalities \cite{RahmanTariq2009Bernstein} give $\abs{G^{(k)}(0)}\le\eps s^k$ and
$1=\abs{G(iB)}\le\eps\cosh(sB)$.
\end{proof}

This is a positive continuous analogue of the equiripple Chebyshev window \cite{Dolph1946Current}.

\begin{theorem}[Optimal support with attenuation; \stU]\label{fil:thm:optsupport}
Under \eqref{fil:eq:count}, for all $B\ge2$ and $L\ge0$,
\begin{multline*}
 \frac LB+\pi a\log B+\pi(b-a\log2)-O_\Lambda\Bigl(\frac{\log^2(B+2)}B\Bigr)\le H_\Lambda(B,L)\\
 \le\frac{\operatorname{arcosh}(e^L)}B+\pi a\log B+O_\Lambda(1).
\end{multline*}
Thus $H_\Lambda(B,L)=L/B+\pi a\log B+O_\Lambda(1)$ uniformly in $L\ge0$. The upper bound is attained by an even
probability measure with $\int e^{ru}\dd\mu(u)\le\exp\bigl(\frac{r^2}2(C_\Lambda+\operatorname{arcosh}(e^L)/B^2)\bigr)$, and by an
absolutely continuous one at an extra half-width $B^{-1}$.
\end{theorem}

\begin{proof}
The lower bound is \cref{fil:thm:slit} with the asymptotic after it. For the upper bound convolve $\nu_B$ of
\cref{fil:thm:cancel} with $\sigma_{B,s}$, $s=\operatorname{arcosh}(e^L)/B$; the zeros persist and
$\abs{\widehat\nu_B}\le1$, so the tail is $e^{-L}$; add half-widths and multiply moment bounds. A further
convolution with the uniform probability on $[-B^{-1},B^{-1}]$ gives a density.
\end{proof}

\begin{corollary}[Logarithmic support for Dirichlet zero filters; \stU]\label{fil:cor:dirsupport}\label{rec:thm:logfilter}
For a fixed finite family $\mathcal F$ of $d$ primitive Dirichlet characters closed under conjugation (the
character of modulus one gives $\zeta$; repetitions allowed), with conductors $f_\chi$, and $\Lambda$ the aggregate
multiset of nonzero ordinates of their nontrivial zeros,
\[
 H_{\mathcal F}(B,L)\ge\frac LB+\frac d2\log B+\frac12\sum_{\chi\in\mathcal F}\log\frac{f_\chi}{4\pi e}-O_{\mathcal F}\Bigl(\frac{\log^2B}B\Bigr),\]
\[
 H_{\mathcal F}(B,L)=\frac LB+\frac d2\log B+O_{\mathcal F}(1)\qquad(B\ge2,\ L\ge0).
\]
In particular $H_{\mathcal F}(B,cB\log B)/\log B\to d/2+c$ for each fixed $c\ge0$: the leading half-width
constant is $d/2$ for cancellation alone and $d/2+c_0$ with the attenuation $e^{-c_0B\log B}$, $c_0=-\log\sin1$.
So there are even probability densities $\mu_B$ supported in $[-H_B,H_B]$ with
$H_B=(d/2+c_0)\log B+O_{\mathcal F}(1)$, vanishing on $\Lambda\cap[-B,B]$ to at least the aggregate
multiplicity, with $\sup_{\abs\omega\ge B}\abs{\widehat\mu_B(\omega)}\le e^{-c_0B\log B}$.
\end{corollary}

\begin{proof}
The Riemann--von Mangoldt formula gives \eqref{fil:eq:count} with $a=d/(2\pi)$ and
$b=\frac1{2\pi}\sum_\chi\log(f_\chi/(2\pi e))$ \cite{Davenport2000Multiplicative,MontgomeryVaughan2007Multiplicative,BennettMartinOBryantRechnitzer2021Counting};
apply \cref{fil:thm:optsupport} and $L\le\operatorname{arcosh}(e^L)\le L+\log2$.
\end{proof}

No simplicity, distinctness across characters or independence of ordinates is used, and the ordinates are
real by definition; \cref{src:thm:dgrh} enters only when they are used as Fourier frequencies of prime errors.
For $\zeta$ the lower bound reads $H\ge\tfrac12\log B-\tfrac12\log(4\pi)-\tfrac12-o(1)$.

\begin{corollary}[Free cutoff; \stU]\label{fil:cor:free}
As $L\to\infty$, $\inf_{B\ge2}H_\Lambda(B,L)=\pi a\log L+O_\Lambda(1)$.
\end{corollary}

\begin{proof}
By \cref{fil:thm:optsupport}, uniformly in $B\ge2$, $H_\Lambda(B,L)$ equals $L/B+\pi a\log B$ up to a bounded
error, and for large $L$ the minimum of this expression over $B$ is attained at $B=L/(\pi a)$.
\end{proof}

Finite near-lattice examples, including repeated ordinates in distinct colours, the stopband factor and the
slit bound at zeta ordinates are evaluated in
\begin{center}\artifact{grh_optimal_positive_filters/numerical_checks.json};\end{center}
a cosine-window example, a clustered multiset and the slit and circular bounds at $B=25,50,100$ for $\zeta$ are in
\begin{center}\artifact{grh_zero_filter_tradeoff/checks.json}.\end{center}
These are floating-point illustrations, not certificates.

\paragraph{Compression of a positive theta density.}
For $\zeta$ and for $L(s,\chi_{-4})$ a second construction compresses a positive density with doubly
exponential tails into a finite interval while keeping prescribed exponential-polynomial moments exactly. It
gives explicit finite-height constants and, for $\zeta$, a density uniformly close to the Riemann kernel.

\begin{lemma}[Gap-free extrapolation; \stU]\label{fil:lem:extrap}
Let $p(t)=\sum_{j}P_j(t)e^{\lambda_jt}$ with distinct $\lambda_j\in\C$, $\deg P_j<d_j$, $\sum d_j=n\ge1$ and
$\abs{\Re\lambda_j}\le a$. For $t\ge h>0$,
\[
 \abs{p(t)}\le\frac{4n^2}h\Bigl(\frac{8et}h\Bigr)^ne^{at}\int_0^h\abs{p(v)}\dd v ;
\]
if $a=0$ the right side is at most $h^{-1}(64et/h)^n\int_0^h\abs p$, since $4n^2\le8^n$. The analogous bound at
$-t$ uses $[-h,0]$.
\end{lemma}

\begin{proof}
Fix $t\ge h$, $k=\ceil{2nt/h}$, and for $v\in[0,h/2]$ put $\Delta=(t-v)/k$, $a_j=p(v+j\Delta)e^{-aj\Delta}$. The
sequence satisfies the recurrence with characteristic roots $e^{(\lambda_\ell-a)\Delta}$ (with multiplicity
$d_\ell$, coincidences allowed), all of modulus at most one, so $\sum_ja_jw^j=\sum_{j<n}b_jw^j/\prod_\ell(1-z'_\ell w)$
with $\abs{b_j}\le2^n\max_{i<n}\abs{a_i}$, and the reciprocal has coefficients at most $\binom{l+n-1}{n-1}$.
Hence $\abs{a_k}\le n2^n\binom{k+n-1}{n-1}\max_{i<n}\abs{a_i}\le n(8et/h)^n\max_{i<n}\abs{a_i}$. Multiply back by at
most $e^{at}$; the points $v+i(t-v)/k$, $i<n$, lie in $[0,h]$ with derivative at least $\tfrac12$ in $v$;
integrate over $v\in[0,h/2]$.
\end{proof}

This is related to Tur\'an-type estimates for exponential polynomials \cite{Nazarov1994Local}, but needs no
separation of the exponents.

\begin{theorem}[Positive moment compression; \stU]\label{fil:thm:compress}
Let $w$ be an even positive density with $e^{-4e^{2t}}\le w(t)\le128e^{-e^{2t}}$ for $t\ge0$, and $V$ a real
vector space of real exponential polynomials as in \cref{fil:lem:extrap} with $a\le\tfrac12$ and order $n$,
containing $1$ and closed under $p(t)\mapsto p(-t)$. With $H=\log(16\sqrt n)$ and $\delta=e^{-100n}$ there is an
even measurable $r$ on $[-H,H]$ with $\abs r\le\delta$ such that $w_H=\mathbf1_{[-H,H]}w(1+r)$ satisfies
$\int pw_H=\int pw$ for all $p\in V$, $\int w_H=\int w$ and $\int\abs{w_H-w}\le2\delta\int w$.
\end{theorem}

\begin{proof}
Put $h=H/2>1$. By \cref{fil:lem:extrap}, $(1+u/H)^n\le e^{nu/H}$, $e^{2u}\ge1+2u$ and $w\ge e^{-4e^H}$ on $[0,h]$,
the tail $\int_{\abs t>H}\abs pw$ is at most $\delta\int_{-H}^H\abs pw$: with $e^{2H}=256n$ the ratio is at most
$2n\exp(n\log(16e)+aH-256n+64\sqrt n)\le e^{-184n}$. So $\ell(p)=\int_{\abs t>H}pw$ is a functional of norm at most
$\delta$ on $V\subset L^1([-H,H],w\dd t)$; extend it by Hahn--Banach, represent it by $r\in L^\infty$ with
$\norm r_\infty\le\delta$, and symmetrize $r$. Then $w_H\ge0$ has the moments of $w$, and the mass moved is at most
$\delta\int_{-H}^Hw$ each way.
\end{proof}

The bound counts modes with polynomial multiplicity, not their separations, so it remains valid as
frequencies merge.

\begin{theorem}[Positive Fourier-jet quadrature; \stU]\label{fil:thm:quad}
Let $w$ be an even probability density with $ce^{-ae^{2t}}\le w(t)\le Ce^{-be^{2t}}$ for $t\ge0$, and let
$v_1,\dots,v_k$ be real even exponential polynomials with real exponents such that every nonzero combination of
the $v_j-\int v_jw$ has order at most $n$. There is an even atomic probability measure $\nu$ on $[-H,H]$,
$H=\tfrac12\log(D(n+1))$ with $D=D(a,b,c,C)$, with at most $k+1$ symmetric pairs of atoms at rational radii and
$\int v_j\dd\nu=\int v_jw$ for all $j$. For real positive frequencies $\lambda$ with jet lengths $m(\lambda)$
whose jets of $\widehat w$ vanish, one may take $n=2\sum_\lambda m(\lambda)$.
\end{theorem}

\begin{proof}
Let $f\ne0$ be a real even exponential polynomial of order at most $n$ with $\int fw=0$ and $f\ge0$ on $[0,H]$;
put $L=H/2\ge1$ and $I=\int_0^Lf>0$. Then $\int_0^Hfw\ge cIe^{-ae^H}$, while \cref{fil:lem:extrap} gives
$\int_H^\infty\abs fw\le\frac{CI}L\exp(n\log(128e)-be^{2H})/(2be^{2H}-n)$. Choosing $D\ge e^4$ with $a\sqrt D\le bD/2$,
$bD\ge1$ and $bD/2>\log(128e)+\max(0,\log(C/c))+2$ makes the tail smaller, a contradiction; so every nonzero
centred combination takes both signs on $[0,H]$. Hence $0$ is in the relative interior of the convex hull of
$V(t)=(v_j(t)-\int v_jw)_j$, $t\in[0,H]$, also for a fine rational mesh by uniform continuity, and
Carath\'eodory's theorem gives at most $k+1$ mesh points; symmetrize. For the jets use
$v_{\lambda,k}=t^k\cos(\lambda t+k\pi/2)$, $k<m(\lambda)$.
\end{proof}

For a general density in this class the coefficient $\tfrac12$ cannot be lowered: matching the even moment of
degree $2m$ of $w$ by a measure on $[-H,H]$ needs $H\ge(\int\abs t^{2m}w)^{1/(2m)}=(\tfrac12-o(1))\log m$, by
integrating the lower envelope over $[R,R+1]$ with $R=(\log m-2\log\log m)/2$.

\begin{lemma}[Positive theta kernels; \stU]\label{fil:lem:kernels}
Let $\Xi(z)=\xi(\tfrac12+iz)$ with $\xi(s)=\tfrac12s(s-1)\pi^{-s/2}\Gamma(s/2)\zeta(s)$, and extend evenly
\[
 \Phi(t)=\sum_{k\ge1}\bigl(4\pi^2k^4e^{9t/2}-6\pi k^2e^{5t/2}\bigr)e^{-\pi k^2e^{2t}}\qquad(t\ge0).
\]
Then $\Xi(z)=\int_\R\Phi(t)e^{izt}\dd t$ for all $z\in\C$, $\Phi>0$, and $e^{-4e^{2t}}\le\Phi(t)\le128e^{-e^{2t}}$ and
$e^{-\pi x}\le\Phi(t)\le8\pi^2e^{-x/2}$ ($x=e^{2t}$) for $t\ge0$. With
$\theta_\beta(x)=\sum_{k\ge0}(-1)^k(2k+1)e^{-\pi(2k+1)^2x/4}$ and
$\mathcal L_\beta(s)=(4/\pi)^{(s+1)/2}\Gamma((s+1)/2)L(s,\chi_{-4})$, the kernel
$K_\beta(u)=2e^{3u/2}\theta_\beta(e^{2u})$ is even, positive, $\mathcal L_\beta(\tfrac12+it)=\int K_\beta(u)e^{itu}\dd u$,
and $e^{-\pi x/4}\le K_\beta(u)\le C_\beta e^{-\pi x/8}$ for $u\ge0$. In particular $\Xi(0)>0$ and
$\mathcal L_\beta(\tfrac12)>0$.
\end{lemma}

\begin{proof}
With $q(t)=e^{t/2}\sum_ke^{-\pi k^2e^{2t}}$ the theta relation makes $e^{t/2}+2q(t)$ even, so $q'(0)=-\tfrac14$,
and splitting the Mellin integral at $1$ gives $\Xi(z)=\tfrac12-2(z^2+\tfrac14)\int_0^\infty q(t)\cos(zt)\dd t$;
two integrations by parts give $\Phi=q''-q/4$ \cite[\S2.16]{Titchmarsh1986Theory}. Each summand is positive as
$2\pi k^2e^{2t}>3$; the term $k=1$ gives the lower bounds, and dropping negative terms with
$\sum_kk^4e^{-\pi(k^2-1)x}<2$ gives the upper ones. For $\chi_{-4}$, $\theta_\beta(x)=x^{-3/2}\theta_\beta(1/x)$ (Gauss
sum $2i$, root number $1$) makes $K_\beta$ even, and for $x\ge1$ the terms $(2k+1)e^{-\pi(2k+1)^2x/4}$ decrease
with ratio $3e^{-2\pi x}<\tfrac12$, so $\tfrac12e^{-\pi x/4}\le\theta_\beta(x)\le e^{-\pi x/4}$.
\end{proof}

\begin{theorem}[An explicit filter for the zeros of $\Xi$; \stU]\label{fil:thm:zetaexplicit}
Let $M(B)$ count the zeros $z$ of $\Xi$ with $\abs z\le B$, with multiplicity. For every $B\ge2$ there is an even
probability density $\mu_B$ on $[-H_B,H_B]$, $H_B=\log(16\sqrt{M(B)+3})$, whose transform vanishes at all these zeros
to at least their multiplicities, with $F_{\mu_B}(\pm i/2)=1/(2\xi(\tfrac12))$,
$0\le\dd\mu_B/\dd t\le(1+e^{-100n})\Phi/\Xi(0)$ and $\int\abs{\dd\mu_B/\dd t-\Phi/\Xi(0)}\le2e^{-100n}$, $n=M(B)+3$.
If the zeros up to $B$ are real (as under RH, \stG), there is also an even atomic filter with at most
$N(B)+1$ symmetric pairs at rational radii and half-width $\tfrac12\log\bigl(256(2N(B)+1)\bigr)$, $N$ the positive
ordinate count. The least half-width $H_*(B)$ for these zeros satisfies
\[
 \tfrac12\log B-\tfrac12\log(2\pi)-1-O\Bigl(\frac{\log^2B}B\Bigr)\le H_*(B)\le\tfrac12\log B+\tfrac12\log\log B+O(1).
\]
\end{theorem}

\begin{proof}
The zeros of $\Xi$ satisfy $\abs{\Im z}<\tfrac12$. Apply \cref{fil:thm:compress} with $w=\Phi$ to the real-valued part of the complex
span of $1$, $e^{\pm t/2}$ and $t^je^{izt}$ ($z$ a zero in the disc, $j<\ord_z\Xi$); it has order $n=M(B)+3$ and is
closed under reflection, the zeros being symmetric under $z\mapsto-z$ and conjugation; then divide by
$\Xi(0)$; the extra moments are $\Xi(\pm i/2)/\Xi(0)=\xi(1)/\xi(\tfrac12)$. For the atomic version apply
\cref{fil:thm:quad} with $a=\pi$, $b=\tfrac12$, $C/c=8\pi^2$, where $D=256$ works:
$16\pi<64$, $256>e^4$ and $64>\log(128e)+\log(8\pi^2)+2$. The disc count is $M(B)=2N(B)+O(\log B)$, which gives the
upper bound; the lower bound is Jensen's formula at radius $B$ with
$\int_0^BM(t)\dd t/t=\frac B\pi(\log(B/2\pi)-2)+O(\log^2B)$.
\end{proof}

For complex zeros the lower bound agrees with \cref{rec:lem:jensensupport} without RH; under RH the two
formulations coincide, and then \cref{fil:cor:dirsupport} with $L=0$ improves the upper bound to
$\tfrac12\log B+O(1)$ and the lower one to $\tfrac12\log B-\tfrac12\log(4\pi)-\tfrac12-o(1)$, by
filters that are not atomic and do not approximate $\Phi$. Numerical illustrations are in
\begin{center}\artifact{sharp_zeta_filters/verification/results.json}\\
\artifact{sharp_zero_filters/examples.json}\end{center}

\begin{corollary}[Powers of $\zeta$ and $L(s,\chi_{-4})$; \stG]\label{fil:cor:zetabeta}
For fixed integers $r,s\ge0$, $r+s>0$, and the multiplicities $rm_\zeta(\gamma)+sm_\beta(\gamma)$ at the nonzero
ordinates up to $B$, there are even atomic filters with rational radii and half-width
$\frac{r+s}2(\log B+\log\log B)+O_{r,s}(1)$.
\end{corollary}

\begin{proof}
By \cref{src:thm:dgrh} the nontrivial zeros are real zeros of the completions in the Fourier variable. Apply
\cref{fil:thm:quad} to each of the kernels of \cref{fil:lem:kernels}, normalized, with all positive zeros up to
$B$, and convolve $r$ and $s$ copies; the counts are
$N_j(T)=\frac T{2\pi}\log\frac{q_jT}{2\pi}-\frac T{2\pi}+O(\log T)$, $q_\zeta=1$, $q_\beta=4$.
\end{proof}

Without the requirement of atoms at rational radii, \cref{fil:cor:dirsupport} gives half-width
$\frac{r+s}2\log B+O_{r,s}(1)$.

\subsection{One positive filter for every noncentral zero}\label{fil:sec:global}
\sources{grh_global_positive_filters_2026-11-27; grh_sharp_endpoint; ProveIt_research_2026-11-27; sharp_global_filters; positive-fourier-zeros; quadratic-turning-points-and-critical-filters}

Exact cancellation of \emph{all} noncentral zeros is impossible with compact support, by
\cref{rec:lem:jensensupport}; the right question is the fastest possible decay. Let $\mathcal F$ be a nonempty
finite multiset of primitive characters closed under conjugation, the trivial character of conductor $1$
allowed, $d=\abs{\mathcal F}$. For $\chi\ne\triv$ with conductor $q_\chi$ and parity $a_\chi\in\set{0,1}$ put
$\Lambda_\chi(s)=(q_\chi/\pi)^{(s+a_\chi)/2}\Gamma((s+a_\chi)/2)L(s,\chi)$, and $\Lambda_\triv=\xi$. Let $m_\chi(\gamma)$
be the multiplicity of $\tfrac12+i\gamma$ as a zero of $\Lambda_\chi$ and $m(\gamma)=\sum_{\chi\in\mathcal F}m_\chi(\gamma)$.

\begin{theorem}[Sharp full-spectrum localization; \stG]\label{fil:thm:global}
There are an even, strictly positive, continuous probability density $p$ on $\R$ and $C,c>0$ with
\[
 p(t)\le C\exp\bigl(-c\exp(2\abs t/d)\bigr)\qquad(t\in\R),
\]
such that $\widehat p$ is entire and $\ord_\gamma\widehat p\ge m(\gamma)$ for every $\gamma\ne0$. Conversely, if a probability
measure $\mu$ has these zeros with at least these multiplicities and
$\int\exp(b\exp(\alpha\abs t))\dd\mu(t)<\infty$ for some $b,\alpha>0$, then $\alpha\le2/d$. The exponent $2/d$ is
therefore exact, including its endpoint. No distinctness, simplicity or independence of ordinates is
assumed; central zeros are excluded because $\widehat p(0)=1$.
\end{theorem}

The construction rests on two lemmas.

\begin{lemma}[Signed-to-positive compact smoothing; \stU]\label{fil:lem:smooth}
Let $h$ be continuous, even, strictly positive for $\abs t>A$, with $\int\abs he^{u\abs t}\dd t<\infty$ for all
$u\ge0$, and suppose $\int h(t)e^{ut}\dd t>0$ for every real $u$. Then there are $\sigma,R>0$ such that
$k(t)=e^{-t^2/(2\sigma^2)}\mathbf1_{[-R,R]}(t)$ has $h*k>0$ on $\R$; a double-exponential bound
$\abs h\le Ce^{-ce^{\alpha\abs t}}$ passes to $h*k$ with the same $\alpha$. Conversely, if the Laplace transform of $h$
vanishes somewhere on $\R$, no nonzero nonnegative compactly supported kernel makes $h*k$ nonnegative and nonzero.
\end{lemma}

\begin{proof}
Take $A\ge1$, $N=\int h_-$, $P=\int_{A+1}^{A+2}h>0$ and $D$ with $e^DP>N$; let $\delta>0$ be the minimum of the
transform on $\abs u\le D+1$, and $B>A+2$ with $\int_{\abs t>B}\abs he^{(D+1)\abs t}<\delta/4$. For $\sigma$ large,
$\int_{-B}^Bh(t)e^{ut}e^{-t^2/(2\sigma^2)}\dd t>\delta/2$ for $\abs u\le D+1$. Let $R=D\sigma^2+A+2+B$. For
$\abs x\le D\sigma^2+A+2$ factor out $e^{-x^2/(2\sigma^2)}$ and use $u=x/\sigma^2$; the rest of the integral is
nonnegative since $B>A$. For $x\ge D\sigma^2+A+2$ every negative contribution has $t\le A$, so
$k(x-t)\le k(x-A)$, while for $t\in[A+1,A+2]$, $k(x-t)\ge k(x-A)e^{(x-A-1/2)/\sigma^2}\ge e^Dk(x-A)$; the
convolution is at least $k(x-A)(e^DP-N)>0$ (or positive near $t=x$ if $k(x-A)=0$). Evenness handles $x<0$.
The tail bound follows from $\abs{t-v}\ge\abs t-R$. The Laplace transform of a convolution is the product, and
that of a nonzero nonnegative function is positive wherever finite.
\end{proof}

\Cref{rec:lem:blurpos} in \cref{sec:chz} is the companion statement for a Laplace transform bounded below by a
positive constant only on an interval $\abs y<A_0$: the kernel is then positive everywhere instead of compactly supported, and $h*k$ is
bounded below by a multiple of $k$. Neither lemma implies the other.

\begin{lemma}[Theta kernels and their convolution; \stU]\label{fil:lem:tails}
For $\chi\ne\triv$ let $K_\chi(t)=2e^{(1/2+a_\chi)t}\sum_{n\ge1}\chi(n)n^{a_\chi}\exp(-\pi n^2e^{2t}/q_\chi)$, and let
$K_\triv=\Phi$ of \cref{fil:lem:kernels}. Then $\int K_\chi(t)e^{wt}\dd t=\Lambda_\chi(\tfrac12+w)$ for all $w$, and
$f=\mathop{\ast}_{\chi\in\mathcal F}K_\chi$ is continuous and real, has bilateral Laplace transform
$F(w)=\prod_{\chi\in\mathcal F}\Lambda_\chi(\tfrac12+w)$, satisfies $f(-t)=\eta f(t)$ for a sign $\eta$, and
\[
 \abs{f(t)}\le C_1\exp\bigl(-c_1\exp(2\abs t/d)\bigr),\qquad f(t)>0\ \text{for all large }t.
\]
\end{lemma}

\begin{proof}
Termwise Mellin integration gives the transform for $\Re w$ large, theta inversion
$K_\chi(-t)=\eps_\chi K_{\bar\chi}(t)$ with $\abs{\eps_\chi}=1$ gives decay on the left, and continuation gives it
everywhere; $K_\triv$ is the Riemann kernel. On the right the term $n=1$ is real and positive and the others are
relatively $O(\exp(-c_\chi e^{2t}))$, so $\abs{K_j(t)}\le Ce^{-ce^{2\abs t}}$ and $K_j(t)=v_j(t)(1+o(1))$ with
$v_j(t)\ge e^{-Ce^{2t}}$, $v_j>0$. Conjugation makes $f$ real and the product functional equation gives
$F(-w)=\eta F(w)$. For the tail write $t_1+\dots+t_d=t$; convexity gives $\sum_je^{2\abs{t_j}}\ge de^{2\abs t/d}$, and
half of the exponent gives the bound while the other half is integrable. For positivity when $d\ge2$, the region
with all $t_j\ge T$ contributes a real part at least $\tfrac12\prod_jv_j(t_j)$, hence at least
$c_2\exp(-C_2e^{2t/d})$ from the box $\abs{t_j-t/d}\le1/d$; on the complement some $t_j<T$, so
$\sum_{i\ne j}e^{2\abs{t_i}}\ge(d-1)e^{2(t-T)/(d-1)}$, which is negligible. For $d=1$ the family is
self-conjugate and the asymptotic suffices.
\end{proof}

\begin{proof}[Proof of \cref{fil:thm:global}]
Let $m=\ord_{w=0}F$ and $H(w)=F(w)/w^m$. For $m\ge1$ put $h(t)=\frac1{(m-1)!}\int_t^\infty(v-t)^{m-1}f(v)\dd v$;
since $F^{(j)}(0)=\int v^jf=0$ for $j<m$, this equals minus the integral from $-\infty$ to $t$, so $h$ has
double-exponential decay at both ends, $h^{(m)}=(-1)^mf$, and $\int h(t)e^{wt}\dd t=H(w)$. Since
$\eta=(-1)^m$, $h$ is even, and it is positive far to the right, hence far to the left. Under
\cref{src:thm:dgrh} $F$ has no real zero except possibly $0$ (the completions are nonzero at $s=0,1$), and
$F(u)>0$ for large $u$, the conjugate factors multiplying to $\abs\cdot^2$ and the real ones being positive there;
so $H>0$ on $\R$. \Cref{fil:lem:smooth} gives $k$, and $p=(h*k)/(H(0)\int k)$ has
$\widehat p(z)=\frac{H(iz)}{H(0)}\frac{\int k(v)e^{izv}\dd v}{\int k}$, which vanishes to order $m(\gamma)$ at each
$\gamma\ne0$.

Conversely let $\int\exp(be^{\alpha\abs t})\dd\mu=M_0<\infty$. Then $\widehat\mu$ is entire, and for $z=Re^{i\theta}$,
$y=R\abs{\sin\theta}$, $\abs{\widehat\mu(z)}\le M_0\exp(\sup_{r\ge0}(yr-be^{\alpha r}))$, the supremum being
$(y/\alpha)(\log(y/(b\alpha))-1)$ for $y\ge b\alpha$. So $\log\abs{\widehat\mu(Re^{i\theta})}\le\frac{R\abs{\sin\theta}}\alpha\log R+O(R)$,
and Jensen's formula gives $\sum_{0<\abs z<R}\ord_z\widehat\mu\log(R/\abs z)\le\frac2{\pi\alpha}R\log R+O(R)$. Under
\cref{src:thm:dgrh} the prescribed zeros are real with $N_+(T)=\sum_{0<\gamma\le T}m(\gamma)=\frac d{2\pi}T\log T+O(T)$,
so their contribution is $2\int_0^RN_+(u)\dd u/u=\frac d\pi R\log R+O(R)$, and $d\le2/\alpha$.
\end{proof}

\begin{remark}[What positivity needs]\label{fil:rem:positivity}
The construction of $h$ is unconditional, and a nonnegative compact smoothing of this central-zero-removed
kernel exists if and only if $F$ has no real noncentral zero: necessity is the last part of
\cref{fil:lem:smooth}, sufficiency the argument above. Full GRH is used to identify the whole nontrivial zero
divisor with real Fourier frequencies.
\end{remark}

\paragraph{The sharp coefficient.}
The exponent $2/d$ leaves the coefficient of the double-exponential moment. Circular Jensen averaging as in
the proof of \cref{fil:thm:global} gives only $b\le2\pi d/Q_*$; dividing the transform by the completed product
keeps the directional growth and gives the exact value. Put
\[
 Q_*=\Bigl(\prod_{\chi\in\mathcal F}q_\chi\Bigr)^{1/d}\ (q_\triv=1),\qquad \alpha=\frac2d,\qquad b_*=\frac{\pi d}{Q_*},
 \qquad I_b(\mu)=\int\exp\bigl(be^{\alpha\abs t}\bigr)\dd\mu(t),
\]
let $n_0$ be the number of trivial characters in $\mathcal F$, $m=\ord_{w=0}F$ and $H(w)=F(w)/w^m$ as in the proof
of \cref{fil:thm:global}, and
\[
 \kappa_{\mathcal F}=2n_0+\frac12\sum_{\chi\in\mathcal F}a_\chi-\frac d4-m .
\]
$H$ is even, real on $\R$, and under \cref{src:thm:dgrh} positive on $\R$ with zero divisor exactly the points
$i\gamma$, $\gamma\ne0$, with multiplicities $m(\gamma)$. Write $M_\mu(w)=\int e^{wt}\dd\mu(t)$; a finite $I_b(\mu)$
makes $M_\mu$ entire. A probability measure $\mu$ with $I_b(\mu)<\infty$ for some $b>0$ is a \emph{full-divisor
filter} if $\ord_\gamma\widehat\mu\ge m(\gamma)$ for every $\gamma\ne0$, that is, if $M_\mu/H$ is entire.

\begin{theorem}[Sharp endpoint coefficient; \stG]\label{fil:thm:endpoint}\label{fil:cor:endpoint}
Assume \cref{src:thm:dgrh}.
\begin{enumerate}[label=\textup{(\roman*)}]
\item If $\mu$ is a full-divisor filter and $I_b(\mu)<\infty$, then $b\le b_*=\pi d/Q_*$.
\item For every $0<b<b_*$ there is a strictly positive continuous probability density $p$ with
$\ord_\gamma\widehat p\ge m(\gamma)$ for all $\gamma\ne0$ and $I_b(p)<\infty$. Thus $b_*$ is the exact supremum, for
every family.
\item If $I_{b_*}(\mu)<\infty$ for a full-divisor filter $\mu$, then $M_\mu=PH/H(0)$ for a real polynomial $P$
with $P(0)=1$ and $P>0$ on $\R$ (so of even degree with positive leading coefficient), and
$\deg P+\kappa_{\mathcal F}<-\tfrac12$. No symmetry of $\mu$ is assumed.
\item Let $h$ be the real even function with $\int h(t)e^{wt}\dd t=H(w)$ from the proof of \cref{fil:thm:global}.
The full-divisor filters with $I_{b_*}<\infty$ are exactly the densities $P(-\partial_t)h/H(0)$ for real
polynomials $P$ with $P(0)=1$, $\deg P+\kappa_{\mathcal F}<-\tfrac12$ and $P(-\partial_t)h\ge0$ on $\R$. In
particular there is none if $\kappa_{\mathcal F}\ge-\tfrac12$, and at most one, $h/H(0)$, if
$-\tfrac52\le\kappa_{\mathcal F}<-\tfrac12$.
\end{enumerate}
\end{theorem}

Part (i) halves the bound of the Jensen argument, which averages the gamma-factor growth over all directions.
The proof of (i), (iii) and (iv) uses factorization, Stirling's formula and a sector maximum principle, and no
information on the spacing of zeros; (ii) uses a separation argument.

\begin{lemma}[Sector asymptotics; \stU]\label{fil:lem:sector}
Uniformly for $w=Re^{i\theta}$, $\abs\theta\le\pi/4$, $x=\Re w$, $y=\Im w$, as $R\to\infty$,
\[
 \log\abs{H(w)}=\frac d2\Bigl(x\log R-y\theta-x+x\log\frac{Q_*}{2\pi}\Bigr)+\kappa_{\mathcal F}\log R+O_{\mathcal F}(1),
\]
and $H(x)\sim C_{\mathcal F}x^{\kappa_{\mathcal F}}\exp\bigl[\frac x\alpha\bigl(\log\frac x{\alpha b_*}-1\bigr)\bigr]$ as $x\to+\infty$,
with $C_{\mathcal F}\ne0$. By evenness the same holds about the negative axis with $\abs x$.
\end{lemma}

\begin{proof}
In this sector every Dirichlet series at $\tfrac12+w$ is $1+o(1)$. With $c_\chi=\tfrac14+\tfrac{a_\chi}2$,
$\Gamma(w/2+c_\chi)=\sqrt{2\pi}(w/2)^{w/2+c_\chi-1/2}e^{-w/2}(1+O(R^{-1}))$; the factor
$(q_\chi/\pi)^{w/2+c_\chi}$ adds $\tfrac x2\log(q_\chi/\pi)+O(1)$, each factor $\tfrac12s(s-1)$ of $\xi$ adds
$2\log R+O(1)$, and the division by $w^m$ subtracts $m\log R$. Sum over $\chi$, using $\alpha b_*=2\pi/Q_*$.
\end{proof}

\begin{lemma}[Entire quotient; \stG]\label{fil:lem:quotient}
If $\mu$ is a full-divisor filter, then $G=M_\mu/H$ is entire and $\log^+\abs{G(w)}=O_\sigma(1+\abs w^\sigma)$ for
every $\sigma>1$. If $I_b(\mu)<\infty$, then for $u=\abs{\Re w}\ge b\alpha$
\begin{equation}\label{fil:eq:legendre}
 \abs{M_\mu(w)}\le I_b(\mu)\exp\Bigl[\frac u\alpha\Bigl(\log\frac u{b\alpha}-1\Bigr)\Bigr],
\end{equation}
and $\abs{M_\mu(w)}\le I_b(\mu)$ for $u<b\alpha$.
\end{lemma}

\begin{proof}
The bound follows by maximizing $uv-be^{\alpha v}$ over $v\ge0$; it also justifies differentiation under the
integral and gives $\log\max_{\abs w=R}\abs{M_\mu}=O(R\log R)$. The completions have order at most one, and so
has $H$. Under \cref{src:thm:dgrh} the zero divisor of $H$ is the prescribed one, so $G$ is entire. Both $M_\mu$
and $H$ are nonzero at $0$; cancelling the divisor of $H$ in their genus-one Hadamard products gives
$G(w)=e^{Aw+C}\prod_{z}E_1(w/z)$, $E_1(v)=(1-v)e^v$, over the surplus zeros of $M_\mu$, whose exponent of
convergence is at most one. For $1<\sigma<2$, $\sum_z\abs z^{-\sigma}<\infty$ and $\log\abs{E_1(v)}\le C_\sigma\abs v^\sigma$.
\end{proof}

\begin{lemma}[Two narrow cones force a polynomial; \stU]\label{fil:lem:cones}
Let $G$ be entire with $\log^+\abs{G(w)}=O_\sigma(1+\abs w^\sigma)$ for every $\sigma>1$, and let $N\ge0$ be an
integer such that for every $\eta>0$ there are $0<\theta_0<\pi/4$ and $C_\eta$ with
$\abs{G(w)}\le C_\eta(1+\abs w)^Ne^{\eta\abs{\Im w}}$ in the two cones $\abs{\arg(\pm w)}\le\theta_0$. Then $G$ is a
polynomial of degree at most $N$.
\end{lemma}

\begin{proof}
In the upper half-plane $g_\eta(w)=G(w)(w+i)^{-N}e^{i\eta w}$ is bounded in the two cones, as
$\abs{w+i}\ge\max(1,\abs w)$ there. On the sector between them, of opening $\pi-2\theta_0<\pi$, it is bounded on the
rays and has growth $\exp(O(\abs w^\sigma))$ with $1<\sigma<\pi/(\pi-2\theta_0)$, so the Phragm\'en--Lindel\"of
principle \cite{Boas1954Entire} bounds it there. Thus $g_\eta$ is bounded on the upper half-plane; on $\R$ it is bounded by
$C_0=\sup_x\abs{G(x)}/\abs{x+i}^N$, independently of $\eta$. Multiplying by $(1-i\delta w)^{-1}$ and using the
maximum principle on half-discs, then $R\to\infty$ and $\delta\downarrow0$, gives $\abs{g_\eta}\le C_0$ there. Let
$\eta\downarrow0$, and treat the lower half-plane with $(w-i)^N$ and $e^{-i\eta w}$; Cauchy's estimates conclude.
\end{proof}

\begin{lemma}[Moment series; \stU]\label{fil:lem:series}
For a probability measure $\mu$ with all real exponential moments put
$S_b=\sum_{n\ge0}\frac{b^n}{n!}\bigl(M_\mu(\alpha n)+M_\mu(-\alpha n)\bigr)$. Then $I_b(\mu)\le S_b\le2I_b(\mu)$. If
$M_\mu=PH/H(0)$ with $P$ a real polynomial and $M_\mu>0$ on $\R$, then $I_b(\mu)<\infty$ for $b<b_*$, $I_b(\mu)=\infty$
for $b>b_*$, and $I_{b_*}(\mu)<\infty$ if and only if $\deg P+\kappa_{\mathcal F}<-\tfrac12$.
\end{lemma}

\begin{proof}
$e^{be^{\alpha\abs t}}\le e^{be^{\alpha t}}+e^{be^{-\alpha t}}\le2e^{be^{\alpha\abs t}}$; expand and use Tonelli. Positivity of
$M_\mu$ and evenness of $H$ force $P$ to have even degree $k$ and positive leading coefficient, so by
\cref{fil:lem:sector} $M_\mu(\alpha n)+M_\mu(-\alpha n)\sim C'n^{k+\kappa_{\mathcal F}}(n/(eb_*))^n$ with $C'>0$, and by Stirling's
formula the $n$th term of $S_b$ is asymptotic to a positive multiple of $(b/b_*)^nn^{k+\kappa_{\mathcal F}-1/2}$. At
$b=b_*$ the series converges exactly when the exponent is below $-1$; the boundary case diverges
logarithmically.
\end{proof}

\begin{proof}[Proof of \cref{fil:thm:endpoint}\textup{(i)}, \textup{(iii)} and \textup{(iv)}]
Let $I_b(\mu)<\infty$ with $b\ge b_*$; then $I_{b_*}(\mu)<\infty$. Subtracting \cref{fil:lem:sector} from
\eqref{fil:eq:legendre} at $b_*$ cancels the terms of order $R\log R$ and the conductor-dependent linear terms:
uniformly for $\abs\theta\le\pi/4$,
\[
 \abs{G(Re^{i\theta})}\le CR^{-\kappa_{\mathcal F}}\exp\Bigl[\frac d2R\bigl(\cos\theta\log\cos\theta+\theta\sin\theta\bigr)\Bigr],
\]
and the same about the negative axis, without any symmetry of $\mu$. As $\cos\theta\log\cos\theta\le0$, the
exponential is at most $e^{(d/2)\theta_0\abs{\Im w}}$ for $\abs\theta\le\theta_0$. \Cref{fil:lem:cones} with
$N=\max(0,\lceil-\kappa_{\mathcal F}\rceil)$ and $\theta_0\le2\eta/d$, together with \cref{fil:lem:quotient}, makes $G$ a polynomial;
put $P=H(0)G$. It is real, $P(0)=1$, and $P>0$ on $\R$ because $M_\mu$ and $H$ are. If $b>b_*$,
\eqref{fil:eq:legendre} and \cref{fil:lem:sector} on the positive axis give
$\abs{P(x)}\ll x^{-\kappa_{\mathcal F}}\exp[-(x/\alpha)\log(b/b_*)]\to0$, which is impossible; this proves (i). The degree
bound in (iii) is \cref{fil:lem:series}.

For (iv), on every vertical strip of bounded width $\abs{H(c+iy)}\ll(1+\abs y)^{C}e^{-\pi d\abs y/4}$, by Stirling and the
polynomial growth of $L$-functions in strips. So $h(t)=\frac1{2\pi}\int H(iy)e^{-iyt}\dd y$ is smooth, and shifting the
contour to $\Re w=\pm c$ shows that $p_P=P(-\partial_t)h/H(0)$ is real with $\abs{p_P(t)}\ll_{P,c}e^{-c\abs t}$ for every
$c$, and has bilateral Laplace transform $PH/H(0)$. If $M_\mu=PH/H(0)$, uniqueness of Fourier transforms gives
$\mu=p_P\dd t$, and $p_P\ge0$ by continuity. Conversely, if $p_P\ge0$ and $P(0)=1$, it is a probability density
with the full divisor, and \cref{fil:lem:series} decides $I_{b_*}$. A positive polynomial has even degree, which
gives the last two assertions.
\end{proof}

For the subcritical construction (ii), expand $h$ by the first terms of its Dirichlet series. For $\chi\in\mathcal F$
put $c_\chi=\tfrac12+a_\chi$, $v_\chi(t)=2e^{c_\chi t}\exp(-\pi e^{2t}/q_\chi)$, $v=\mathop{\ast}_\chi v_\chi$,
$Rg(t)=\int_t^\infty g$, $v_m=R^mv$, $\beta_0=d^{-1}\sum_\chi c_\chi$, $\mathcal B(D)=((D^2-\tfrac14)/2)^{n_0}$ with
$D=\dd/\dd t$, and $\sum_n\alpha_nn^{-s}=\prod_\chi L(s,\chi)$ ($\Re s>1$), so $\alpha_1=1$ and $\abs{\alpha_n}$ grows at most
polynomially.

\begin{lemma}[Leading Dirichlet term; \stU]\label{fil:lem:leading}
\begin{enumerate}[label=\textup{(\alph*)}]
\item The function $h$ has the expansion $h(t)=\sum_{n\ge1}\alpha_nn^{-1/2}\,\mathcal B(D)v_m(t+\log n)$.
\item There is a positive measure $\rho$ on $[b_*,\infty)$ with $b_*\in\supp\rho$ and
$v_m(t)=e^{\beta_0t}\int e^{-ue^{\alpha t}}\dd\rho(u)$, finite for every $t$.
\item $h(t)>0$ for all large $t$, and $\lim_{t\to+\infty}e^{-\alpha t}\log h(t)=-b_*$.
\item For $0<\eps<b_*$, a compact set $K$ in $\set{s:\Re e^{-\alpha s}>0}$ and $J\in\set{1,2}$, the sum over $n\ge J$ in
\textup{(a)} with $t$ replaced by $t-s$ is bounded in absolute value by
$C\exp[-(b_*-\eps)J^\alpha\Re(e^{-\alpha s})e^{\alpha t}]$ for $s\in K$ and large real $t$, and is holomorphic in $s$.
\end{enumerate}
\end{lemma}

\begin{proof}
(a) The bilateral transform of $v_\chi$ is the gamma and conductor part of $\Lambda_\chi(\tfrac12+w)$; $R$ divides
the transform by $w$ for real $w>0$, $\mathcal B(D)$ supplies the polynomial factors of $\xi$, and translation by
$\log n$ supplies $n^{-w}$. All of this converges absolutely on a far right vertical line, and Fourier uniqueness
there gives (a). (b) In $v(t)$ write $t_\chi=t/d+y_\chi$ with $\sum y_\chi=0$; the exponent becomes
$-e^{\alpha t}\sum_\chi(\pi/q_\chi)e^{2y_\chi}$ and the weight $e^{\beta_0t}e^{\sum c_\chi y_\chi}$. By the arithmetic--geometric
mean inequality the sum $\sum_\chi(\pi/q_\chi)e^{2y_\chi}$ has minimum $b_*$ on the hyperplane, and $\rho$ is the
image of the weight under it (a point mass if $d=1$). For $m>0$,
$R^mv(t)=\int_0^\infty z^{m-1}v(t+z)\dd z/(m-1)!$ mixes $\rho$ under $u\mapsto ue^{\alpha z}$ with positive weights, which keeps
the infimum of the support. (c) Under $e^{-ue^{\alpha t}}$, $\mathcal B(D)$ acts as multiplication by a polynomial in
$ue^{\alpha t}$ of degree $2n_0$ with positive leading coefficient, positive for large $t$ uniformly in $u\ge b_*$; the
mass of $\rho$ near $b_*$ gives the lower rate and the damping of polynomials the upper one. The terms $n\ge2$
have the faster rate from $n^\alpha e^{\alpha t}$, since $(b_*-\eps)2^\alpha>b_*$ for small $\eps$, by (d). (d) Each term
carries a fixed polynomial in $n^\alpha ue^{\alpha t}$ and a power of $n$; a fraction of the exponential absorbs these
and $e^{\beta_0t}$, a further small fraction makes the sum over $n$ bounded and gives an integrable factor $e^{-u}$
against $\rho$, and the compact lower bound for $\Re e^{-\alpha s}$ makes this uniform. The sum converges locally
uniformly in $s$, hence is holomorphic.
\end{proof}

\begin{proof}[Proof of \cref{fil:thm:endpoint}\textup{(ii)}]
By \cref{fil:lem:smooth}, as in the proof of \cref{fil:thm:global}, there is a nonnegative compactly supported
$k_0$ with $h*k_0>0$. Fix $\delta>0$ and a strictly positive continuous $W$ on $\R$ equal to $h(\abs t-\delta)$
for large $\abs t$. By \cref{fil:lem:leading}(c) and evenness, $f_s(t)=h(t-s)/W(t)$ ($\abs s\le\delta$) extends
continuously to $[-\infty,+\infty]$; it vanishes at both ends for $\abs s<\delta$, while $f_\delta$ is $1$ at $+\infty$ and
$0$ at $-\infty$, and $f_{-\delta}$ the reverse.

We claim that the real span of the $f_s$ contains a strictly positive function on $[-\infty,+\infty]$. If not,
separating the span from the open cone of strictly positive continuous functions and the Riesz representation
theorem give a nonzero positive measure $\sigma$ annihilating the span. Put $\dd\lambda=\dd\sigma/W$ on $\R$; then
\begin{equation}\label{fil:eq:orth}
 \int_\R h(t-s)\dd\lambda(t)=0\qquad(\abs s<\delta).
\end{equation}
Let $r_\pm=\inf\set{r>0:\int_{\pm t\ge0}\exp(-re^{\alpha\abs t})\dd\lambda(t)<\infty}$ and $a_\pm=\alpha^{-1}\log(b_*/r_\pm)$
($a_\pm=\infty$ if $r_\pm=0$). By \cref{fil:lem:leading}(c) for $W$, $r_\pm\le b_*e^{-\alpha\delta}$, so $a_\pm\ge\delta$; by
\cref{fil:lem:leading}(d) and evenness the left side of \eqref{fil:eq:orth} is holomorphic near every point of
$(-a_-,a_+)$ and so vanishes there. Suppose $a_+<\infty$. Near $s=a_+$ the part of the integral over $t<0$ is
holomorphic, and so is the part of the terms $n\ge2$ over $t\ge0$, by \cref{fil:lem:leading}(d), because
$2^\alpha b_*e^{-\alpha a_+}>r_+$. Hence the $n=1$ part over $t\ge0$, which is $\mathcal B(-\partial_s)V(s)$ with
$V(s)=\int_{t\ge0}v_m(t-s)\dd\lambda(t)$, continues holomorphically across $a_+$, and so does $V$, as a solution of a
linear differential equation with constant coefficients and holomorphic right side. But
\cref{fil:lem:leading}(b) gives $V(s)=e^{-\beta_0s}\int e^{-z\nu'}\dd\nu(\nu')$ with $z=e^{-\alpha s}$, where $\nu$ is the
image of $e^{\beta_0t}\mathbf1_{t\ge0}\dd\lambda(t)\dd\rho(u)$ under $\nu'=ue^{\alpha t}$, a positive measure whose real Laplace
abscissa is exactly $z_0=r_+/b_*$ (above it use $u\ge b_*$; below it the mass of $\rho$ near $b_*$ forces divergence).
A Laplace integral of a positive measure cannot be continued holomorphically across its real abscissa: the
Taylor series at a point $z_1>z_0$ has nonnegative terms towards $z_2<z_0$, and monotone convergence would
make the divergent integral at $z_2$ finite. As $s\mapsto e^{-\alpha s}$ is locally invertible, this contradiction
gives $r_+=0$, and likewise $r_-=0$. Then \eqref{fil:eq:orth} holds for all real $s$, with absolute convergence
locally uniformly in $s$; integrating against $k_0(s)\dd s$ gives $\int(h*k_0)\dd\lambda=0$, so $\lambda=0$, and pairing
$\sigma$ with $f_{\pm\delta}$ removes its masses at $\pm\infty$. This contradicts $\sigma\ne0$.

So some finite combination $g=\sum_\ell c_\ell h(\cdot-s_\ell)$, $\abs{s_\ell}\le\delta$, has $g/W$ strictly positive
and bounded on $[-\infty,+\infty]$. Its transform is $H(w)\sum_\ell c_\ell e^{s_\ell w}$, so $g/\int g$ is a strictly
positive continuous probability density with the full divisor, and $g\le CW$ with \cref{fil:lem:leading}(c) gives
$I_b<\infty$ for every $b<b_*e^{-\alpha\delta}$. Choose $\delta<\alpha^{-1}\log(b_*/b)$.
\end{proof}

The separation gives no bound on the number or size of the coefficients $c_\ell$. The upper bound (i) and the
rigidity (iii) need no Riemann hypothesis if the full complex zero divisor of $H$ is prescribed instead of
the real ordinates; \cref{src:thm:dgrh} identifies the two conditions and gives $H>0$ on $\R$, which (ii) and
(iv) use. Requiring only one vanishing per distinct ordinate is a different problem, to which the division
argument does not apply.

\paragraph{Exact examples.}
For $\zeta$ and $L(s,\chi_{-4})$ the kernels of \cref{fil:lem:kernels} are positive; so is the kernel of the even
quadratic character $\chi_5$ of conductor $5$.

\begin{lemma}[A third positive kernel; \stU]\label{fil:lem:chifive}
The kernel of \cref{fil:lem:tails} for $\chi_5$,
\[
 K_5(t)=2e^{t/2}\sum_{n\ge1}\chi_5(n)e^{-\pi n^2e^{2t}/5},
\]
is even and strictly positive, with $\int K_5(t)e^{wt}\dd t=\Lambda_{\chi_5}(\tfrac12+w)$, and $\log K_5(t)=-\tfrac\pi5e^{2\abs t}+O(1+\abs t)$.
\end{lemma}

\begin{proof}
Evenness and the transform are \cref{fil:lem:tails}, the root number being $1$ (Gauss sum $\sqrt5$). For $t\ge0$,
relative to the first term the later terms have total absolute value at most
$\sum_{n\ge2}e^{-\pi(n^2-1)e^{2t}/5}\le\sum_{n\ge2}e^{-3\pi(n-1)/5}=e^{-3\pi/5}/(1-e^{-3\pi/5})<1$, using
$n^2-1\ge3(n-1)$; this also gives the asymptotic.
\end{proof}

\begin{theorem}[Exact supremum and attainment for $\zeta$, $\chi_{-4}$ and $\chi_5$; \stG]\label{fil:thm:endpointex}
Let $\mathcal F$ consist of $r$ copies of the trivial character, $s$ of $\chi_{-4}$ and $u$ of $\chi_5$, $d=r+s+u>0$.
Then $m=0$, $\kappa_{\mathcal F}=(7r+s-u)/4$, the supremum of the admissible coefficients is
$b_*=\pi d/(4^s5^u)^{1/d}$, and the normalized convolution $\mathbf h$ of the $d$ positive kernels has
$I_b(\mathbf h)<\infty$ for every $b<b_*$. The supremum is attained if and only if $u>7r+s+2$; there is exactly one
endpoint measure when $7r+s+2<u\le7r+s+10$, and there are infinitely many when $u>7r+s+10$. For $u$ copies of
$\chi_5$ alone every even degree $2j<u/4-\tfrac12$ occurs for $P$ in \cref{fil:thm:endpoint}(iii).
\end{theorem}

\begin{proof}
The kernels have positive integrals, so there is no central zero and $\mathbf h=h/H(0)$. By the leading terms,
$\mathbb Ee^{c_je^{2\abs{X_j}}}<\infty$ for independent $X_j$ with these densities and every $c_j<\pi/q_j$; for
$b<b_*$ choose $c_j=(b/b_*)\pi/q_j$, and the arithmetic--geometric mean and triangle inequalities give
$\sum_jc_je^{2\abs{X_j}}\ge b\,e^{\alpha\abs{\sum_jX_j}}$, so $I_b(\mathbf h)\le\prod_j\mathbb Ee^{c_je^{2\abs{X_j}}}<\infty$. With
$P=1$, \cref{fil:thm:endpoint}(iv) gives attainment exactly when $\kappa_{\mathcal F}<-\tfrac12$, which is $u>7r+s+2$,
and uniqueness for $\kappa_{\mathcal F}\ge-\tfrac52$. If $u\ge1$, let $p$ be the normalized $K_5$. With $x=\pi e^{2t}/5$, the
derivatives of its first term are $e^{t/2}e^{-x}R_k(x)$ with $R_0=1$, $R_{k+1}=(\tfrac12-2x)R_k+2xR_k'$, of degree $k$ and
leading coefficient $(-2)^k$, and the other terms contribute $O_k(e^{t/2}x^ke^{-4x})$. So $p^{(2j)}>0$ for large
$\abs t$, and $p_\eps=p+\eps p^{(2j)}$ is a positive probability density for small $\eps>0$, with transform
$(1+\eps w^{2j})H_5(w)$, $H_5$ the normalized transform of $p$. Replacing one factor $p$ by $p_\eps$ realizes
$P=1+\eps w^{2j}$, admissible when $2j+\kappa_{\mathcal F}<-\tfrac12$; distinct $\eps$ give distinct measures. For $j=1$ this
is $u>7r+s+10$.
\end{proof}

\begin{center}
\small
\begin{tabular}{@{}lcccl@{}}
\toprule
Family & $d$ & $\kappa_{\mathcal F}$ & $b_*$ & endpoint\\
\midrule
$\zeta$ & $1$ & $7/4$ & $\pi$ & not attained\\
$L(s,\chi_{-4})$ & $1$ & $1/4$ & $\pi/4$ & not attained\\
$\zeta\,L(s,\chi_{-4})$ & $2$ & $2$ & $\pi$ & not attained\\
$L(s,\chi_5)^2$ & $2$ & $-1/2$ & $2\pi/5$ & not attained\\
$L(s,\chi_5)^3$ & $3$ & $-3/4$ & $3\pi/5$ & unique\\
$L(s,\chi_5)^{10}$ & $10$ & $-5/2$ & $2\pi$ & unique\\
$L(s,\chi_5)^{11}$ & $11$ & $-11/4$ & $11\pi/5$ & infinitely many\\
$\zeta\,L(s,\chi_5)^{10}$ & $11$ & $-3/4$ & $11\pi/5^{10/11}$ & unique\\
\bottomrule
\end{tabular}
\end{center}

For $\zeta$ the range of admissible coefficients in $\int\exp(be^{2\abs t})\dd\mu<\infty$ is therefore exactly
$0<b<\pi$, realized by the single density $\Phi/\xi(\tfrac12)$; for $u$ copies of $\chi_5$ a new degree $2j$
appears at $u=8j+3$. The theta kernels of the real even primitive characters of conductors $8$, $12$, $13$ and
$17$ are also positive \cite{LouboutinMunsch2013Positive}; each such factor contributes $-\tfrac14$ to $\kappa_{\mathcal F}$, and
for any family made of these, $\zeta$, $\chi_{-4}$ and $\chi_5$ the normalized convolution of the kernels is the
canonical density $h/H(0)\ge0$, so by \cref{fil:thm:endpoint}(iv) the endpoint is attained exactly when
$\kappa_{\mathcal F}<-\tfrac12$. The degree and parity bookkeeping of the table is checked in
\artifact{grh_sharp_endpoint/verification.json}.

\begin{corollary}[A weighted endpoint for the zeta zeros; \stG]\label{fil:cor:zetaweighted}
For $\mathcal F=\set\triv$ and $\lambda\in\R$ put $J_\lambda(\mu)=\int\exp(\pi e^{2\abs t}-\lambda\abs t)\dd\mu(t)$. No full-divisor
filter has $J_\lambda<\infty$ for $\lambda\le\tfrac92$; for $\tfrac92<\lambda\le\tfrac{17}2$ exactly one does, the normalized
Riemann kernel $\Phi/\xi(\tfrac12)$; for $\lambda>\tfrac{17}2$ infinitely many do.
\end{corollary}

\begin{proof}
Here $\alpha=2$, $b_*=\pi$, $\kappa_{\mathcal F}=\tfrac74$. Maximizing $(u+\lambda)t-\pi e^{2t}$ instead of $ut-\pi e^{2t}$ changes the
bound \eqref{fil:eq:legendre} by $\tfrac\lambda2\log u+O(1)$, so the proof of \cref{fil:thm:endpoint}(iii) again makes
$M_\mu=PH/H(0)$ with $P$ a positive polynomial of even degree $k$. Expanding $e^{\pi e^{2\abs t}}$, the terms of
$J_\lambda$ are comparable to $\frac{\pi^n}{n!}(M_\mu(2n-\lambda)+M_\mu(\lambda-2n))$, and
$H(2n-\lambda)/H(2n)\sim(n/\pi)^{-\lambda/2}$ by \cref{fil:lem:sector}, so they are of order $n^{k+5/4-\lambda/2}$:
convergence means $\lambda>\tfrac92+2k$. Degree $0$ is $\Phi/\xi(\tfrac12)$; for degree $2$, $\Phi''/\Phi\sim4\pi^2e^{4t}$ at
both ends, so $\Phi+\eps\Phi''$ is positive for small $\eps>0$, with transform $(1+\eps w^2)\xi(\tfrac12+w)$.
\end{proof}

\paragraph{Exact tails and truncation.}
For a fixed filter put $T_\mu(H')=\mu(\set{\abs t>H'})$ and $c(\mu)=\sup\set{b:I_b(\mu)<\infty}$. Markov's inequality and
layer-cake integration give $c(\mu)=\liminf_{H'\to\infty}e^{-\alpha H'}\log(1/T_\mu(H'))$.

\begin{proposition}[Canonical tail and truncation constant; \stG]\label{fil:prop:tail}
Let $\mathbf h$ be the normalized convolution of \cref{fil:thm:endpointex}. Then
\[
 \lim_{\abs t\to\infty}-e^{-\alpha\abs t}\log\mathbf h(t)=b_*,\qquad
 \lim_{H'\to\infty}e^{-\alpha H'}\log\frac1{T_{\mathbf h}(H')}=b_*,
\]
and for $c>0$ the least $H_{x,c}$ with $T_{\mathbf h}(H_{x,c})\le x^{-c}$ satisfies
$H_{x,c}=\frac d2\log\log x+\frac d2\log(c/b_*)+o(1)$ as $x\to\infty$. No full-divisor filter has
$T_\mu(H')\le\exp(-(b_*+\eta)e^{\alpha H'})$ for all large $H'$, for any $\eta>0$.
\end{proposition}

\begin{proof}
Write $\mathbf h=h_1*\dots*h_d$ with $\log h_i(t)=-c_ie^{2\abs t}+O(1+\abs t)$, $c_i=\pi/q_i$. The arithmetic--geometric
mean inequality on $\sum t_i=t$, keeping a fraction $\eps$ of each exponent for integrability, gives
$\mathbf h(t)\le C_\eps e^{-(1-2\eps)b_*e^{\alpha\abs t}}$. For the lower bound integrate over the box
$\abs{t_i-t_i^*}\le\delta/d$ ($i<d$) about $t_i^*=t/d+\frac12\log(q_i/Q_*)$, where $c_ie^{2t_i^*}=b_*e^{\alpha t}/d$; this gives
$\mathbf h(t)\ge C_\delta\exp(-e^{2\delta}b_*e^{\alpha t}-O(1+t))$. Integrating the two bounds gives the tail, positivity
makes $T_{\mathbf h}$ continuous and strictly decreasing, and inversion gives $H_{x,c}$. The last claim is the formula
for $c(\mu)$ with \cref{fil:thm:endpoint}(i).
\end{proof}

With \cref{fil:cor:globaltrunc} this identifies, for this fixed rule, the constant in the truncation radius:
truncating at $H'$ and renormalizing changes the average of a bounded function $f$ by at most
$2\norm f_\infty T_{\mathbf h}(H')$, so error $x^{-c}$ in normalized bounded tests costs exactly
$H_{x,c}$. This is not an optimum over filters chosen depending on $x$.

For $\zeta$ alone, $\Phi/\xi(\tfrac12)$ is a full-divisor filter without central division or smoothing (a classical
positive kernel \cite{Titchmarsh1986Theory}), and so is the normalized convolution for $\zeta$ and $L(s,\chi_{-4})$. All
primitive inducers modulo $q$ give exponent $2/\varphi(q)$. The checks in
\artifact{grh_global_positive_filters_2026-11-27/}\allowbreak\artifact{examples_output.txt} are floating-point illustrations.

\paragraph{Arithmetic consequence.}
Let $W\in C_c^\infty((1,2))$ be real, $S_\chi(x)=\sum_n\Lambda(n)\chi(n)W(n/x)$, $\delta_\chi=1$ for conductor $1$ and
$0$ otherwise, and $E_\chi(t)=e^{-t/2}(S_\chi(e^t)-\delta_\chi e^t\widetilde W(1))$. The smooth explicit formula and
\cref{src:thm:dgrh} give $E_\chi(t)=-\sum_\gamma m_\chi(\gamma)\widetilde W(\tfrac12+i\gamma)e^{i\gamma t}+r_\chi(t)$ with an
absolutely and uniformly convergent series (the central term included) and $\abs{r_\chi(t)}\le C_{\chi,W}e^{-t/2}$
for all real $t$ (the trivial zeros contribute $O(1)$ to $S_\chi$, and $E_\chi$ is bounded for $t\le0$).

\begin{theorem}[Exact cancellation of the whole oscillatory spectrum; \stG]\label{fil:thm:globalprime}
Let $\mu=p\,\dd t$ with $p$ from \cref{fil:thm:global} and $M_\mu(v)=\int e^{vu}\dd\mu(u)$. For every $\chi\in\mathcal F$
and every such $W$,
\[
 \int E_\chi(t+u)\dd\mu(u)=-m_\chi(0)\widetilde W(\tfrac12)+O_{\mathcal F,W}(e^{-t/2}),\]
\[
 \int e^{-u/2}S_\chi(xe^u)\dd\mu(u)=\delta_\chi x\widetilde W(1)M_\mu(\tfrac12)-\sqrt x\,m_\chi(0)\widetilde W(\tfrac12)+O_{\mathcal F,W}(1).
\]
The rule depends only on the family; the zero-dependent oscillation is cancelled exactly, not up to a height.
\end{theorem}

\begin{proof}
By absolute summability the term at $\gamma$ is multiplied by $\widehat\mu(\gamma)$, which is $0$ for $\gamma\ne0$ and
$1$ at $0$; the remainder is at most $C_{\chi,W}e^{-t/2}M_\mu(-\tfrac12)$. Multiply by $\sqrt x$, $x=e^t$.
\end{proof}

\begin{corollary}[A polylogarithmic truncation; \stG]\label{fil:cor:globaltrunc}
For large $x$ let $H_x=\frac d2\log\log x+C_0$ with $C_0$ large depending on $\mathcal F$, and
$\mu_x=\mathbf1_{[-H_x,H_x]}\mu/\mu([-H_x,H_x])$. The second formula of \cref{fil:thm:globalprime} holds with
$\mu_x$ on the left and $M_\mu(\tfrac12)$ on the right, so only the dilations
$e^{-C_0}x(\log x)^{-d/2}\le xe^u\le e^{C_0}x(\log x)^{d/2}$ are used.
\end{corollary}

\begin{proof}
The tail bound gives $\int_{\abs u>H}e^{vu}\dd\mu\ll_v\exp(-c'e^{2H/d})$, which is $O(x^{-2})$ for $v=0,\pm\tfrac12$;
renormalization changes the bounded averages by $O(x^{-2})$.
\end{proof}

\begin{corollary}[All reduced residue classes; \stG]\label{fil:cor:globalres}
Let $\mathcal F$ consist of the primitive inducers of the characters modulo $q$, so $d=\varphi(q)$, and put
$S_a(x)=\sum_{n\equiv a\bmod q}\Lambda(n)W(n/x)$ and $B_a=\varphi(q)^{-1}\sum_{\chi\bmod q}\overline{\chi(a)}m_{\chi^*}(0)$.
The same $\mu$ works for every $a$ and $W$:
\[
 \int e^{-u/2}S_a(xe^u)\dd\mu(u)=\frac{x\widetilde W(1)}{\varphi(q)}M_\mu(\tfrac12)-\sqrt x\,B_a\widetilde W(\tfrac12)+O_{q,W}(1),
\]
and the same holds for the truncation of \cref{fil:cor:globaltrunc}.
\end{corollary}

\begin{proof}
Orthogonality expresses $S_a$ through the $S_\chi$; passing to primitive inducers changes $S_\chi(y)$ only at
powers of primes dividing $q$, at most two of each in $[y,2y]$, so by $O_{q,W}(1)$, and $M_\mu(-\tfrac12)<\infty$.
\end{proof}

\subsection{Averaging prime sums with one positive rule}\label{fil:sec:primes}
\sources{logarithmic_dirichlet_zero_filters_2026-11-26; grh_optimal_positive_filters; sharp_zero_filters; sharp_zeta_filters; grh_bounded_support_2026-11-27}

\begin{theorem}[A uniform positive averaging rule; \stU]\label{rec:thm:uniformavg}
Let $f(t)=b+\sum_{\gamma\in\Gamma}c_\gamma e^{i\gamma t}$ with $\Gamma\subset\R\setminus\set0$ locally finite and symmetric, coefficients in a
finite-dimensional normed space, $\mathcal A=\sum\norm{c_\gamma}<\infty$, and let $\mu_B$ be a probability measure with
$\widehat\mu_B=0$ on $\Gamma\cap[-B,B]$ and $\abs{\widehat\mu_B}\le e^{-c_0B\log B}$ on $\abs\omega\ge B$. Then
$\sup_t\norm{\int f(t+u)\dd\mu_B(u)-b}\le\mathcal Ae^{-c_0B\log B}$. If $H_B\le C\log B$, one rule of half-width
$O(\log\log(\mathcal A/\eps))$ has error at most $\eps$ at every centre.
\end{theorem}

\begin{proof}
The difference from $b$ is $\sum_{\abs\gamma>B}c_\gamma e^{i\gamma t}\widehat\mu_B(\gamma)$. For the second claim take
$B=2L/(c_0\log L)$, $L=\log(\mathcal A/\eps)$.
\end{proof}

By \cref{rec:thm:logfilter} the hypothesis $H_B\le C\log B$ holds with $C=d/2+c_0+O(1/\log B)$. If each low frequency
is known only to within $\eta$ in a factor of half-width at most $A$, the bound becomes
$\mathcal A(A\eta+e^{-c_0B\log B})$, by $\abs{\phi(\omega)-\phi(\lambda)}\le A\abs{\omega-\lambda}$ for a probability measure on $[-A,A]$;
a numerical interpolation vector $\widetilde a$ with residual $r=M\widetilde a+b$ in \cref{rec:lem:sepfilter} satisfies
$\norm{\widetilde a-a}_1\le\tfrac{12}{11}\norm r_1$, which certifies finite kernels from interval data. Exact
cancellation of $t^je^{i\gamma t}$, $j<m$, at the low frequencies also persists, by expanding $(t+u)^j$.

The explicit formula can be continued to any fixed power accuracy. In the notation of \cref{zer:lem:ramp}, for
$\chi$ modulo $q$ with primitive inducer $\chi^*$ of conductor $f_{\chi^*}$ put
\[
 E_\chi(x)=\sum_{\substack{p\mid q\\p\nmid f_{\chi^*}}}\sum_{p^k\le x}\chi^*(p)^k\log p\,(p^{-k}-x^{-1}),\qquad
 E_\chi^\infty=\sum_{\substack{p\mid q\\p\nmid f_{\chi^*}}}\frac{\chi^*(p)\log p}{p-\chi^*(p)} ;
\]
let $-L'(s,\chi^*)/L(s,\chi^*)=r_{-1}/s+r_0+O(s)$ at $s=0$, $P_\chi(u)=-r_0-r_{-1}(1+u)$, $n_\chi(k)=\ord_{s=-k}L(s,\chi^*)$
for $k\ge1$, and for an integer $m\ge0$
\[
 C_{\chi,m}(x)=E_\chi^\infty-E_\chi(x)+x^{-1}P_\chi(\log x)-\sum_{k=1}^m\frac{n_\chi(k)}{k(k+1)}x^{-k-1},
\]
\[
 R^{\mathrm{corr}}_{a,m}(x)=R_a(x)-\frac1{n_q}\sum_\chi\overline{\chi(a)}C_{\chi,m}(x),
\]
with $R_a(x)=S_a(x)-\log x/n_q-c_a$. For $\zeta$, $P_\chi=\log(2\pi)$ and the Euler correction is empty.

\begin{lemma}[Corrected ramp formula; \stG]\label{fil:lem:rampcorr}
Under \cref{src:thm:dgrh}, for $x\ge2$,
$S_\chi(x)-\delta_\chi\log x-c_\chi-C_{\chi,m}(x)=x^{-1/2}K_\chi(\log x)+O_{q,m}(x^{-m-3/2})$.
\end{lemma}

\begin{proof}
For primitive $\chi$ shift the contour in the proof of \cref{zer:lem:ramp} to $\Re s=-m-\tfrac12$. On the new line
$-L'/L=O_{q,m}(\log(2+\abs{\Im s}))$, the horizontal segments can be placed at heights $\gg_q1/\log(2+\abs t)$ from all
ordinates, and the new vertical integral is $O_{q,m}(x^{-m-3/2})$. The residue at $1$ is $\delta_\chi\log x+c_\chi$; a
zero $\tfrac12+i\gamma$ of order $n$ contributes $nx^{-1/2+i\gamma}/(\tfrac14+\gamma^2)$; at $0$ the kernel is
$-x^{-1}(s^{-1}+1+\log x+O(s))$, giving $x^{-1}P_\chi(\log x)$; at $-k$ the residue is
$-n_\chi(k)x^{-k-1}/(k(k+1))$. For imprimitive $\chi$, $S_\chi=S_{\chi^*}-E_\chi$ and $c_\chi=c_{\chi^*}-E_\chi^\infty$.
\end{proof}

\begin{theorem}[Simultaneous averaging at the sharp scale; \stG]\label{rec:thm:polylogavg}
Fix $q$ and $m\ge0$, and let $\beta_a=4v_a(0)$. For all large $X$ there is one even smooth probability density
$\mu_{X,m}$, independent of $a$, supported in $[-H_X,H_X]$ with
\[
 H_X\le\frac{\varphi(q)}2\log\log X+O_{q,m}(1),
\]
such that for every reduced class $a$
\[
 \int e^{u/2}R^{\mathrm{corr}}_{a,m}(Xe^u)\dd\mu_{X,m}(u)=\frac{\beta_a}{\sqrt X}+O_{q,m}(X^{-m-3/2}),
\]
and without correction
\[
 \int e^{u/2}R_a(Xe^u)\dd\mu_{X,0}(u)=\frac{\beta_a}{\sqrt X}+O_q\Bigl(\frac{1+\log X}X\Bigr).
\]
The scales lie in $[X/(C_{q,m}(\log X)^{\varphi(q)/2}),\,C_{q,m}X(\log X)^{\varphi(q)/2}]$, and at each $X$ the average
can be realized by at most $\varphi(q)+1$ scales with nonnegative weights. Central zeros are retained in
$\beta_a$; no nonvanishing at $\tfrac12$ is assumed.
\end{theorem}

\begin{proof}
Apply \cref{fil:thm:optsupport} (through \cref{fil:cor:dirsupport}) to the $\varphi(q)$ primitive inducers of the
characters modulo $q$, a conjugation-closed multiset, with $T=\log X$, $B=T$ and $L=(m+1)T$, and convolve with a
smooth even bump on $[-1/B,1/B]$. Then $L/B=m+1$ gives the bound for $H_X$, and the moment bound has variance
proxy $O_{q,m}(1)$, since $\operatorname{arcosh}(e^L)/B^2=O_m(1/T)$. By \cref{fil:lem:rampcorr}, orthogonality and
absolute convergence, the corrected average is $X^{-1/2}\sum_\gamma v_a(\gamma)e^{i\gamma T}\widehat\mu_{X,m}(\gamma)/(\tfrac14+\gamma^2)$
plus $O_{q,m}(X^{-m-3/2}\int e^{-(m+1)u}\dd\mu_{X,m})$, and the integral is bounded by the moment estimate. The mode
$\gamma=0$ gives $\beta_a/\sqrt X$, the modes $0<\abs\gamma\le B$ vanish, and the others contribute
$O_q(X^{-1/2}e^{-L})=O_q(X^{-m-3/2})$. For $m=0$, $C_{\chi,0}(x)=O_q((1+\log x)/x)$, whose average is
$O_q((1+T)/X)$, giving the uncorrected formula. For the atoms, the integral of the continuous vector
$(e^{u/2}R^{\mathrm{corr}}_{a,m}(Xe^u))_a$ lies in the convex hull of its values, and Carath\'eodory's theorem gives at
most $\varphi(q)+1$ points; this representation need not cancel the spectrum at other centres.
\end{proof}

The coefficient $\varphi(q)/2$ is sharp within this frequency-filtering model: for $L=(m+1)T$, the value used in the proof, or
$L=T/2$, and any cutoff $B$, \cref{fil:cor:free} gives at least $(\varphi(q)/2)\log T-O_{q,m}(1)$. This is no lower bound for every method of
balancing prime errors; a measure chosen at each $X$ for the actual finite vector of errors need not satisfy the
stopband and multiplicity constraints. It is a translation-independent rule and does not compete with
centre-dependent exact balancing in bounded windows (\cref{zer:sec:balancing}). With $B=\log X$ it supersedes
the consequence of the explicit sinc-product kernel of \cref{zer:thm:kernel} below, which has remainder
$O_{q,\eps}(X^{-1+\eps})$ on the scales $X\exp(\pm O_q((\log\log X)^2))$. For general $B$,
\cref{rec:thm:uniformavg} with \cref{rec:thm:logfilter} gives one rule of half-width
$(\frac d2+c_0)\log B+O(1)$ and high-frequency tail $e^{-c_0B\log B}$ for a family of $d$ characters,
against the support $O_q(\log^2B)$ and tail $K\pi^{-N_B}$ of that kernel; since $N_B\gg B\log B$, the tails
are not directly comparable. For two classes modulo $4$ the compression filters of \cref{fil:cor:zetabeta} with $\alpha=1$ attenuation
give the window $\log T+\log\log T+O(1)$ with remainder $O(T^{3/2}(\log T)^{1/2}/X)$, which this theorem
supersedes. An exact rational certificate for the clustered multiset $\set{3/2,3/2,7/4,11/4,17/5,19/5,47/10}$ and
numerical zeta filters for $B\le200$ are
\begin{center}
\artifact{logarithmic_dirichlet_zero_filters_2026-11-26/evidence/exact_certificate.json}\\
\artifact{logarithmic_dirichlet_zero_filters_2026-11-26/evidence/numerical_results.json}
\end{center}
At these small cutoffs the colouring construction is not shorter than one uniform factor per ordinate, and its
improvement is in asymptotic order.

\begin{remark}[An explicit sinc-product kernel; \stG]\label{zer:thm:kernel}
In the notation of \cref{rec:thm:polylogavg} put $F(v)=(e^{v/2}R_a(e^v))_{a}$, $\beta=(\beta_a)_a$, and let
$K=\sum_\gamma\norm{c(\gamma)}$ with $c(\gamma)=(v_a(\gamma)/(\tfrac14+\gamma^2))_a$, as in \cref{zer:cor:smoothbalance}.
For $B\ge2$ let $\mathcal Z_B$ be the multiset of the absolute values $0<\abs\gamma\le B$ of the
ordinates of all $L(s,\chi^*)$, $\chi\bmod q$, with multiplicity and both signs, $N_B=\abs{\mathcal Z_B}$,
\[
 H_B=\pi\sum_{g\in\mathcal Z_B}\frac1g,\qquad
 \mu_B=\mathop{\ast}_{g\in\mathcal Z_B}\operatorname{Unif}[-\pi/g,\pi/g].
\]
Then $\mu_B$ is a probability measure on $[-H_B,H_B]$, independent of the centre and of the class, and
for $T\ge H_B$ and fixed $0<\delta<\tfrac12$,
\[
 \Bigl\|\int F(T+w)\,\dd\mu_B(w)-\beta\Bigr\|_2\le K\pi^{-N_B}+C_{q,\delta}e^{-(1/2-\delta)(T-H_B)}.
\]
Moreover $H_B=O_q(1+\log^2B)$ and $N_B\gg B\log B$. The support is not optimal, but the kernel and its
high-frequency bound are explicit.
\end{remark}
\begin{proof}
Convolution adds support radii and preserves mass, and
\[
 \widehat\mu_B(\omega)=\prod_{g\in\mathcal Z_B}\frac{\sin(\pi\omega/g)}{\pi\omega/g}.
\]
This vanishes at every
nonzero ordinate of modulus at most $B$, and for $\abs\omega>B$ each factor has modulus at most
$g/(\pi\abs\omega)\le1/\pi$. Integrating \eqref{zer:eq:Zexp} termwise, the high frequencies
contribute at most $K\pi^{-N_B}$, and the remainder is bounded on the support. The zero count
$O_q(B\log(B+2))$ and partial summation give $\sum_{g\le B}1/g=O_q(1+\log^2B)$, with a fixed
positive minimum among the finitely many small ordinates; the zeta factor and the
Riemann--von Mangoldt formula give $N_B\gg B\log B$.
\end{proof}

For $\zeta$ alone, the compressed theta density of \cref{fil:thm:zetaexplicit} also matches the secondary term of
the ramp formula and gives a sharper remainder.

\begin{remark}[Ramp formula with secondary term; \stG]\label{fil:lem:zetaramp}
Let $R(x)=S(x)-\log x+1+\zergE$, with $S$ the ramp sum of \cref{zer:lem:mass}. Under \cref{src:thm:dgrh},
\eqref{zer:eq:zetaexact} gives, for $x\ge1$,
\[
 R(x)=x^{-1/2}\sum_\gamma\frac{m_\gamma e^{i\gamma\log x}}{1/4+\gamma^2}+\frac{\log(2\pi)}x-\mathcal T(1/x),\qquad
 0\le\mathcal T(1/x)\le x^{-3},
\]
the sum over the distinct signed zeta ordinates with multiplicities $m_\gamma$, absolutely and uniformly
convergent; the bound on $\mathcal T$ follows from its power series and $\mathcal T(1)=1-\log2$.
\end{remark}

\begin{theorem}[A short exact zero-cancelling prime average; \stG]\label{fil:thm:zetaprime}
For large $X$ put $B=4\log X$ and let $\mu_B$ be the density of \cref{fil:thm:zetaexplicit}. Then
\[
 \int e^{u/2}R(Xe^u)\dd\mu_B(u)=\frac{\log(2\pi)}{2\xi(\tfrac12)\,X}+O(X^{-3/2}),
\]
with all scales in $[X/L_X,XL_X]$, $L_X=O(\sqrt{\log X\log\log X})$. The measure is independent of the centre once
$B$ is fixed and cancels each zero contribution up to $B$.
\end{theorem}

\begin{proof}
For $\abs y\ge2$, Stirling and a polynomial bound for $\zeta$ give $\abs{\Xi(y)/\Xi(0)}\ll e^{-\abs y/2}$, so with
the total-variation bound of \cref{fil:thm:zetaexplicit}, $\sup_{\abs y\ge B}\abs{F_{\mu_B}(y)}\ll e^{-B/2}$. Insert
\cref{fil:lem:zetaramp}: the zeros up to $B$ vanish, the others give $O(X^{-1/2}e^{-B/2})=O(X^{-5/2})$, the secondary
term gives $\log(2\pi)X^{-1}F_{\mu_B}(i/2)$, and the remainder $\mathcal T(1/x)\le x^{-3}\le x^{-3/2}$ integrates to
$O(X^{-3/2}\int e^{-u}\dd\mu_B)=O(X^{-3/2})$ by the domination $\dd\mu_B\le2\Phi\,\dd u/\Xi(0)$; the crude bound $e^{H_B}$ is
not needed. Finally $e^{H_B}=16\sqrt{M(4\log X)+3}$.
\end{proof}

Matching finitely many further moments $e^{ct}$ in \cref{fil:thm:compress} (with $H=\tfrac12\log n+C_a$) and shifting
the contour farther left gives variants with any fixed number of lower-order residue terms.

\subsection{Finite positive grids and the central multiplicity}\label{fil:sec:grids}
{\raggedright\sources{grh_local_spectral_detection; grh_multiscale_prime_fluctuations; finite_scale_prime_products_research; grh_binomial_filter; grh_sharp_extensions_2026-11-27}}

For a family $\anaC$ of \cref{ana:sec:family}, with weights $b_\chi$, let $S$ be the ramp of $\Lambda a_\anaC$ and $R$ its centred
remainder (\cref{ana:def:ramp}), let $\delta=\sum_{\chi\ \mathrm{principal}}b_\chi$, and put
\[
 d(t)=\frac{e^{(t+1)/2}}{e^{1/2}-1}\bigl\{\delta+S(e^t)-S(e^{t+1})\bigr\}
 =\frac{e^{(t+1)/2}}{e^{1/2}-1}\bigl(R(e^t)-R(e^{t+1})\bigr),
\]
the normalized finite difference of \cref{ana:sec:fd} with $\lambda=e$. It is computed from
prime powers alone, with no centring constant. Let
\[
 M_\anaC=\sum_{\chi\in\anaC}b_\chi\ord_{s=1/2}L(s,\chi)=m_\anaC(\tfrac12),\qquad
 K(w)=\frac{e^{1/2}-e^w}{e^{1/2}-1},
\]
the second equality because imprimitive Euler factors do not vanish at $1/2$; and for real
$\gamma$ write $\nu_\gamma=m_\anaC(\frac12+i\gamma)$.

\begin{lemma}[Spectral expansion of the difference; \stG]\label{ana:lem:gridexp}
Uniformly for $t\ge0$,
\begin{equation}\label{ana:eq:gridexp}
 d(t)=4M_\anaC+\sum_{\gamma\ne0}\frac{K(i\gamma)\nu_\gamma}{1/4+\gamma^2}e^{i\gamma t}
 +O_\anaC(e^{-t/4}),
\end{equation}
over the distinct nonzero ordinates, and $\abs{K(i\gamma)}\le B_K:=(e^{1/2}+1)/(e^{1/2}-1)$.
\end{lemma}

\begin{proof}
Put $D(s)=-\sum_\chi b_\chi L'/L(s,\chi)$. By \cref{ana:lem:mellin} the transform $\Phi$ of $R$ is
regular at $1$, and at a nontrivial zero $\rho$ it has residue $m_\anaC(\rho)/(\rho(1-\rho))$,
which is $\nu_\gamma/(1/4+\gamma^2)$ at $\rho=\frac12+i\gamma$. By \cref{src:thm:dgrh} all
nontrivial zeros lie on the critical line. Mellin inversion for the continuous ramp and a
contour shift to $\Re s=1/4$ give
\[
 R(x)=\sum_\gamma\frac{\nu_\gamma x^{-1/2+i\gamma}}{1/4+\gamma^2}+O_\anaC(x^{-3/4}).
\]
The integral on the new line converges absolutely, since the logarithmic derivatives are
$O_\anaC(\log(2+\abs{\Im s}))$ there by the partial-fraction expansion and the local zero counts,
and the kernel decays quadratically. The horizontal edges are taken at heights avoiding the
zeros, where the logarithmic derivatives are $O_\anaC(\log^2T)$, so their integrals tend to zero.
Finite Euler factors have no zero in $\Re s>0$. Local zero counts with multiplicity are in
\cite{BennettMartinOBryantRechnitzer2021Counting}. With $r(t)=e^{t/2}R(e^t)$ one has
$d(t)=(e^{1/2}r(t)-r(t+1))/(e^{1/2}-1)$, which multiplies $e^{i\gamma t}$ by $K(i\gamma)$; the
ordinate $0$ occurs with weight $M_\anaC$ and $K(0)/(1/4)=4$.
\end{proof}

\emph{The averaging operator.} For integers $k\ge2$ and $m\ge1$ define
\[
 U_k(w)=\frac1{k+1}\sum_{j=0}^ke^{jw/k},\qquad
 \anaM_{k,m}f(t)=\sum_{j=0}^{mk}p_jf(t+j/k),\qquad
 \sum_{j=0}^{mk}p_jz^j=\Bigl(\frac{1+z+\dots+z^k}{k+1}\Bigr)^m .
\]
All $p_j$ are positive rationals with $\sum_jp_j=1$; the operator has $mk+1$ distinct taps and
support length $m$, and its multiplier on $e^{wt}$ is $U_k(w)^m$. The polynomial is the
transfer polynomial of a cascade of uniform averages, a classical digital filter
\cite{Hogenauer1981Economical}.

\begin{lemma}[Contraction away from the aliases]\label{ana:lem:contraction}
For $0<\theta\le\pi$, $k\ge2$ and $\operatorname{dist}(\gamma,2\pi k\Z)\ge\theta$,
\[
 \abs{U_k(i\gamma)}\le\exp\Bigl(-\frac{\theta^2}{6\pi^2}\Bigr).
\]
For $a>0$, $0<U_k(-a)\le v_a:=(1+e^{-a})/2<1$.
\end{lemma}

\begin{proof}
Reduce $\gamma$ modulo $2\pi k$ to $\delta'\in[-\pi k,\pi k]$, so $\abs{\delta'}\ge\theta$. If
$\abs{\delta'}\le\pi$, let $X,Y$ be independent and uniform on $\set{0,1/k,\dots,1}$. Since
$1-\cos u\ge2u^2/\pi^2$ for $\abs u\le\pi$ and $\operatorname{Var}X=(k+2)/(12k)$,
\[
 1-\abs{U_k(i\delta')}^2=\anaEx\bigl[1-\cos(\delta'(X-Y))\bigr]
 \ge\frac{2\delta'^2}{\pi^2}\anaEx(X-Y)^2=\frac{k+2}{3k\pi^2}\delta'^2\ge\frac{\delta'^2}{3\pi^2},
\]
so $\abs{U_k(i\delta')}\le e^{-\delta'^2/(6\pi^2)}$. If $\pi\le\abs{\delta'}\le\pi k$, the
geometric sum gives
\[
 \abs{U_k(i\delta')}\le\frac1{(k+1)\sin(\abs{\delta'}/(2k))}\le\frac1{(k+1)\sin(\pi/(2k))}
 \le\frac1{\sqrt2},
\]
using $\sin x\ge(2\sqrt2/\pi)x$ on $[0,\pi/4]$, and $2^{-1/2}<e^{-1/6}$. Finally $e^{-ax}$ lies
below its chord on $[0,1]$, and the uniform average has mean $1/2$.
\end{proof}

\begin{theorem}[Alias-controlled positive spectral projection; \stU]\label{ana:thm:alias}
Let $\Gamma\subset\R\setminus\set0$ be locally finite with $\abs\gamma\ge\eta>0$ on $\Gamma$, let
$\nu_\gamma\ge0$, and let $A,B,E,a>0$. Suppose that
\[
 \sum_{\gamma\in\Gamma\cap[u-\pi,u+\pi]}\nu_\gamma\le A\log(2+\abs u)\quad(u\in\R),\qquad
 \abs{c_\gamma}\le\frac{B\nu_\gamma}{1/4+\gamma^2},
\]
and that $f(t)=c_0+\sum_{\gamma\in\Gamma}c_\gamma e^{i\gamma t}+E_0(t)$ with
$\abs{E_0(t)}\le Ee^{-at}$ for $t\ge0$. Put $W=\sum_\gamma\nu_\gamma/(1/4+\gamma^2)$,
$\theta=\min(\pi,\eta)$, $\kappa=\min\set{\theta^2/(6\pi^2),-\log v_a}>0$ and $C_*=2+2\pi$. Then
$W<\infty$, and for all $k\ge2$, $m\ge1$ and $t\ge0$,
\begin{equation}\label{ana:eq:aliasbound}
 \abs{\anaM_{k,m}f(t)-c_0}\le BWe^{-\theta^2m/(6\pi^2)}+Ee^{-at}v_a^m
 +\frac{2AB}{\pi^2k^2}\bigl\{\zeta(2)\log(C_*k)-\zeta'(2)\bigr\}.
\end{equation}
In particular, with $m_k=\lceil2\log k/\kappa\rceil$,
\[
 \sup_{t\ge0}\abs{\anaM_{k,m_k}f(t)-c_0}
 \le\frac{BW+E+(2AB/\pi^2)\bigl[\zeta(2)\log(C_*k)-\zeta'(2)\bigr]}{k^2}.
\]
\end{theorem}

\begin{proof}
Summing $O(\log(2+\abs n)/(1+n^2))$ over unit intervals shows $W<\infty$, so the series converges
absolutely and may be averaged termwise, uniformly in $t$. Away from the intervals
$I_\ell=(2\pi k\ell-\theta,2\pi k\ell+\theta)$, \cref{ana:lem:contraction} gives the first term
of \eqref{ana:eq:aliasbound}. The interval $I_0$ contains no frequency, since $\theta\le\eta$.
For $\ell\ne0$ and $\gamma\in I_\ell$ one has $\abs\gamma\ge\pi k\abs\ell$, and the local count
with $\abs{U_k(i\gamma)}\le1$ gives
\[
 \sum_{\gamma\in I_\ell}\abs{c_\gamma}\abs{U_k(i\gamma)}^m\le
 \frac{AB\log(2+2\pi k\abs\ell)}{\pi^2k^2\ell^2}.
\]
Use $2+2\pi k\ell\le C_*k\ell$ for $\ell\ge1$, sum over both signs of $\ell$, and use
$\sum\ell^{-2}=\zeta(2)$ and $\sum(\log\ell)\ell^{-2}=-\zeta'(2)$. The averaged remainder is at
most $Ee^{-at}U_k(-a)^m$. The choice of $m_k$ makes both exponential factors at most $k^{-2}$.
\end{proof}

The same proof gives $O(e^{-\kappa m}+\log k/k^r)$ when
$\abs{c_\gamma}\le B\nu_\gamma(1+\abs\gamma)^{-r}$ with $r>1$, the alias constant now involving
$\zeta(r)$ and $-\zeta'(r)$. With $m\ge(r/\kappa)\log k$ and $m\asymp\log k$, the error in
$N\asymp k\log k$ taps is $O((\log N)^{r+1}/N^r)$.

\Cref{rec:thm:uniformavg} is the continuous analogue of \cref{ana:thm:alias}, and \cref{fil:thm:globalprime} and
\cref{rec:thm:polylogavg} read off the same central term with other positive rules; none of these statements
implies another.

\begin{proposition}[The alias term is necessary; \stU]\label{ana:prop:alias}
For every integer $k\ge2$, every probability measure supported on $k^{-1}\Z$ has Fourier
transform $1$ at $2\pi k$. There are symmetric spectra satisfying the count and gap hypotheses
of \cref{ana:thm:alias}, with constants independent of $k$, for which every such averaging rule
has error at least $c\log k/k^2$ at $t=0$, with an absolute $c>0$.
\end{proposition}

\begin{proof}
Put mass $M_k=\floor{\log(2+2\pi k)}$ at each of $\gamma=\pm2\pi k$, with coefficients
$M_k/(1/4+(2\pi k)^2)$, and $E_0=0$. These spectra satisfy the local count with one absolute
constant and have a fixed positive gap. Every lattice sample sees both exponentials as $1$, so
every probability average leaves their total $2M_k/(1/4+4\pi^2k^2)$ unchanged.
\end{proof}

This is a lower bound for the abstract spectral class, not a statement that Dirichlet ordinates
lie at these aliases.

\begin{theorem}[Central multiplicity from a finite grid; \stG]\label{ana:thm:grid}
For the fixed family $\anaC$, apply \cref{ana:thm:alias} to $f=d$, with the constants supplied
by \cref{ana:lem:gridexp}, the local zero count and $a=1/4$, and with
$\eta$ the least positive ordinate of the family, and set
\[
 Z_k(t)=\tfrac14\anaM_{k,m_k}d(t),\qquad N_k=(m_k+1)k+1 .
\]
\begin{enumerate}[label=\textup{(\roman*)}]
\item Uniformly for $t\ge0$,
 \[
  Z_k(t)=M_\anaC+O_\anaC(\log k/k^2)=M_\anaC+O_\anaC\bigl((\log N_k)^3/N_k^2\bigr).
 \]
 The statistic uses exactly the $N_k$ ramp values $S(e^{t+j/k})$, $0\le j\le(m_k+1)k$. At $t=0$
 it uses prime powers only up to $e^{m_k+1}=N_k^{O_\anaC(1)}$, and its logarithmic span is
 $m_k+1=O_\anaC(\log N_k)$. If $M_\anaC$ is an integer, rounding $Z_k(t)$ to the nearest
 integer is valid as soon as the right side of
 \eqref{ana:eq:aliasbound}, divided by $4$, is below $1/2$.
\item Without any information on the gap, the prescribed depth
 $m_k=\lceil\log k\,\log\log(k+e^e)\rceil$ gives, for every fixed family,
 \[
  Z_k(t)-M_\anaC=O_\anaC\Bigl(\frac{(\log N_k)^3(\log\log N_k)^2}{N_k^2}\Bigr),\qquad
  e^{m_k+1}=N_k^{O(\log\log N_k)} .
 \]
\end{enumerate}
\end{theorem}

\begin{proof}
By \cref{ana:lem:gridexp}, $d$ has the form of \cref{ana:thm:alias} with $c_0=4M_\anaC$,
$B=B_K$, $E_0=O_\anaC(e^{-t/4})$, $A$ from the local zero count, and $\eta>0$ by discreteness
of the ordinates; dividing by $4$ gives (i). The tap $j$ of $\anaM_{k,m}$ applied to $d$ uses
$S(e^{t+j/k})$ and $S(e^{t+j/k+1})$, that is, the indices $j$ and $j+k$; and
$N_k\asymp_\anaC k\log k$ converts $\log k/k^2$ into $(\log N_k)^3/N_k^2$. In (ii) the
prescribed $m_k$ exceeds $2\log k/\kappa$ for all large $k$, so both exponential terms in
\eqref{ana:eq:aliasbound} are eventually at most $k^{-2}$, and now $N_k\asymp k\log k\log\log k$.
\end{proof}

The onset in (ii) depends on the family. Neither part claims uniformity in an unbounded
conductor, or a practical cutoff without explicit values of $A$, $E$ and $\eta$.

\begin{remark}[Noise]\label{ana:rem:gridnoise}
If each observed value of $d$ has error at most $\eps$, positivity and $\sum_jp_j=1$ give an
error at most $\eps/4$ in $Z_k$. If instead the normalized ramp values
$Y_j=e^{(t+j/k)/2}S(e^{t+j/k})$ have error at most $\eps$, then
$\abs{Z_k^{\mathrm{obs}}-Z_k}\le B_K\eps/4$, because
$d(t+j/k)=\bigl[e^{1/2}Y_j-Y_{j+k}+\delta e^{(t+j/k+1)/2}\bigr]/(e^{1/2}-1)$. No bound independent
of the span is claimed for a uniform absolute error in the unnormalized values of $S$; such errors
are multiplied by the exponential factors displayed.
\end{remark}

\begin{corollary}[Individual central orders; \stG]\label{ana:cor:gridchar}
Let $\chi$ be a nonprincipal character modulo $q$ and $\chi_0$ the principal one. The real
nonnegative weight
\[
 a_\chi(n)=\Lambda(n)\bigl\{\chi_0(n)+\Re\chi(n)\bigr\}\ge0
\]
is $\Lambda a_\anaC$ for the family $\chi_0,\chi,\bar\chi$ with weights $1,\frac12,\frac12$
(repeated characters combined), and for it $M_\anaC=\ord_{s=1/2}L(s,\chi)\in\Z_{\ge0}$. Hence
\cref{ana:thm:grid} recovers the central order of an individual, possibly nonreal, character
from a finite positive grid of real prime statistics.
\end{corollary}

\begin{proof}
Conjugation gives equal central orders for $\chi$ and $\bar\chi$, and the principal factor does
not vanish at $1/2$: $\zeta(1/2)=\eta(1/2)/(1-\sqrt2)<0$, since the alternating series
$\eta(1/2)$ is positive, and the finitely many omitted Euler factors are nonzero at $1/2$.
\end{proof}

For noncentral zeros the same statistic sees all three factors. No assumption of central
nonvanishing, simple zeros, separation of distinct noncentral ordinates or rational independence
is used.

\begin{theorem}[Every fixed grid retains one-sided detection; \stU]\label{ana:thm:gridconverse}
Fix $k\ge2$ and $m\ge1$. For a family $\anaC$ as above, GRH for all the $L(s,\chi)$, $\chi\in\anaC$,
is equivalent to an eventual upper bound for $\anaM_{k,m}d(t)$, and to an eventual two-sided
bound. If one of them has a zero $\rho$ with $\Re\rho>1/2+\sigma$, $\sigma\ge0$, then
\[
 \limsup_{t\to\infty}e^{-\sigma t}\anaM_{k,m}d(t)=+\infty .
\]
\end{theorem}

\begin{proof}
The Laplace transform of $r(t)=e^{t/2}R(e^t)$ is $\Phi(w+1/2)$, whose poles in $\Re w>0$ are at
$w=\rho-1/2$ with residue $m_\anaC(\rho)/(\rho(1-\rho))$. A forward shift by $h$ multiplies the
continuation by $e^{hw}$ up to an entire function, the contribution of an initial compact
interval. So the transform of $\anaM_{k,m}d$ is $K(w)U_k(w)^m$ times that of $r$, plus an entire
function. The zeros of $U_k$ satisfy $\Re w=0$: with $z=e^{w/k}$ they are the nontrivial
$(k+1)$st roots of unity. The zeros of $K$ satisfy $\Re w=1/2$. So no pole in $0<\Re w<1/2$ is
cancelled, and at real $w$ in that interval both factors are positive. Put
$\tilde R(x)=x^{-1/2}\anaM_{k,m}d(\log x)$. Its Mellin transform is meromorphic on $\Re s>0$,
with simple poles at the zeros in $1/2<\Re\rho<1$, nonzero residues there and positive residues
at real ones. An eventual bound $\anaM_{k,m}d(t)\le Ce^{\sigma t}$ is
$\tilde R(x)\le Cx^{\sigma-1/2}$, so \cref{ana:thm:onesided} with \cref{ana:cor:forms}(a)
excludes zeros with $\Re\rho>1/2+\sigma$; at $\sigma=0$ the functional equation excludes the rest.
Under GRH, \eqref{ana:eq:gridexp} is bounded, and so is every positive finite average of it.
\end{proof}

\paragraph{A binomial grid.}
The two-step grid with incommensurable steps $1$ and $\sqrt2$ also projects onto the central
zeros, at a slower rate. Let $X$ be the group of characters modulo $q$ trivial on a subgroup
$H\le(\Z/q\Z)^\times$, a family as above with all $b_\chi=1$, so that
$\Lambda a_X=[(\Z/q\Z)^\times:H]\Lambda\anaone_{n\bmod q\in H}$ and $\delta=1$. With $T_hf(t)=f(t+h)$
and $A_h=(I+T_h)/2$, put
\[
 B_md(t)=A_1^mA_{\sqrt2}^md(t)=4^{-m}\sum_{j,k=0}^m\binom mj\binom mkd(t+j+k\sqrt2),
\]
with $(m+1)^2$ shifts and positive weights of sum one, and write $M_X=M_\anaC$.

\begin{theorem}[Binomial filter; \stU\ and \stG]\label{ana:thm:binomial}
\leavevmode
\begin{enumerate}[label=\textup{(\roman*)}]
\item \textup{(\stU)} For every fixed $m\ge0$, GRH for the characters in $X$ is equivalent to
 an upper bound for $B_md(t)$ for large $t$, and to a two-sided bound. If some $L(s,\chi)$,
 $\chi\in X$, has a zero with $\Re\rho>1/2+\sigma$, $\sigma\ge0$, then
 $\limsup_{t\to\infty}e^{-\sigma t}B_md(t)=+\infty$.
\item \textup{(\stG)} There is $C_{q,H}$ with $\abs{B_md(t)-4M_X}\le C_{q,H}\log(2m)/\sqrt m$
 for all $m\ge1$ and $t\ge0$. Hence $M_X=\frac14\lim_mB_md(0)$, which uses only prime powers
 up to $\exp(1+(1+\sqrt2)m)$ and no constant, $L$-value or ordinate, and $M_X$ is the nearest
 integer to $B_md(t)/4$ for all $t\ge0$ once $m$ is large.
\end{enumerate}
\end{theorem}

\begin{proof}
(i) is the proof of \cref{ana:thm:gridconverse} with the multiplier
$K(w)\bigl(\frac{1+e^w}2\cdot\frac{1+e^{\sqrt2w}}2\bigr)^m$, which has no zero in
$0<\Re w<1/2$ and is positive on $(0,1/2)$.

(ii) By \cref{ana:lem:gridexp} it suffices to damp the nonzero ordinates. On the imaginary axis
the Bernoulli part has modulus $\abs{\cos(\gamma/2)\cos(\sqrt2\gamma/2)}^m$. For $\abs u\ge\eta$,
\[
 \abs{\cos(u/2)\cos(\sqrt2u/2)}\le\exp\bigl(-c_\eta/(1+\abs u)^2\bigr):
\]
with $v=u/2\pi$ and nearest integers $a$ to $v$ and $b$ to $\sqrt2v$, the nonzero integer
$b^2-2a^2$ gives $1\le\abs{b^2-2a^2}\ll(\abs{a-v}+\abs{b-\sqrt2v})(1+\abs u)$, and
$\abs{\cos\pi d}\le e^{-2d^2}$ for $\abs d\le1/2$. Summing
$\exp(-c_*m/(1+\abs\gamma)^2)/(1/4+\gamma^2)$ over the ordinates, dyadically above and below
$\sqrt m$, gives $O(\log(2m)/\sqrt m)$. The error $O(e^{-t/4})$ is multiplied by at most
$\bigl(\frac{1+e^{-1/4}}2\frac{1+e^{-\sqrt2/4}}2\bigr)^m$.
\end{proof}

In terms of samples, $B_m$ uses $(m+1)^2$ values of $d$, or $(m+1)(m+2)$ distinct ramp values,
for an error $O(\log N/N^{1/4})$ in $N$ observations over logarithmic span $(1+\sqrt2)m$, while
\cref{ana:thm:grid} gives $O((\log N)^3/N^2)$ over span $O(\log N)$. The inverse-square damping
of $B_m$ cannot be improved uniformly: the solutions of $b^2-2a^2=\pm1$ give $u=2\pi a$ with
damping $1-O(a^{-2})$. Equal Bernoulli weights matter: $p+(1-p)e^{hw}$ with $p>1/2$ vanishes at
real part $h^{-1}\log(p/(1-p))>0$, which would cancel poles. If the ordinates up to $U$ are known,
the product $\prod_{0<\gamma\le U}A_{\pi/\gamma}$ annihilates them exactly and gives
$4M_X+O(\log(2U)/U)+O(e^{-t/4})$ with total shift $O(\log^2(2U))$, and the converse in (i) holds
for any fixed finite product of such operators.

\subsection{Open questions}\label{fil:sec:open}
\sources{quadratic_spectra; dirichlet_zero_filters; grh_positive_spectral_extensions_2026-11-27; logarithmic_dirichlet_zero_filters_2026-11-26; grh_global_positive_filters_2026-11-27; grh_optimal_positive_filters; optimal_dirichlet_zero_filters; optimal_zero_filters; sharp_positive_zero_filters; sharp_zero_filters; sharp_zeta_filters; grh_extensions_2026-11-27_2; grh_zero_filter_tradeoff; gram-euler-and-zero-filters; grh_cayley_and_filters; grh_hidden_forms_and_filters_2026-11-27; grh_binomial_filter; grh_sharp_endpoint; ProveIt_research_2026-11-27; sharp_global_filters; positive-fourier-zeros; grh_bounded_support_2026-11-27; erdelyi_ellipse_and_proveit_extensions; deep_padic_roots_2026-11-27; proveit_three_extensions_2026-11-27_1}

\begin{enumerate}
\item \emph{Global phase diagrams.} \Cref{rec:thm:twofreq,rec:thm:cusp} settle two arbitrary frequencies with
simple vanishing, \cref{fil:cor:pair} equal orders up to ratio $3$, and \cref{fil:thm:reflect,fil:thm:perturb,fil:thm:split}
open neighbourhoods of strictly positive configurations. What are the minimum width, the support transitions
and the fractional powers at resonances for three noncommensurate frequencies outside these neighbourhoods, for
equal orders $m\ge2$ above ratio $3$, and for unequal orders at two frequencies? \Cref{fil:thm:central} decides which
paired bands with all angles in $[\pi/3,\pi/2]$ are rescued by central roots, and \cref{fil:thm:centralphase} gives the
first-order central multiplicity for one repeated pair at any angle. What is the least central multiplicity for
several distinct angles below $\pi/3$, and what is the size of $R_m(c)-m\tau(c)$ for $c>\tfrac12$? For real
frequencies these are questions about the equality region of \cref{fil:thm:meanbound}.
\item \emph{Effective neighbourhoods and degenerate mirror configurations.} Make the neighbourhoods of
\cref{fil:thm:perturb,fil:thm:split,fil:thm:mirrorhessian} and the constants of \cref{fil:cor:sixth} explicit in terms of
root separation, minimum weight and certificate curvature, so that interval certificates like \cref{fil:thm:three} can be
produced systematically; no uniformity in the cluster size $N$ is known, and by \cref{fil:ex:mirror}\textup{(ii)} none
can hold up to a vanishing weight. \Cref{fil:thm:mirrorhessian} gives the Hessian at every mirror configuration with
strictly positive optimal weights: positive definite transverse to the reflection locus, with exactly its tangent
space as kernel. At mirror configurations where some optimal weights vanish, what are the local order of the width
excess and the support transitions? At a general base, which coordinates, if any, give a fixed sign pattern of
the minors, as the frequency coordinates do at a coalesced base (\cref{fil:thm:hessian}), given that mixed entries
change sign in the coefficient coordinates (\cref{fil:ex:mirror}\textup{(iv)})?
\item \emph{Second-order term of the support.} By \cref{fil:thm:optsupport} the excess
$H_\Lambda(B,L)-L/B-\pi a\log B$ is bounded. Does it converge, and is its limit the slit-plane constant
$\pi(b-a\log2)$ of the lower bound? For a finite prescribed divisor, what is the optimizer of the joint problem,
including possible atoms at the endpoints of the support, and can \cref{fil:thm:slit} be sharpened to an exact
constrained minimax formula like that of \cref{fil:prop:cheb}?
\item \emph{Distinct ordinates and uniformity.} If only one vanishing at each distinct ordinate is required, what
are the leading support constant and the optimal decay exponent of a global filter, without assuming simplicity
or a density asymptotic for distinct ordinates? Can the constructions be made uniform as the conductor or the
number of functions grows, and can the support, variance proxy and number of atoms be optimized jointly?
\item \emph{Global filters.} By \cref{fil:thm:endpoint} the supremal coefficient is $\pi d/Q_*$ for every family,
and the endpoint is a finite-dimensional positivity problem. When the canonical kernel $h$ changes sign, which
polynomials $P$ of the admissible degree make $P(-\partial_t)h$ nonnegative, and for many copies of $\chi_5$ where is
the boundary of that cone of coefficients? Can the number and size of the translations in the subcritical
construction be bounded as $b\uparrow\pi d/Q_*$, and is there a rigidity statement for near-extremizers under a
uniform moment normalization? What is the least width of a truncated approximate filter chosen depending on $x$,
beyond the fixed rule of \cref{fil:prop:tail}? Does the division argument extend to completed Dedekind or
automorphic $L$-functions with several gamma factors, under their own critical-line hypothesis? Which further
completed $L$-functions have positive theta kernels or explicit compression constants like those of
\cref{fil:thm:zetaexplicit}, and can the rational mesh and weights of \cref{fil:thm:quad} be bounded explicitly from zero
enclosures? Beyond the residual certificate for separated classes after \cref{rec:thm:uniformavg}, can the positive
interpolation of \cref{rec:lem:sepfilter,fil:lem:cosine} be certified when clustered or multiple ordinates are known
only through interval enclosures, with control of both residual and width?
\item \emph{Approximate cancellation.} \Cref{rec:thm:polylogavg} subtracts the finite Euler-factor terms and the
residues at $0$ and at the trivial zeros explicitly and reaches any fixed power accuracy. Does approximate
cancellation alone have a universal leading support coefficient under a natural robustness requirement?
\item \emph{Uniform versions of the family statistics.} The constants in \cref{ana:thm:grid} are for a fixed
family. Can they be made effective and uniform in the conductor, with explicit values of the
local-count, remainder and gap constants $A$, $E$ and $\eta$ of \cref{ana:thm:alias} for a
chosen modulus, and hence an explicit rounding threshold in \cref{ana:thm:grid}?
\Cref{ana:prop:alias} shows that the alias term $\log k/k^2$ cannot be
removed on the lattice $k^{-1}\Z$ within the abstract spectral class; is the rate of
\cref{ana:thm:grid} optimal among irregular sampling schedules, or for actual Dirichlet zero
spectra? Do the converses of \cref{ana:thm:gridconverse,ana:thm:binomial}(i) hold for a depth
that grows slowly with $t$?
\end{enumerate}
